% concatenated from MQF_20260506.tex
\documentclass[a4paper, 11pt]{amsart}
\usepackage{amssymb,amsmath,comment,mathrsfs,amsfonts,mathtools}
\usepackage{color}
\usepackage{enumitem}
\usepackage[top=30truemm,bottom=30truemm,left=25truemm,right=25truemm]{geometry}%
\usepackage[normalem]{ulem}

\newcommand{\notmid}{\mathrel{\ooalign{$\mkern-5mu\not$\crcr$|$}}}
\newcommand{\A}{\mathbb A}
\newcommand{\N}{\mathbb N}
\newcommand{\Z}{\mathbb Z}
\newcommand{\Q}{\mathbb Q}
\newcommand{\C}{\mathbb C}
\newcommand{\F}{\mathbb F}
\newcommand{\PS}{\mathbb P}
\newcommand{\Gal}{\mathrm{Gal}}
\newcommand{\Gm}{{\mathbb G}_m}
\newcommand{\Hom}{\mathrm{Hom}}
\newcommand{\Tr}{\mathrm{Tr}}

\theoremstyle{plain}
\newtheorem{thm}{Theorem}[section]
\newtheorem*{thm*}{Theorem}
\newtheorem{prob}[thm]{Problem}
\newtheorem{prop}[thm]{Proposition}
\newtheorem{lem}[thm]{Lemma}
\newtheorem{cor}[thm]{Corollary}

\theoremstyle{definition}
\newtheorem{defn}[thm]{Definition}
\newtheorem{remark}[thm]{Remark}
\newtheorem{example}[thm]{Example}
\newtheorem*{acknowledgments}{Acknowledgments}
\newtheorem*{acknowledgment}{Acknowledgment}

\def\vector#1{\mbox{\boldmath $#1$}} 
%\makeindex

\renewcommand{\labelenumi}{\textup{(\arabic{enumi})}}
\renewcommand{\thefootnote}{\fnsymbol{footnote}}

\makeatletter
\renewcommand*\l@subsection{\@tocline{2}{0pt}{1.8em}{3.0em}{}}
\makeatother

\numberwithin{equation}{section}

%\raggedbottom
\allowdisplaybreaks[4]

\begin{document}
\title[Multiplicative $f$-ic forms on algebraic varieties]
{Multiplicative $f$-ic forms on algebraic varieties arising from Thaine's generalized Jacobi sums}

\author[A. Hoshi]{Akinari Hoshi}
\address[A. Hoshi]{Department of Mathematics, Niigata University, Niigata 950-2181, Japan}
\email{hoshi@math.sc.niigata-u.ac.jp}

\author[K. Kanai]{Kazuki Kanai}
\address[K. Kanai]{Department of Computer and Information Sciences, Ibaraki University, 4-12-1, Ibaraki, 316-8511, Japan}
\email{kazuki.kanai.du62@vc.ibaraki.ac.jp}

\footnote[0]{\textit{$2020$ Mathematics Subject Classification}. 
11D57, 11R18, 11T22, 11T24, 14M10.}
%11D09 Quadratic and bilinear Diophantine equations
%11D57 Multiplicative and norm form equations
%11R18 Cyclotomic extensions
%11T22 Cyclotomy
%11T23 Exponential sums
%11T24 Other character sums and Gauss sums
%14M10 Complete intersections
\footnote[0]{\textit{Keywords and phrases}. 
Jacobi sums, Gauss sums, Gaussian periods, multiplication matrices, 
cyclotomy, Dickson's system, multiplicative quadratic forms.}
%finite Fourier transform
\footnote[0]{This work was partially supported by JSPS KAKENHI Grant Numbers 19K03418, 20H00115, 24K00519.}


%%%%%%%%%%%%%%%%%%%%%%%%%%%%%%%%%%%%%%%%%
\begin{abstract}
We study generalized Jacobi sums, cyclotomic numbers, and $d$-compositions in Thaine's framework, and prove new multiplicative identities extending Davenport and Hasse's lifting theorem from the classical prime-power setting to products of prime powers. 
As applications, we construct multiplicative forms of degree $f\ge2$, i.e. $f$-ic forms, on complete intersections of $f$-ics. 
This places Pfister's theory of multiplicative quadratic forms over fields within the broader setting of multiplicative $f$-ic forms on affine algebraic varieties, where new phenomena arise.
Moreover, a dense open subset $W \subset V$ carries the structure of an algebraic torus, and the multiplicative form is compatible with the induced group law on $W$.
\end{abstract}

\maketitle
%%%%%%%%%%%%%%%%%%%%%%%%%%%%%%%%%%%%%%%%%
\tableofcontents

%------------------------------------S1
\section{Introduction}\label{S1}

This paper has two closely related themes. 
On the arithmetic side, we establish a multiplicative law for Thaine's generalized Jacobi sums and generalized multiplication matrices, extending Davenport and Hasse's lifting phenomenon from the classical prime-power setting to products of prime powers. 
On the geometric side, we apply these identities over a field $K$ to construct multiplicative forms on complete intersections, 
first in the quadratic case and then in higher degree, and we show that a dense open subset carries the structure of an algebraic $K$-torus.
The key point is that Thaine's $d$-composition is not only a formal operation on cyclotomic data:
it yields multiplicative identities for generalized Jacobi sums and, at the same time, produces
algebraic composition laws on complete intersections, whose invertible locus is an algebraic torus.

Thus the algebraic torus appearing at the end of the construction is not introduced as an independent object. 
It is the invertible part of the coefficient variety obtained by geometrizing the relations satisfied by Jacobi sums and their generalized counterparts.
The coefficients of Jacobi sums provide the link between the cyclotomic identities
and the algebraic varieties constructed in this paper.

We work in the framework of generalized Jacobi sums, generalized multiplication matrices,
and $d$-compositions introduced by Thaine~\cite{Tha04}, together with the study of the $(-1)$-composition and its lifting properties by Hoshi and Kanai \cite{HK22}. 
The present paper extends Thaine's multiplicative formalism from prime powers to products and places the lifting results of~\cite{HK22} in the framework of $d$-compositions.
This also leads to an extension of the Diophantine systems studied by Katre and Rajwade \cite{KR85a} and van Wamelen \cite{Wam02} beyond the classical prime-power case. 
From the geometric point of view, the paper continues the study of multiplicative forms on algebraic varieties over $K$ initiated by Hoshi \cite{Hos03}, moving from quadratic examples to multiplicative $f$-ic forms on complete intersections and to the algebraic $K$-group structure carried by a dense open subset.

\subsection{Arithmetic results: generalized Jacobi sums and $d$-compositions}\label{S1-1}

Let $D$ be an integrally closed  domain with char $D=0$, and let $K$ be its quotient field. 
Let $e \geq 2$ be an integer and 
$P(X)=\prod_{i=0}^{e-1}(X-\theta_i), 
P^{\prime}(X)=\prod_{i=0}^{e-1}(X-\theta^{\prime}_i) \in D[X]$ 
be cyclic polynomials of degree $e$, 
i.e. irreducible polynomials with cyclic Galois groups over $K$. 
Then $L=K(\theta_i)$ and $L^{\prime}=K(\theta^{\prime}_i)$ 
are cyclic extensions of $K$ of degree $e$ with 
$\Gal(L/K)=\langle\tau\rangle$, 
$\tau(\theta_i)=\theta_{i+1}$ 
and 
$\Gal(L^{\prime}/K)=\langle\tau^{\prime}\rangle$, 
$\tau^{\prime}(\theta^{\prime}_i)=\theta^{\prime}_{i+1}$ 
where subscripts are taken modulo $e$. 

We assume that $L \cap L^{\prime}= K$. 
We also 
assume that 
$\{\theta_0,\ldots,\theta_{e-1}\}$ 
and 
$\{\theta^{\prime}_0,\ldots,\theta^{\prime}_{e-1}\}$ 
are linearly independent over $K$. 
Then $[L L^{\prime}:K]=e^2$ with 
$G={\rm Gal}(L L^{\prime}/K)\simeq \langle \tau \rangle \times \langle \tau^{\prime} \rangle$ 
and $\{\theta_i\theta^{\prime}_j\mid 0\leq i,j\leq e-1\}$ 
becomes a normal basis of $L L^{\prime}$ over $K$. 
%
We make identifications $\tau = (\tau,1)$ 
and $\tau^{\prime} = (1, \tau^{\prime})$.  
For $1\leq d\leq e-1$, 
we have intermediate fields 
$L_d=(L L^{\prime})^{\langle \tau {\tau^{\prime}}^{d} \rangle}=K(\theta_{d,i})$ which are cyclic extensions of $K$ of degree $e$ 
with a normal basis $\{\theta_{d,0},\ldots,\theta_{d,e-1}\}$ 
where $\theta_{d,i} = \mathrm{Tr}_{L L^{\prime}/L_d}(\theta_0\theta^{\prime}_i)= \sum^{e-1}_{s=0}\theta_s \theta^{\prime}_{ds+i}$. 

Let $M_e(K)$ be the algebra of $e\times e$ matrices over $K$. 
Define the matrix $A=[a_{i,j}]_{0\leq i,j \leq e-1}\in M_e(K)$ 
by $\theta_0\theta_i=\sum^{e-1}_{j=0}a_{i,j}\theta_j$
which is called 
{\it the multiplication matrix of $\theta_0,\ldots,\theta_{e-1}$}. 
Note that the $\theta_i$'s are eigenvalues of $A$ and hence 
$P(X)={\rm Char}_X(A)$, the characteristic polynomial of $A$. 

%%%%%%%%%%%%%%%%%%%%%%%%%%%%%%%%%%%%%%%%%%%%%
Thaine \cite[Section 2]{Tha04} defined 
the $d$-composition $A\overset{d}{\ast}B$ 
of the matrices $A$ and $B$ as follows. 
For $A=[a_{i,j}]_{0\leq i,j\leq e-1}$, 
$B=[b_{i,j}]_{0\leq i,j\leq e-1}\in M_e(K)$ 
and $d\in(\Z/e\Z)\setminus\{0\}$, 
{\it the $d$-composition $A \overset{d}{\ast} B$ of $A$ and $B$} is defined by 
\[
A \overset{d}{\ast} B:= \left[\sum^{e-1}_{s=0}\sum^{e-1}_{t=0}a_{s,t}b_{ds+i,dt+j}\right]_{0\leq i,j\leq e-1}.
\]
Thaine \cite[Proposition~6]{Tha04} (see also \cite[Example 4]{Tha04}) 
proved that the multiplication matrix $A_d$ 
of $\theta_{d,0},\ldots,\theta_{d,e-1}$ is given by 
\begin{align*}
A_d=A\overset{d}{\ast}A^{\prime}
\end{align*}
where $A=[a_{i,j}]_{0\leq i,j\leq e-1}$ and 
$A^{\prime}=[a^{\prime}_{i,j}]_{0\leq i,j\leq e-1}$ are 
the multiplication matrices of 
$\theta_{0},\ldots,\theta_{e-1}$ 
and $\theta^{\prime}_{0},\ldots,\theta^{\prime}_{e-1}$ respectively.  
It turns out that 
the multiplication matrix $A_d$ corresponds to the intermediate field 
$L_d=(L L^{\prime})^{\langle \tau {\tau^{\prime}}^{d} \rangle}=K(\theta_{d,i}) 
\subset LL^{\prime}$ as explained above. 

Hoshi and Kanai \cite[Section 3]{HK22} studied the $d$-compositions 
$\overset{d}{\ast}$.
In general, the $d$-composition $\overset{d}{\ast}$ 
is neither associative nor commutative, 
although the $(-1)$-composition is associative and commutative. 
Thus we can define the $n$-fold product $A^{(n)}$ of $A$ as 
\begin{align*}
A^{(n)} := A \overset{-1}{\ast} \cdots \overset{-1}{\ast} A\quad(\text{$n$-fold}).
\end{align*}
By using Davenport and Hasse's lifting theorem for Gaussian periods, 
Hoshi and Kanai \cite[Theorem~1.3 and Theorem~1.4]{HK22} proved that
\begin{align} \label{eq:HKmainThm}
C_{p^{nr}} = (-1)^{n-1}C_{p^r}^{(n)}
\end{align}
where $C_{p^r}$ is the multiplication matrix of Gaussian periods $\eta_{p^r}(0),\ldots,\eta_{p^r}(e-1)$ (see Section \ref{S2} for Gaussian periods $\eta_{p^r}(i)$). 
This is the matrix analogue of Davenport and Hasse's lifting theorem.

Let 
$\zeta_n=e^{2\pi i/n}$ be a primitive $n$-th root of unity. 
Thaine \cite{Tha04} introduced generalized Jacobi sums
$E_{{\bf p}}(a,b)\in K(\zeta_e)$
and generalized multiplication matrices $H_{{\bf p}}=[h_{i,j}]_{0\leq i,j \leq e-1}
\in M_e(K)$ for ${\bf p} \in K^{\times}$, 
by abstracting the properties of classical Jacobi sums and multiplication matrices
(see Definition \ref{def:T9and10} for generalized Jacobi sums $E_{{\bf p}}(a,b)$ 
and Definition \ref{def:T9and10_2} for generalized multiplication matrices $H_{{\bf p}}$ for ${\bf p} \in K^{\times}$). 

Thaine's formalism behaves well with respect to products: generalized Jacobi sums and generalized multiplication matrices attached to ${\bf p}_1$ and ${\bf p}_2$ canonically yield those attached to ${\bf p}_1{\bf p}_2$.
In particular, in the prime-power case, classical Jacobi sums for $p^r$ and $q^s$ combine to yield a generalized
Jacobi sum attached to $p^rq^s$.

The main theorem of this paper gives generalized Jacobi sums for the product ${\bf p}_1{\bf p}_2 $, together with generalized multiplication matrices for ${\bf p}_1{\bf p}_2$ corresponding to ${\rm Cyc}_{{\bf p}_1}\overset{d}{\ast}{\rm Cyc}_{{\bf p}_2}$ where ${\bf p}_1, {\bf p}_2 \in K^\times$:

\begin{thm}\label{mainth1}
Let $e \geq 2$ be an integer and $d \in (\Z/e\Z)^{\times}$. 
Let $K$ be a field with 
char $K=0$ and $\Gal(K(\zeta_e)/K)\simeq (\Z/e\Z)^{\times}$, 
and 
let $\sigma_{-d}: K(\zeta_e) \to K(\zeta_e)$, $\zeta_e \mapsto \zeta_e^{-d}$ 
be a $K$-automorphism of $K(\zeta_e)$. 
Let 
$E_{{\bf p}_u}(a,b)$ be generalized Jacobi sums 
for ${\bf p}_u \in K^\times$ with respect to
an integer $v_u$ satisfying $ev_u/2 \in \Z$, for $u=1,2$.  
Then we have:
\begin{enumerate}
\item $(E_{{\bf p}_1}\overset{d}{\ast}E_{{\bf p}_2})(a,b):= -\sigma_{-d}(E_{{\bf p}_1}(a,b))E_{{\bf p}_2}(a,b)$ are generalized Jacobi sums for ${\bf p}_1{\bf p}_2 \in K^\times$ with respect to $v_1+v_2$;
\item for $f_u = ({\bf p}_u -1)/e$, $v_u \in \Z$ with $\frac{ev_u}{2} \in \Z$ and
\begin{align*}
{\rm Cyc}_{{\bf p}_u}&:=H_{{\bf p}_u} + [f_u\delta_{i,\frac{ev_u}{2}}]_{0\leq i,j\leq e-1} \ (u=1,2)
\end{align*}
where the $(i,j)$ entry $H_{{\bf p}_u}[i,j]$ of the matrix $H_{{\bf p}_u}$ 
is given by 
\begin{align*}
H_{{\bf p}_u}[i,j]&=-f_u\delta_{i,\frac{ev_u}{2}}+\frac{1}{e^2}\sum^{e-1}_{a=0} \sum^{e-1}_{b=0} (-1)^{v_ua} \zeta_e^{-(ai+bj)}E_{{\bf p}_u}(a,b),
\end{align*}
we have
\begin{align*}
(-1)^{(v_1+ v_2)a+1}\sum_{i=0}^{e-1}\sum_{j=0}^{e-1}({\rm Cyc}_{{\bf p}_1}\overset{d}{\ast}{\rm Cyc}_{{\bf p}_2})[i,j]\,\zeta_e^{ai+bj}
=(E_{{\bf p}_1}\overset{d}{\ast}E_{{\bf p}_2})(a,b).
\end{align*}
\end{enumerate}

In particular, for prime power ${\bf p}_1=p^r$ $($resp.  ${\bf p}_2=q^s$$)$ 
and Jacobi sums
$E_{{\bf p}_1}(a,b)=J^\ast_{p^r}(a,b)$ 
$($resp. $E_{{\bf p}_2}(a,b)=J^\ast_{q^s}(a,b)$$)$ for $\F_{p^r}$ 
$($resp. $\F_{q^s}$$)$, we have:
\begin{enumerate}
\item $(J^{*}_{p^r}\overset{d}{\ast}J^{*}_{q^s})(a,b):= -\sigma_{-d}(J^\ast_{p^r}(a,b))J^\ast_{q^s}(a,b)$ are generalized Jacobi sums for $p^rq^s$;
\item for $v_1=f_1$, $v_2=f_2$, we have
\[
(-1)^{\lambda a +1} \sum_{i=0}^{e-1}\sum_{j=0}^{e-1}({\rm Cyc}_{p^r}\overset{d}{\ast}{\rm Cyc}_{q^s})[i,j]\,\zeta_e^{ai+bj}
= (J^{*}_{p^r}\overset{d}{\ast}J^{*}_{q^s})(a,b)
\]
where $f:=(p^rq^s-1)/e =f_1f_2e+f_1+f_2$ and
\[
\lambda = 
\begin{cases}
f & (p,q)\neq(2,2),\\ 
0 & (p,q)=(2,2)
\end{cases}
\]

i.e. the $d$-compositions of cyclotomic numbers ${\rm Cyc}_{p^r}\overset{d}{\ast}{\rm Cyc}_{q^s}$ 
correspond to generalized Jacobi sums $(J^{*}_{p^r}\overset{d}{\ast}J^{*}_{q^s})(a,b)$ up to sign 
$($cf. the equation \eqref{eq:Jr-Cycr} in Section $\ref{S2}$$)$.
\end{enumerate}
\end{thm}

Theorem~\ref{mainth1} shows that classical Jacobi sums are intrinsically attached to prime powers, whereas Thaine's generalized Jacobi sums are multiplicative with respect to products. 
Thus the $d$-composition combines the classical Jacobi sums for $p^r$ and $q^s$ into a generalized Jacobi sum attached to $p^rq^s$, extending the Diophantine systems considered by Katre and Rajwade \cite{KR85a} and van Wamelen \cite{Wam02} beyond the classical prime-power case.
These multiplicative identities are the arithmetic counterpart of the geometric composition maps constructed below, which send points of the fibers over $p$ and $q$ to points of the fiber over $pq$.

%%%%%%%%%%%%%%%%%%%%%%%%%%%%%%%%%%%%%%%%%%%%%%%%%%%%%%%%%%%%%%%%%%
\begin{remark}\label{rem:Sto}
Storer \cite{Sto67}, \cite{Sto69} studied cyclotomic numbers 
with respect to the direct sum $\F_{p^r} \oplus \F_{q^s}$ combinatorially. 
Theorem~\ref{mainth1} gives a generalization of his results 
using the $d$-composition $\overset{d}{\ast}$.
\end{remark}

Applying Theorem~\ref{mainth1} to the case where
${\bf p}_1={\bf p}_2=p^r$ is a prime power
and $d=-1$, 
we obtain the following equivalence between 
(1) Davenport and Hasse's lifting theorem for Jacobi sums and 
(2) the lift of multiplication matrices of Gaussian periods.
%
\begin{thm}\label{thm:corMT1}
Let $e \geq 2$ be an integer, $p^r$ be a prime power with 
$p^r\equiv 1\ ({\rm mod}\ e)$ 
and $\chi$ be the character of order $e$ 
on $\F_{p^r}$ with $\chi(\gamma)=\zeta_e$ and $\chi(0)=0$ 
where $\langle\gamma\rangle=\F_{p^r}^\times$. 
Let $\chi^{\prime}$ be the lift of $\chi$ 
from $\F_{p^r}$ to $\F_{p^{nr}}$ defined by 
$\chi^{\prime}(\alpha) = \chi(\mathrm{Nr}(\alpha))$ 
and $\mathrm{Nr}$ is the norm map.
Let $C_{p^r} = [{\rm Cyc}_{p^r}(i,j) - D_i f]_{0\leq i,j\leq e-1}$ be the multiplication matrix of Gaussian periods $\eta_r(0),\ldots,\eta_r(e-1)$ of degree $e$ for $\F_{p^{r}}$. 
For any integer $n \ge 1$,
the following conditions are equivalent:
\begin{enumerate}
\item $J_{p^{nr}}^{\ast}((\chi^{\prime})^a, (\chi^{\prime})^b)=(-1)^{n-1}J_{p^r}^{\ast}(a, b)^n$ for all $a,b \in \Z$ such that $e\nmid a,\  e\nmid b$ and $e\nmid a+b$;
\item $C_{p^{nr}} = (-1)^{n-1}C_{p^r}^{(n)}$
\end{enumerate}
where $C_{p^r}^{(n)}$ is the $n$-fold product of $C_{p^r}$ with respect to the $(-1)$-composition $\overset{-1}{\ast}$. 
\end{thm}

Based on Theorem~\ref{thm:corMT1}, Theorem~\ref{mainth1} (1) for generalized Jacobi sums may be regarded as, in the sense of Thaine \cite{Tha04}, a generalization of Davenport and Hasse's lifting theorem for (classical) Jacobi sums.


\subsection{Geometric results: multiplicative forms on algebraic varieties}\label{S1-2}

In the geometric direction, we use the arithmetic identities above to construct multiplicative forms on algebraic $K$-varieties, first in the quadratic case and then in higher degree. 
This extends the study of multiplicative quadratic forms on algebraic $K$-varieties in Hoshi~\cite{Hos03}.

As a first geometric application of Theorem~\ref{mainth1}, we construct multiplicative quadratic forms of dimension $l-1$ on algebraic $K$-varieties given by intersections of $\frac{l-1}{2}-1$ quadrics.
This provides the quadratic prototype for the higher-degree construction of Theorem~\ref{mainth3}.

\begin{thm}\label{mainth2}
Let $l$ be an odd prime. 
Let $K$ be a field with char $K\neq 2$, $l$ and 
$\Gal(K(\zeta_l)/K)\simeq (\Z/l\Z)^{\times}$. 
Let $V \subset K^{l-1}$ be
an 
algebraic $K$-variety defined by
\begin{align*}
f_m({\bf x})=0\quad (m=1,\dots, \tfrac{l-1}{2}-1)
\end{align*}
where
\[
f_m({\bf x})=f_m(x_1,\ldots,x_{l-1})
=\sum_{i=1}^{l-1}x_ix_{i+m}-\sum_{i=1}^{l-1}x_ix_{i+(m+1)}
\]
and each subscript of the $x_{i}$'s should be taken modulo $l$ with the convention $x_0=0$.
Let $q({\bf x})=q(x_1,\ldots,x_{l-1})$ be a regular quadratic form defined by
\[
q({\bf x})=f_0({\bf x})
=\sum_{i=1}^{l-1}x_i^2-\sum_{i=1}^{l-1}x_ix_{i+1}.
\]
Then $q({\bf x})$ is multiplicative on $V$, 
i.e. there exists a bilinear map 
$\varphi :K^{l-1}\times K^{l-1} \rightarrow K^{l-1}$ such that
\[
\varphi(V \times V) \subset V \quad \mbox{and}\quad q({\bf x})q({\bf y})=q(\varphi({\bf x,y}))\mbox{ for any }{\bf x,y} \in V. 
\]
Moreover, for $d \in (\Z/l\Z)^{\times}$, $\varphi_d ({\bf x},{\bf y})=(z_{1,d},\ldots,z_{l-1,d}) \in V$ 
can be given explicitly by
\[
z_{k,d}=-\left(\sum_{i=1}^{l-1}x_{-d^{-1}i}y_{k-i} - \sum_{i=1}^{l-1}x_{-d^{-1}i}y_{-i}\right)\quad (k=1,\ldots,l-1)
\]
where each subscript of the $x_i$'s and the $y_j$'s should be taken modulo $l$ with the convention $x_0=y_0=0$.
\end{thm}

It is also known that the cyclotomic numbers ${\rm Cyc}_{p^r}(a,b)$ are determined by the same system of Diophantine equations as the Jacobi sums for $p^r$. 
In particular, when $e=3$, the corresponding Diophantine equation reduces to the classical binary quadratic form $4p^r=x_1^2+27x_2^2$ appearing in Gauss's \emph{Disquisitiones Arithmeticae} (see Section~\ref{S4-3}). 
In \cite[Section 7]{HK22}, by computing the $n$-fold $(-1)$-product $C_{p^r}^{(n)}$ of the multiplication matrix $C_{p^r}$ via Thaine's $d$-composition, we obtained explicit lifts of solutions of the corresponding Diophantine systems. 
Thus, for a fixed prime power $p^r \equiv 1 \pmod e$, three lift phenomena run in parallel: Davenport and Hasse's lifting for Jacobi sums, the matrix lifting formula \eqref{eq:HKmainThm} for multiplication matrices of Gaussian periods (and hence for cyclotomic numbers), and explicit lifts of solutions of the associated Diophantine systems. 
We also extend this parallelism from the prime-power case to the more general setting of $d$-compositions.

Theorem~\ref{mainth2} serves as the quadratic model for the higher-degree case.
The passage to higher degree is motivated by the observation that a relative norm relation in a cyclotomic extension yields a multiplicative form on a complete intersection.
To state the higher-degree result, recall that an $n$-ic form
$h(x_1,\dots,x_m)$ over $K$ is called regular if the associated symmetric $n$-linear form
$\theta$ is nondegenerate in the sense that
\[
\theta(v_1,\dots,v_n)=0 \quad \text{for all } v_2,\dots,v_n \in K^m
\]
implies $v_1=0$; see Shapiro \cite[Chapter~16]{Sha00}.

For the higher-degree construction, we also assume that $\operatorname{char} K \nmid f!$ where $f!$ denotes the factorial of  an integer $f$. 
This hypothesis is used in passing from a homogeneous $f$-ic form to its associated symmetric $f$-linear form, and in the proof of regularity.

In the quadratic case, the classical relation $J^{*}_{p^r}(a,b)\overline{J^{*}_{p^r}(a,b)} = p^r$ may be interpreted as a relative norm relation from the cyclotomic field $L=\mathbf{Q}(\zeta_l)$ to its maximal real subfield $M=\mathbf{Q}(\zeta_l+\zeta_l^{-1})$. 
Motivated by this picture, let $l$ be an odd prime and $f \ge 2$ with $f \mid (l-1)$, and let $M$ be the unique intermediate field of $K(\zeta_l)/K$ such that $[K(\zeta_l):M]=f$. 
For an element
\[
H=\sum_{k=0}^{l-1} a_k \zeta_l^k \qquad (a_k \in K)
\]
satisfying
\[
N_{K(\zeta_l)/M}(H)={\bf p} \in K^\times,
\]
we construct multiplicative $f$-ic forms on algebraic $K$-varieties realized as complete intersections
of $e-1$ hypersurfaces of degree $f$:

\begin{thm}\label{mainth3}
Let $f\geq2$ be an integer and $l$ be an odd prime with $l \equiv 1\ ({\rm mod}\ f)$. 
Write $l=ef+1$.
Let $K$ be a field with char $K \notmid f!$, $l$ and 
$\Gal(K(\zeta_l)/K)\simeq (\Z/l\Z)^{\times}=\langle \gamma \rangle$. 
Define $\beta := \gamma^e$ and $M=L^{\langle \beta \rangle}$ where $L=K(\zeta_l)$, 
i.e., $\Gal(L/M)\simeq \Z/f\Z = \langle \beta \rangle \leq  (\Z/l\Z)^{\times}$.
Let $V \subset K^{l-1}$ be
an algebraic $K$-variety defined by
\begin{align*}
h_m({\bf x})=0\quad (m = 1,\ldots,e-1)
\end{align*}
where
\begin{align*}
h_m({\bf x})&=h_m(x_1,\ldots,x_{l-1})\\
&=\left(\sum_{i_0 + \beta i_1 + \cdots +\beta^{f-1}i_{f-1} \equiv \gamma^m \ (l)}x_{i_0}x_{i_1}\cdots x_{i_{f-1}}\right) - \left(\sum_{i_0 + \beta i_1 + \cdots +\beta^{f-1}i_{f-1} \equiv \gamma^{m+1} \ (l)}x_{i_0}x_{i_1}\cdots x_{i_{f-1}}\right)
\end{align*}
and each subscript of the $x_{i}$'s should be taken modulo $l$ with the convention $x_0=0$.
Let $h({\bf x})=h(x_1,\ldots,x_{l-1})$ be a regular $f$-ic form defined by
\begin{align*}
h({\bf x})
&=\left(\sum_{i_0 + \beta i_1 + \cdots +\beta^{f-1}i_{f-1}\equiv 0 \ (l)}x_{i_0}x_{i_1}\cdots x_{i_{f-1}}\right)
-\left(\sum_{i_0 + \beta i_1 + \cdots +\beta^{f-1}i_{f-1}\equiv \gamma \ (l)}x_{i_0}x_{i_1}\cdots x_{i_{f-1}}\right).
\end{align*}
Then $h({\bf x})$ is multiplicative on $V$.
Moreover, for $d \in (\Z/l\Z)^{\times}$, $\varphi_d ({\bf x},{\bf y})=(z_{1,d},\ldots,z_{l-1,d}) \in V$ 
can be given explicitly by
\[
z_{k,d}=(-1)^{f-1} \left(\sum_{i=1}^{l-1}x_{-d^{-1}i}y_{k-i} - \sum_{i=1}^{l-1}x_{-d^{-1}i}y_{-i}\right)\quad (k=1,\ldots,l-1)
\]
where each subscript of the $x_i$'s and the $y_j$'s should be taken modulo $l$ with the convention $x_0=y_0=0$.
\end{thm}

For each $d \in (\Z/l\Z)^\times$, we define binary operations ${\bf x}\overset{d}{\ast}{\bf y}:=\varphi_d({\bf x},{\bf y})$ on $V$, where $\varphi_d$ is defined as in Theorem~\ref{mainth3}.
The operation $\overset{d}{\ast}$ is commutative and associative if and only if $d=-1$ (see Proposition~\ref{prop:CA}).
 
The compositions depend on $d$ and are generally neither commutative nor associative. 
Only for $d=-1$ do they define a commutative algebraic group law, 
and this group law is defined not on the whole variety $V$ but on a dense open subset $W \subset V$.

\begin{thm}\label{thm:structureV}
Let $K$, $L=K(\zeta_l)$, $M$ and the algebraic $K$-variety $V:=\{{\bf x}\in K^{l-1} \mid h_{m}({\bf x})=0\quad (m = 1,\ldots,e-1) \} \subset K^{l-1}$ be as in Theorem~\ref{mainth3}. 
For ${\bf x}=(x_1,\dots,x_{l-1})\in V$, define
\[
  \alpha_{\bf x}:=(-1)^{f-1}\sum_{i=1}^{l-1} x_i \zeta_l^{\,i}\in L.
\]
Let
\[
  W:=\{{\bf x}\in V \mid N_{L/K}(\alpha_{\bf x})\neq 0\}=V^\times=V\setminus S
  \quad {\text where}\quad S:=\{{\bf x}\in V \mid N_{L/K}(\alpha_{\bf x})=0\}.
\]
Then the variety $W$ equipped with the operation $\overset{-1}{\ast}$ is isomorphic to an algebraic $K$-torus and we have
\[
  W \simeq N_{L/M}^{-1}(\mathbb{G}_{m,K})= \left\{ \alpha \in R_{L/K}(\mathbb{G}_m) \ \Bigm|\  N_{L/M}(\alpha) \in \mathbb{G}_{m,K} \right\}
\]
where $R_{L/K}$ denotes the Weil restriction. 
Here $\mathbb{G}_{m,K}$ is regarded as the diagonal subtorus of $R_{M/K}(\mathbb{G}_m)$ induced by $K \subset M$.
\end{thm}

For background on algebraic $K$-tori, we refer to Voskresenskii \cite[Chapter~2]{Vos98} and Manin \cite[Section~30]{Man74}.
As a consequence of the torus description of the dense open subset $W$, Corollary~\ref{cor:CI} will be deduced in Section~\ref{S6}.

\begin{cor}\label{cor:CI}
The algebraic $K$-variety $V \subset K^{l-1}$ defined by the system of equations $h_{m}(x)=0$ for $m=1,\dots,e-1$ is a complete intersection in $\mathbb{A}^{l-1}_K$, cut out by $e-1$ $f$-ics over $K$.
\end{cor}

For ${\bf p}\in K^\times$ we define
\[
W_{\bf p}(K):=\{\,x\in W(K)\mid h(x)={\bf p}\,\}.
\]
Since $h$ is multiplicative with respect to $\overset{-1}{\ast}$, i.e.\ $h(x\overset{-1}{\ast}y)=h(x)h(y)$ for $x,y\in W(K)$, it follows that $W_{\bf p}(K)\overset{-1}{\ast}W_{\bf q}(K)\subseteq W_{{\bf p}{\bf q}}(K)$.

\begin{cor}\label{cor:lifting}
With notation as above,  for any $\bf{p},\bf{q} \in K^\times$, the group law $\overset{-1}{\ast}$ on $W(K)$ induces
a well-defined map
\[
  W_{\bf{p}}(K) \times W_{\bf{q}}(K) \longrightarrow W_{\bf{pq}}(K),
  \qquad (x,y) \longmapsto x \overset{-1}{\ast} y.
\]
\end{cor}

In terms of the fibers $W_p$ introduced before Corollary~\ref{cor:lifting},
classical Jacobi sums in the prime-power case furnish integral points on $W_{p^r}$. 
The composition law then yields integral points on $W_{p^r q^s}$ from those on $W_{p^r}$ and $W_{q^s}$.

\medskip
\medskip
\noindent {\bf Organization of the paper.}~~
We organize this paper as follows.
In Section~\ref{S2} we collect basic definitions and known results on Gaussian periods, Jacobi sums, cyclotomic numbers, and the associated multiplication matrices, and we also review Thaine's generalized Jacobi sums and generalized multiplication matrices of Gaussian periods following Thaine \cite{Tha04}.
In Section~\ref{S3}, we prove Theorems~\ref{mainth1} and~\ref{thm:corMT1} and interpret Theorem~\ref{mainth1} in terms of a twisted double finite Fourier transform, which also yields an alternative proof in a restricted case.
In Section~\ref{S4} we develop the quadratic part of the paper: 
after recalling multiplicative quadratic forms over fields, Pfister forms, the level of fields,
and multiplicative quadratic forms on algebraic $K$-varieties due to Hoshi~\cite{Hos03}, 
we prove Theorem~\ref{mainth2} and apply Theorem~\ref{thm:BF} to obtain explicit lifts of solutions of the Diophantine systems attached to generalized Jacobi sums and to multiplicative quadratic forms on algebraic $K$-varieties.
In Section~\ref{S5} we construct multiplicative $f$-ic forms on affine algebraic $K$-varieties.
Finally, in Section~\ref{S6} we study the induced composition law on the variety $V$ of Section~\ref{S5} and show that an open subset $W\subset V$ carries a commutative algebraic $K$-group structure (indeed, an algebraic  $K$-torus).

\medskip
\medskip
\noindent {\bf Notation.}~~
We record here the principal notational conventions used in the statements of the main results, especially where several closely related objects appear in parallel.

\medskip
\textit{Basic parameters and fields:}
\par\smallskip
\begin{tabular}{@{} r @{\quad} l}
 $l$: & An odd prime number. \\
 $f$: & An integer $f \ge 2$ dividing $l-1$, representing the degree of the form $h(x)$. \\
 $e$: & An integer determined by $l-1 = ef$. \\
 $\zeta_n$: & A primitive $n$-th root of unity. \\
 $K$: & A base field with $\text{char}\, K \notmid f!, l$ and $\text{Gal}(K(\zeta_l)/K) \simeq (\mathbb{Z}/l\mathbb{Z})^\times$. \\
 $L$: & The cyclotomic extension $K(\zeta_l)$. \\
 $M$: & The unique intermediate field $K \subset M \subset L$ such that $[L:M]=f$.\\
 $\sigma_a$: & The automorphism in $\text{Gal}(L/K)$ defined by $\sigma_a(\zeta_l) = \zeta_l^a$. \\
 $N_{L/K}$: & The field norm map from $L$ to $K$. \\
\end{tabular}

\medskip
\textit{Sums, matrices, and operations:}
\par\smallskip
\begin{tabular}{@{} r @{\quad} l}
 $J_{p^r}^*(a,b)$: & The modified Jacobi sum for the finite field $\mathbb{F}_{p^r}$. \\
 $E_{\mathbf{p}}(a,b)$: & Thaine's generalized Jacobi sum for an element $\mathbf{p} \in K^\times$. \\
 $C_{p^r}$: & The multiplication matrix of Gaussian periods. \\
 ${\rm Cyc}_{p^r}$: & The matrix of cyclotomic numbers. \\
 $\overset{d}{\ast}$: & Thaine's $d$-composition operation (for matrices or points on $V$). \\
\end{tabular}

%------------------------------------S2
\section{Preliminaries on Jacobi sums and Thaine's framework}\label{S2}
\subsection{Classical Jacobi sums and multiplication matrices}\label{S2-1}
We recall some basic definitions, e.g. 
Gaussian periods, period polynomials, Jacobi sums and cyclotomic numbers, 
multiplication matrices of Gaussian periods, 
and their relations (see Berndt, Evans and Williams \cite{BEW98}). 

Let $e \geq 2$ be an integer and 
$p^r$ be a prime power with $p^r \equiv 1\ ({\rm mod}\ e)$. 
Write $p^r=ef+1$. 
Let $\F_{p^r}$ be the finite field of $p^r$ elements 
and $\gamma$ be a fixed generator of 
$\F_{p^r}^{\times}=\F_{p^r}\setminus \{0\}$. 
Let $\zeta_n=e^{2\pi i/n}$ be a primitive $n$-th root of unity. 
For $0 \leq i \leq e-1$, 
{\it Gaussian $f$-periods $\eta_{p^r}(i)$ of degree $e$ for $\F_{p^r}$} 
are defined by
\begin{align*}
\eta_{p^r}(i) := \sum^{f-1}_{j=0}\zeta_p^{\mathrm{Tr}(\gamma^{ej+i})}
\end{align*}
where $\mathrm{Tr}$ is the trace map $\mathrm{Tr} : \F_{p^r} \to \F_p$, 
and {\it the period polynomial $P_{e,p^r}(X)$ of degree $e$ for $\F_{p^r}$} 
is defined by $P_{e,p^r}(X)= \prod^{e-1}_{i=0}\left(X-\eta_{p^r}(i)\right)\in\Z[X]$. 
Note that $\eta_{p^r}(i)$ depends on the choice of $\gamma$. 

For a nontrivial character $\psi$ on $\F_{p^r}^{\times}$ and 
the trivial character $\varepsilon$ on $\F_{p^r}^{\times}$, 
we extend them to $\F_{p^r}$ by setting $\psi(0)=0$ and $\varepsilon(0)=1$. 
Let $\psi_1$, $\psi_2$ be multiplicative characters on $\F_{p^{r}}^{\times}$, extended to $\F_{p^r}$. 
{\it Jacobi sums $J_{p^r}^{\ast}(\psi_1,\psi_2)$ for $\F_{p^r}$} are defined by
\[
J_{p^r}^{\ast}(\psi_1,\psi_2):=\sum_{\substack{\alpha \in \F_{p^r}\\ \alpha \neq 0,1}}\psi_1(\alpha)\psi_2(1-\alpha).
\]
If $\psi_1, \psi_2 \neq \varepsilon$, then 
$J_{p^r}^{\ast}(\varepsilon,\psi_2)=J_{p^r}^{\ast}(\psi_1, \varepsilon)=-1$ and 
$J_{p^r}^{\ast}(\varepsilon, \varepsilon)=p^r-2$.

From now on, we take the character $\chi$ of order $e$ 
on $\F_{p^r}$ with $\chi(\gamma)=\zeta_e$ and $\chi(0)=0$ 
where $\langle\gamma\rangle=\F_{p^r}^\times$. 
For $0 \leq a,b \leq e-1$, 
we simply write 
\[
J_{p^r}^{\ast}(a,b):=J_{p^r}^{\ast}(\chi^a,\chi^b).
\]
Davenport and Hasse's lifting theorem for Jacobi sums claims that 
for any integer $n \geq 1$, if $e$ $\notmid$ $a$, $e$ $\notmid$ $b$ and $e$ $\notmid$ $a+b$ , then
\begin{align*}
J_{p^{nr}}^{\ast}((\chi^{\prime})^a, (\chi^{\prime})^b)=(-1)^{n-1}J_{p^r}^{\ast}(a, b)^n
\end{align*}
where $\chi^{\prime}$ is the lift of $\chi$ 
from $\F_{p^r}$ to $\F_{p^{nr}}$ defined by 
$\chi^{\prime}(\alpha) = \chi(\mathrm{Nr}(\alpha))$ 
and $\mathrm{Nr}$ is the norm map 
(see Davenport and Hasse \cite[Relation (0.8)]{DH35}, 
see also Weil \cite[pages 503--506]{Wei49}, \cite[pages 360--362, Theorem~11.5.2, Corollary~11.5.3]{BEW98}, Hoshi and Kanai \cite[Section 1]{HK22}). 

For $0 \leq i,j \leq e-1$, 
{\it cyclotomic numbers ${\rm Cyc}_{p^r}(i,j)$ of order $e$ for $\F_{p^r}$} 
are defined by
\begin{align*}
{\rm Cyc}_{p^r}(i,j) :=\# \{ (v_1, v_2) \mid 0 \leq v_1, v_2 \leq f-1, \ 1+\gamma^{ev_1+i} \equiv \gamma^{ev_2+j} \ ({\rm mod}\ p^r) \}.
\end{align*}
We have the following well-known 
relations between cyclotomic numbers ${\rm Cyc}_{p^r}(i,j)$ 
and Jacobi sums $J_{p^r}^{\ast}(a,b)$ 
(see \cite[page 79, Theorem~2.5.1]{BEW98}):
\begin{align}\label{eq:Jr-Cycr}
(-1)^{fa}\sum_{i=0}^{e-1} \sum_{j=0}^{e-1} {\rm Cyc}_{p^r}(i,j) \zeta_e^{ai+bj}=J_{p^r}^{\ast}(a,b)
\end{align}
and
\begin{align}\label{eq:Cycr-Jr}
\sum_{a=0}^{e-1} \sum_{b=0}^{e-1} (-1)^{fa} \zeta_e^{-(ai+bj)}J_{p^r}^{\ast}(a,b)=e^2{\rm Cyc}_{p^r}(i,j) 
\end{align}
Note that both ${\rm Cyc}_{p^r}(i,j)$ and $J_{p^r}^{\ast}(a,b)$ 
depend on a choice of a fixed generator $\gamma$ of $\F_{p^{r}}^{\times}$. 

We see that a product of Gaussian periods is represented 
by a linear combination of Gaussian periods again 
and these coefficients are given in terms of the cyclotomic numbers 
(see \cite[page 328, Lemma~10.10.2,  page 437, Exercise 12.23]{BEW98}):

\begin{align*}
\eta_{p^r}(m)\eta_{p^r}(m+i)=\sum^{e-1}_{j=0} ({\rm Cyc}_{p^r}(i,j) -D_if)\eta_{p^r}(m+j)
\end{align*}
where $p\neq 2$ and $D_i=\delta_{0,i}$ (resp. $\delta_{e/2 ,i}$) if $f$ is even
(resp. $f$ is odd) and $\delta_{i,j}$ is Kronecker's delta.
For $p=2$, we should adopt $D_i=\delta_{0,i}$.
Note that if $p=2$, then $f$ and $e$ are odd.

It follows that the Gaussian periods $\eta_{p^r}(0),\ldots,\eta_{p^r}(e-1)$ 
are the eigenvalues of the $e \times e$ matrix 
$C_{p^r} = [C_{p^r}[i,j]]_{0\leq i,j\leq e-1} 
:= [{\rm Cyc}_{p^r}(i,j) - D_if]_{0\leq i,j\leq e-1}$ 
which is 
called {\it the multiplication matrix of Gaussian periods 
$\eta_{p^r}(0),\ldots,\eta_{p^r}(e-1)$}. 
Hence the period polynomial $P_{e,p^r}(X)$ can be obtained 
as the characteristic polynomial ${\rm Char}_X(C_{p^r})$ of 
the multiplication matrix $C_{p^r}$. 
By the equations \eqref{eq:Cycr-Jr} and \eqref{eq:Jr-Cycr}, 
we have the following 
relations between $C_{p^r}[i,j]$ 
and  $J_{p^r}^{\ast}(a,b)$: 
\begin{align*}
C_{p^r}[i,j] = -fD_i + \frac{1}{e^2}\sum_{a=0}^{e-1} \sum_{b=0}^{e-1} (-1)^{fa} \zeta_e^{-(ai+bj)}J_{p^r}^{\ast}(a,b)
\end{align*}
and
\begin{align*}
J_{p^r}^{\ast}(a,b)= (p^r-1)\delta_{0,b} + (-1)^{fa}\sum_{i=0}^{e-1} \sum_{j=0}^{e-1} C_{p^r}[i,j] \zeta_e^{ai+bj}. 
\end{align*}

\subsection{Thaine's generalized Jacobi sums and multiplication matrices}

F. Thaine investigated various properties and characterizations of 
Gaussian periods, cyclotomic numbers and Jacobi sums 
with applications to the construction of cyclic polynomials 
in his series of papers \cite{Tha96}, \cite{Tha99}, 
\cite{Tha00}, \cite{Tha01}, \cite{Tha04}, \cite{Tha08} 
(see also Lehmer \cite{Leh88}, 
Schoof and Washington \cite{SW88}, 
Hashimoto and Hoshi \cite{HH05a}, \cite{HH05b}). 

Thaine \cite[Section 3]{Tha04} introduced generalized Jacobi sums $E_p(a,b)$ and generalized multiplication matrices $H_p$ for ${\bf p}\in K^\times$ by axiomatizing the formal properties of the classical objects recalled in Section \ref{S2-1}. 
We briefly review this framework and record the results that will be used in the proof of Theorem~\ref{mainth1}.

We first recall some basic properties of Jacobi sums and 
the multiplication matrix of Gaussian periods.
%
\begin{prop}[{see Thaine \cite[Proposition~7]{Tha04}}]\label{prop:T7}
Let $e \geq 2$ be an  
integer and $p^r$ be a prime power with $p^r\equiv 1\ ({\rm mod}\ e)$. 
Define $f:=(p^r-1)/e$. 
For all $a,b,c \in \Z$,
\begin{enumerate}
  \setlength{\parskip}{0cm}
  \setlength{\itemsep}{0cm}
\item $J^\ast_{p^r}(a+e,b)=J^\ast_{p^r}(a,b+e)=J^\ast_{p^r}(a,b)$,
\item $J^\ast_{p^r}(a,b)=J^\ast_{p^r}(b,a)$,
\item $J^\ast_{p^r}(a,b)=(-1)^{fb}J^\ast_{p^r}(-a-b,b)$,
\item $J^\ast_{p^r}(a,0)=-1$ if $e$ $\notmid$ $a$,
\item $J^\ast_{p^r}(a,b)J^\ast_{p^r}(-a,-b) = p^r$ 
if $e$ $\notmid$ $a$, $e$ $\notmid$ $b$ and $e$ $\notmid$ $a+b$,
\item $J^\ast_{p^r}(a,b)J^\ast_{p^r}(-a,-c) = (-1)^{f(b+c)}J^\ast_{p^r}(-(a+b+c),b)J^\ast_{p^r}(a+b+c,-c)$ if $e$ $\notmid$ $a+b$, 
$e$ $\notmid$ $a+c$ and $e$ $\notmid$ $a+b+c$,
\item if ${\rm gcd}(c,e)=1$ then $\sigma_c (J^\ast_{p^r}(a,b))=J^\ast_{p^r}(ca,cb)$ where $\sigma_c \in \Gal(\Q(\zeta_e)/\Q)$ such that $\sigma_c(\zeta_e)=\zeta_e^c$.
\end{enumerate}
\end{prop}

\begin{remark}\label{rem:DefJacobi}
Thaine \cite{Tha04} uses a different definition of Jacobi sums 
\[
J_{a,b} := -\sum_{k=2}^{p^r-1}\zeta_e^{a\ {\rm ind}_{s}(k)+b\ {\rm ind}_{s}(1-k)}
\]
where  
$s$ is a primitive root modulo $p^r$ such that $s^f \equiv \zeta_e \pmod{P}$ 
($P$ is a prime ideal of $\Z[\zeta_e]$ above $p^r$) and 
${\rm ind}_{s}(k)$ is the least non-negative integer such that 
$s^{{\rm ind}_s(k)}\equiv k \pmod{p^r}$. 
For our notation $J^\ast_{p^r}(a,b)$, 
if we take the same primitive root $s$ modulo $p^r$, 
then $J_{a,b}=-J^\ast_{p^r}(a,b)$.
\end{remark}

We take an $e \times e$ circulant matrix $\mathcal{K}=[\delta_{i+1,j}]_{0\leq i,j\leq e-1}$, i.e. 
\begin{align*}
\mathcal{K} =
\left(
\begin{array}{cccc}
 & 1 & &  \\
 & &  \ddots &  \\
 & &  & 1 \\
1 & & & 
\end{array}\right). 
\end{align*}

\begin{prop}[{Thaine \cite[Proposition~8]{Tha04}}]\label{prop:T8}
Let $e \geq 2$ be an integer and 
$p^r$ be a prime power with $p^r\equiv 1\ ({\rm mod}\ e)$. 
Define $f:=(p^r-1)/e$.
Let $C_{p^r}=[c_{i,j}]_{0\leq i,j \leq e-1}$ be the multiplication matrix of the Gaussian periods $\eta_{p^r}(0),\ldots,\eta_{p^r}(e-1)$
with $c_{i,j} \in \Z$ and $\eta_{p^r}(0)\eta_{p^r}(i)=\sum_{j=0}^{e-1}c_{i,j}\eta_{p^r}(j)$. 
Then the elements $c_{i,j}$ satisfy the following conditions: 
\renewcommand{\labelenumi}{\textnormal{(\roman{enumi})}}
\begin{enumerate}
  \setlength{\parskip}{0cm}
  \setlength{\itemsep}{0cm}
\item $c_{i+e,j}=c_{i,j+e}=c_{i,j}$,
\item $c_{i,j}=c_{j+(p^r-1)/2,i+(p^r-1)/2} + f(\delta_{0,j}-\delta_{ef/2, i})$,
\item $\sum^{e-1}_{k=0}c_{i,k}=f-p^r\delta_{ef/2, i}$,
\item $\sum^{e-1}_{k=0}c_{k,j}=-\delta_{0,j}$,
\item $c_{i,j}=c_{-i,j-i}$,
\item $\sum^{e-1}_{k=0}c_{i,k}c_{k-j,l-j}=\sum^{e-1}_{k=0}c_{j,k}c_{k-i,l-i}$,
\item  for all $k \in \Z$, $C_{p^r}(\mathcal{K}^{-k}C_{p^r}\mathcal{K}^{k})=(\mathcal{K}^{-k}C_{p^r}\mathcal{K}^{k})C_{p^r}$ where $\mathcal{K}=[\delta_{i+1,j}]_{0\leq i,j\leq e-1}$ is an $e \times e$ circulant matrix.
\end{enumerate}
\end{prop}

Thaine \cite{Tha04} 
defines generalized Jacobi sums $E_{a,b}$ with $E_{0,0}=1, E_{a,0}=1$ 
and generalized multiplication matrices $H_{\bf p}=[h_{i,j}]_{0\leq i,j \leq e-1}\in M_e(K)$ for ${\bf p}\in K^\times$. 
However, we will use the modified version 
$E_{\bf p}(a,b)$ of $E_{a,b}$ 
with $E_{\bf p}(0,0)={\bf p}-2$, $E_{\bf p}(a,0)=-1$ as follows:

\begin{defn}[{see Thaine \cite[Proposition~9 and Proposition~10]{Tha04}}]\label{def:T9and10}
Let $e \geq 2$ be an integer
and $v$ be an integer with $ve$ even.
Let $K$ be a field with char $K=0$ and ${\bf p}\in K^{\times}$. 
We consider a map $E_{\bf p}: \Z \times \Z \longrightarrow K(\zeta_e)$ which satisfies the following conditions:
\begin{enumerate}
  \setlength{\parskip}{0cm}
  \setlength{\itemsep}{0cm}
\item $E_{\bf p}(a+e,b)=E_{\bf p}(a,b+e)=E_{\bf p}(a,b)$,
\item $E_{\bf p}(a,b)=E_{\bf p}(b,a)$,
\item $E_{\bf p}(a,b)=(-1)^{vb}E_{\bf p}(-a-b,b)$,
\item $E_{\bf p}(a,0)=-1$ if $e$ $\notmid$ $a$,
\item $E_{\bf p}(a,b)E_{\bf p}(-a,-b) = {\bf p}$ 
if $e$ $\not{|}$ $a$, $e$ $\notmid$ $b$ and $e$ $\notmid$ $a+b$,
\item $E_{\bf p}(a,b)E_{\bf p}(-a,-c) 
= (-1)^{v(b+c)}E_{\bf p}(-(a+b+c),b)E_{\bf p}(a+b+c,-c)$ \\
if $e$ $\notmid$ $a+b$, $e$ $\notmid$
$a+c$ 
and $e$ $\notmid$ $a+b+c$,
\item if ${\rm gcd}(c,e)=1$, then $\sigma_c (E_{\bf p}(a,b))=E_{\bf p}(ca,cb)$.
\end{enumerate}
We call these elements $E_{\bf p}(a,b) \in K(\zeta_e)$ {\it generalized Jacobi sums  for ${\bf p}$}. 
\end{defn}

\begin{defn}[{see Thaine \cite[Proposition~9 and Proposition~10]{Tha04}}]\label{def:T9and10_2}
Let $e \geq 2$ be an integer
and $v$ be an integer with $ve$ even.
Let $K$ be a field with char $K=0$ and 
${\bf p}\in K^{\times}$. 
Define $f:=({\bf p}-1)/e$. 
A matrix $H_{\bf p}=[h_{i,j}]_{0\leq i,j\leq e-1} \in M_e(K)$
is called a {\it generalized multiplication matrix for ${\bf p}$} 
if it satisfies the following conditions:
\begin{enumerate}
\item[(i)] $h_{i+e,j}=h_{i,j+e}=h_{i,j}$,
\item[(iii)] $\sum^{e-1}_{k=0}h_{i,k}=f-{\bf p}\delta_{i,\frac{ev}{2}}$,
\item[(v)] $h_{i,j}=h_{-i,j-i}$,
\item[(vii)]  for all $k \in \Z$, $H_{\bf p}(\mathcal{K}^{-k}H_{\bf p}\mathcal{K}^{k})
=(\mathcal{K}^{-k}H_{\bf p}\mathcal{K}^{k})H_{\bf p}$ where 
$\mathcal{K}=[\delta_{i+1,j}]_{0\leq i,j\leq e-1}$ 
is an $e \times e$ circulant matrix.
\end{enumerate}
\end{defn}

\begin{remark}\label{rem:DefGenJacobi}
In Thaine \cite[Proposition~9 and Proposition~10]{Tha04}, 
the conditions {\rm (i), (iii), (v), (vii)} in Definition \ref{def:T9and10_2} 
are considered together with
\begin{enumerate}
\item[(ii)] $h_{i,j}=h_{j+ve/2,i+ve/2} + f(\delta_{0,j}-\delta_{i,\frac{ev}{2}})$,
\item[(iv)] $\sum^{e-1}_{k=0}h_{k,j}=-\delta_{0,j}$,
\item[(vi)] $\sum^{e-1}_{k=0}h_{i,k}h_{k-j,l-j}
=\sum^{e-1}_{k=0}h_{j,k}h_{k-i,l-i}$
\end{enumerate}
where $v$ is an integer with $ve$ even. 
However, the conditions {\rm (ii), (iv), (vi)} follow from
{\rm (i), (iii), (v), (vii)} 
as mentioned in Thaine \cite[page 271, Observation of Proposition~9]{Tha04}. 
\end{remark}

The following result, which is a version in terms of $E_p(a,b)$ rather than $E_{a,b}$, is due to Thaine~\cite{Tha04}:

\begin{prop}[{Thaine \cite[Proposition~9 and Proposition~10]{Tha04}}]\label{prop:T9andT10}
Let $e \geq 2$ be an integer and $v$ be an integer with $ve$ even. 
Let $K$ be a field with char $K=0$ and ${\bf p}\in K^{\times}$. 
Define $f:=({\bf p}-1)/e$. 
Let $E_{\bf p}(a,b)$ be generalized Jacobi sums for ${\bf p}$. 
For $i,j \in \Z$, we define
\begin{equation}\label{eq:3.1}
h_{i,j}=-f\delta_{i,\frac{ev}{2}}+\frac{1}{e^2}\sum^{e-1}_{a=0} \sum^{e-1}_{b=0} (-1)^{va} \zeta_e^{-(ai+bj)}E_{\bf p}(a,b). \tag{A}
\end{equation}
Then $H_{\bf p}=[h_{i,j}]_{0\leq i,j\leq e-1}$ 
becomes a generalized multiplication matrix for ${\bf p}$
which satisfies 
\begin{equation}\label{eq:3.2}
E_{\bf p}(a,b)=({\bf p}-1)\delta_{0,b} + (-1)^{va}\sum^{e-1}_{i=0} \sum^{e-1}_{j=0} \zeta_e^{ai+bj}h_{i,j}. \tag{B}
\end{equation}
Conversely, if $H_{{\bf p}}=[h_{i,j}]_{0\leq i,j\leq e-1}$ 
is a generalized multiplication matrix for ${\bf p}$ 
and $E_{\bf p}(a,b)$ is defined by \eqref{eq:3.2}, then 
$E_{\bf p}(a,b)$ are generalized Jacobi sums for ${\bf p}$ 
which satisfy \eqref{eq:3.1}. 
\end{prop}

%------------------------------------S3
\section{{Proof of Theorem~\ref{mainth1}} and {Theorem~\ref{thm:corMT1}: Multiplicative identities of generalized Jacobi sums}}\label{S3}

In this section, we prove Theorem~\ref{mainth1} and Theorem~\ref{thm:corMT1} and then interpret Theorem~\ref{mainth1} in terms of a twisted double finite Fourier transform. 
The last viewpoint also gives an alternative proof of Theorem~\ref{mainth1} in a restricted case.
\subsection{Proof of Theorem~\ref{mainth1}}

(1) Define
\[
E_{{\bf p}_1{\bf p}_2,d}(a,b)
:=
-\sigma_{-d}(E_{{\bf p}_1}(a,b))E_{{\bf p}_2}(a,b).
\]
We verify that $E_{{\bf p}_1{\bf p}_2,d}(a,b)$ satisfies the conditions
in Definition~\ref{def:T9and10} with respect to the parameter $v=v_1+v_2$.

First, condition (1) follows immediately from the periodicity of
$E_{{\bf p}_1}$ and $E_{{\bf p}_2}$:
\[
E_{{\bf p}_1{\bf p}_2,d}(a+e,b)
=
E_{{\bf p}_1{\bf p}_2,d}(a,b),
\qquad
E_{{\bf p}_1{\bf p}_2,d}(a,b+e)
=
E_{{\bf p}_1{\bf p}_2,d}(a,b).
\]

For condition (2), using the symmetry of $E_{{\bf p}_1}$ and
$E_{{\bf p}_2}$, we have
\[
\begin{aligned}
E_{{\bf p}_1{\bf p}_2,d}(a,b)
&=
-\sigma_{-d}(E_{{\bf p}_1}(a,b))E_{{\bf p}_2}(a,b)\\
&=
-\sigma_{-d}(E_{{\bf p}_1}(b,a))E_{{\bf p}_2}(b,a)\\
&=
E_{{\bf p}_1{\bf p}_2,d}(b,a).
\end{aligned}
\]

For condition (3), we compute
\[
\begin{aligned}
E_{{\bf p}_1{\bf p}_2,d}(a,b)
&=
-\sigma_{-d}(E_{{\bf p}_1}(a,b))E_{{\bf p}_2}(a,b)\\
&=
-\sigma_{-d}\bigl((-1)^{v_1b}E_{{\bf p}_1}(-a-b,b)\bigr)
  (-1)^{v_2b}E_{{\bf p}_2}(-a-b,b)\\
&=
(-1)^{(v_1+v_2)b}
E_{{\bf p}_1{\bf p}_2,d}(-a-b,b).
\end{aligned}
\]

For condition (4), if $e\nmid a$, then
\[
\begin{aligned}
E_{{\bf p}_1{\bf p}_2,d}(a,0)
&=
-\sigma_{-d}(E_{{\bf p}_1}(a,0))E_{{\bf p}_2}(a,0)\\
&=
-\sigma_{-d}(-1)(-1)\\
&=
-1.
\end{aligned}
\]

For condition (5), assume that $e\nmid a$, $e\nmid b$, and
$e\nmid a+b$. Then
\[
\begin{aligned}
&E_{{\bf p}_1{\bf p}_2,d}(a,b)
 E_{{\bf p}_1{\bf p}_2,d}(-a,-b)\\
&\quad=
\bigl(-\sigma_{-d}(E_{{\bf p}_1}(a,b))E_{{\bf p}_2}(a,b)\bigr)
\bigl(-\sigma_{-d}(E_{{\bf p}_1}(-a,-b))E_{{\bf p}_2}(-a,-b)\bigr)\\
&\quad=
\sigma_{-d}\bigl(E_{{\bf p}_1}(a,b)E_{{\bf p}_1}(-a,-b)\bigr)
E_{{\bf p}_2}(a,b)E_{{\bf p}_2}(-a,-b)\\
&\quad=
{\bf p}_1{\bf p}_2.
\end{aligned}
\]

For condition (6), assume that $e\nmid a+b$, $e\nmid a+c$, and
$e\nmid a+b+c$. Then
\[
\begin{aligned}
&E_{{\bf p}_1{\bf p}_2,d}(a,b)
 E_{{\bf p}_1{\bf p}_2,d}(-a,-c)\\
&\quad=
\sigma_{-d}\bigl(E_{{\bf p}_1}(a,b)E_{{\bf p}_1}(-a,-c)\bigr)
E_{{\bf p}_2}(a,b)E_{{\bf p}_2}(-a,-c)\\
&\quad=
(-1)^{v_1(b+c)}
\sigma_{-d}\bigl(
E_{{\bf p}_1}(-(a+b+c),b)E_{{\bf p}_1}(a+b+c,-c)
\bigr)\\
&\qquad\qquad\times
(-1)^{v_2(b+c)}
E_{{\bf p}_2}(-(a+b+c),b)E_{{\bf p}_2}(a+b+c,-c)\\
&\quad=
(-1)^{(v_1+v_2)(b+c)}
E_{{\bf p}_1{\bf p}_2,d}(-(a+b+c),b)
E_{{\bf p}_1{\bf p}_2,d}(a+b+c,-c).
\end{aligned}
\]

Finally, for condition (7), let $\gcd(c,e)=1$. Since $\sigma_c$ and
$\sigma_{-d}$ commute in $\operatorname{Gal}(K(\zeta_e)/K)$, we have
\[
\begin{aligned}
\sigma_c(E_{{\bf p}_1{\bf p}_2,d}(a,b))
&=
-\sigma_c\sigma_{-d}(E_{{\bf p}_1}(a,b))
 \sigma_c(E_{{\bf p}_2}(a,b))\\
&=
-\sigma_{-d}\sigma_c(E_{{\bf p}_1}(a,b))
 E_{{\bf p}_2}(ca,cb)\\
&=
-\sigma_{-d}(E_{{\bf p}_1}(ca,cb))
 E_{{\bf p}_2}(ca,cb)\\
&=
E_{{\bf p}_1{\bf p}_2,d}(ca,cb).
\end{aligned}
\]
Thus $E_{{\bf p}_1{\bf p}_2,d}(a,b)$ are generalized Jacobi sums for ${\bf p}_1{\bf p}_2$.


(2) Putting $i^{\prime}:=ds+i$ and $j^{\prime}:=dt+j$, we have
\begin{align*}
\sum_{i=0}^{e-1}\sum_{j=0}^{e-1}\zeta_e^{ai+bj}({\rm Cyc}_{{\bf p}_1}\overset{d}{\ast}{\rm Cyc}_{{\bf p}_2})[i,j]
&=\sum_{i=0}^{e-1}\sum_{j=0}^{e-1}\zeta_e^{ai+bj}
\sum^{e-1}_{s=0}\sum^{e-1}_{t=0}{\rm Cyc}_{{\bf p}_1}[s,t]{\rm Cyc}_{{\bf p}_2}[ds+i,dt+j]\\
&=\sum_{i^\prime=0}^{e-1}\sum_{j^\prime=0}^{e-1}\zeta_e^{a(i^\prime-ds)+b(j^\prime-dt)}
\sum^{e-1}_{s=0}\sum^{e-1}_{t=0}{\rm Cyc}_{{\bf p}_1}[s,t]{\rm Cyc}_{{\bf p}_2}[i^\prime,j^\prime]\\
&=\left(\sum^{e-1}_{s=0}\sum^{e-1}_{t=0}\zeta_e^{-(ads+bdt)}{\rm Cyc}_{{\bf p}_1}[s,t]\right)
\left(\sum_{i^\prime=0}^{e-1}\sum_{j^\prime=0}^{e-1}\zeta_e^{ai^\prime+bj^\prime}{\rm Cyc}_{{\bf p}_2}[i^\prime,j^\prime]\right).
\end{align*}
Multiplying both sides by $(-1)^{(v_1+v_2)a+1}$, we get
\begin{align*}
&(-1)^{(v_1+v_2)a+1}\sum_{i=0}^{e-1}\sum_{j=0}^{e-1}\zeta_e^{ai+bj}({\rm Cyc}_{{\bf p}_1}\overset{d}{\ast}{\rm Cyc}_{{\bf p}_2})[i,j]\\
&=-\left((-1)^{v_1a}\sum^{e-1}_{s=0}\sum^{e-1}_{t=0}\zeta_e^{-(ads+bdt)}{\rm Cyc}_{{\bf p}_1}[s,t]\right)
\left((-1)^{v_2a}\sum_{i^\prime=0}^{e-1}\sum_{j^\prime=0}^{e-1}\zeta_e^{ai^\prime+bj^\prime}{\rm Cyc}_{{\bf p}_2}[i^\prime,j^\prime]\right)\\
&=-E_{{\bf p}_1}(-da,-db)E_{{\bf p}_2}(a,b)\\
&=-\sigma_{-d}(E_{{\bf p}_1}(a,b))E_{{\bf p}_2}(a,b)\\
&=E_{{{\bf p}_1{\bf p}_2},d}(a,b). 
\end{align*}

For ${\bf p}_1=p^r$ and ${\bf p}_2=q^s$ prime powers, and $v_1=f_1=(p^r-1)/e, v_2=f_2=(q^s-1)/e$,
we have $f=f_1f_2e+f_1+f_2$ and $(-1)^{fa}=(-1)^{f_1f_2ea}(-1)^{(f_1+f_2)a}$. 
If $p>2$ or $q>2$, then $f_1f_2e$ is even. 
Thus we have $(-1)^{(f_1+f_2)a}=(-1)^{fa}$. 
If $p=q=2$, then $f_1, f_2, e$ are odd. 
Hence $f_1+f_2$ is even and $(-1)^{(f_1+f_2)a}=1$.
\qed\\

\subsection{Proof of Theorem~\ref{thm:corMT1}}

(1) $\Rightarrow$ (2) is proved in 
Hoshi and Kanai \cite[Theorem~1.4]{HK22}.\\

(2) $\Rightarrow$ (1). 
%
In Theorem~\ref{mainth1} (2), if we take ${\bf p}_1={\bf p}_2=p^r$ 
and $d=-1$, then we have
\begin{align*}
(-1)^{f a +1} \sum_{i=0}^{e-1}\sum_{j=0}^{e-1}\zeta_e^{ai+bj}({\rm Cyc}_{p^r}\overset{-1}{\ast}{\rm Cyc}_{p^r})[i,j] = -J^\ast_{p^r}(a,b)^2. 
\end{align*}
We define $D_{i}:=[\delta_{ef/2 ,i}]_{0\leq i, j \leq e-1}$. 
By ${\rm Cyc}_{p^r} = C_{p^r} + D_{i}f$, we have
\begin{align*}
{\rm Cyc}_{p^r}\overset{-1}{\ast}{\rm Cyc}_{p^r}
&= (C_{p^r} + D_{i}f) \overset{-1}{\ast} (C_{p^r} + D_{i}f)\\
&= C_{p^r}^{(2)} + 2f(C_{p^r} \overset{-1}{\ast} D_{i}) + f^2(D_{i} \overset{-1}{\ast} D_{i})\\
&=-C_{p^{2r}} + 2f([f]_{0\leq i, j \leq e-1}-p^rD_i) + f^2(eD_{i})\\
&=-C_{p^{2r}} - \frac{1}{e}(p^{2r}-1)D_{i} + [2f^2]_{0\leq i, j \leq e-1}
\end{align*}
where the third equality follows from 
the assumption
(2) $C_{p^{2r}} = -C_{p^r}^{(2)}$. 
Because 
\begin{align*}
(-1)^{f a +1} \sum_{i=0}^{e-1}\sum_{j=0}^{e-1}\zeta_e^{ai+bj}\left(- \frac{1}{e}(p^{2r}-1)D_{i}[i,j]\right)
&=\frac{1}{e}(p^{2r}-1)(-1)^{f a}\sum_{i=0}^{e-1}\sum_{j=0}^{e-1}\zeta_e^{ai+bj}D_{i}[i,j]\\
&=\frac{1}{e}(p^{2r}-1)(-1)^{f a}\sum_{j=0}^{e-1}\zeta_e^{a\frac{ef}{2}+bj}\\
&=\frac{1}{e}(p^{2r}-1)(-1)^{f a}\zeta_e^{a\frac{ef}{2}}\sum_{j=0}^{e-1}\zeta_e^{bj}\\
&=(p^{2r}-1)\delta_{0,b}
\end{align*}
and
\begin{align*}
(-1)^{f a +1} \sum_{i=0}^{e-1}\sum_{j=0}^{e-1}\zeta_e^{ai+bj}2f^2
=2f^2(-1)^{f a +1} \sum_{i=0}^{e-1}\sum_{j=0}^{e-1}\zeta_e^{ai+bj}
=-2e^2f^2\delta_{a,0}\delta_{0,b},
\end{align*}
we obtain that 
\begin{align*}
-J^\ast_{p^r}(a,b)^2
&=(-1)^{f a +1} \sum_{i=0}^{e-1}\sum_{j=0}^{e-1}\zeta_e^{ai+bj}({\rm Cyc}_{p^r}\overset{-1}{\ast}{\rm Cyc}_{p^r})[i,j]\\
&=(-1)^{f a +1} \sum_{i=0}^{e-1}\sum_{j=0}^{e-1}\zeta_e^{ai+bj}((-1)C_{p^{2r}}[i,j] - \frac{1}{e}(p^{2r}-1)\delta_{ef/2 ,i} + 2f^2)\\
&=-2e^2f^2\delta_{a,0}\delta_{0,b} + (p^{2r}-1)\delta_{0,b} + (-1)^{f a}\sum_{i=0}^{e-1}\sum_{j=0}^{e-1}\zeta_e^{ai+bj}C_{p^{2r}}[i,j]\\
&= -2e^2f^2\delta_{a,0}\delta_{0,b} + (p^{2r}-1)\delta_{0,b}  + J_{p^{2r}}^{\ast}((\chi^{\prime})^a, (\chi^{\prime})^b).
\end{align*}
Hence,
for $a$, $b \neq 0$, we have
\[
J_{p^{2r}}^{\ast}((\chi^{\prime})^a, (\chi^{\prime})^b)=-J^\ast_{p^r}(a,b)^2. 
\]
By repeating the same process, we get the desired result:
\[
J_{p^{nr}}^{\ast}((\chi^{\prime})^a, (\chi^{\prime})^b)=(-1)^{n-1}J^\ast_{p^r}(a,b)^n. 
\]
\qed

\subsection{A twisted double finite Fourier transform}
%
In this subsection, we explain that Theorem~\ref{mainth1} may be viewed as a twisted version of the convolution theorem for the double finite Fourier transform relating cyclotomic numbers and Jacobi sums; see Hoshi and Kanai \cite[Section 4]{HK22}. 
In the restricted case $d=-1$ and $f$ even, the twist disappears, and the convolution theorem yields an alternative proof of Theorem~\ref{mainth1} (2).

We briefly recall the Fourier transform on finite abelian groups; see Terras \cite[Chapter 10]{Ter99}.

Let $G$ be a finite abelian group and 
\[
L^2(G) = \{f : G \to \C\}
\]
be the vector space of all $\C$-valued functions on $G$ with 
the inner product 
$\langle f,g\rangle=\sum_{x\in G}f(x)\overline{g(x)}$. 
Let $\widehat{G}={\rm Hom}(G,\C^{\times})$ be the dual 
of $G$. 
For $f\in L^2(G\times G)$, 
{\it the double finite Fourier transform $\mathscr{F}(f)=\widehat{f}\in L^2(\widehat{G}\times \widehat{G})$ of $f$} is defined to be 
\begin{align*}
(\mathscr{F}(f))(\chi, \psi) = \widehat{f}(\chi,\psi) =\sum_{x \in G}\sum_{y \in G} f(x,y)\overline{\chi(x)\psi(y)}.
\end{align*} 
Then $\mathscr{F}:L^2(G\times G)\rightarrow L^2(\widehat{G}\times \widehat{G})$ becomes a 
bijective linear transformation with the inverse 
\begin{align*}
(\mathscr{F}^{-1}(\widehat{f}))(x,y) = f(x,y) 
= \frac{1}{(\# G)^2}\sum_{\chi \in \widehat{G}}\sum_{\psi \in \widehat{G}}\widehat{f}(\chi,\psi)\chi(x)\psi(y).
\end{align*}
For $f, g\in L^2(G\times G)$, we define {\it the convolution $f \ast g\in L^2(G\times G)$ 
of $f$ and $g$} by
\[
(f \ast g)(x,y) =\sum_{a \in G}\sum_{b \in G} f(a,b)g(x-a,y-b).
\]
We have the compatibility of the convolution $\ast$ and 
the double finite Fourier transform $\mathscr{F}(f)=\widehat{f}$: 
\begin{align}\label{eq:comb}
\widehat{(f\ast g)}(\chi, \psi)=\widehat{f}(\chi, \psi)\widehat{g}(\chi, \psi). 
\end{align}
(see Terras \cite[page 168, Theorem~2]{Ter99}). 

From now on, we will apply the convolution theorem to our case.
We regard the cyclotomic numbers ${\rm Cyc}_{p^r}(i,j)$ of degree $e$ for $\F_{p^r}$ as the function ${\rm Cyc}_{p^r}:\Z/e\Z \times \Z/e\Z \rightarrow \C$, 
$(i,j)\mapsto {\rm Cyc}_{p^r}(i,j)$ 
and Jacobi sums $J_{p^r}^{\ast}(a,b)$ for $\F_{p^r}$ as the function 
$J_{p^r}^{\ast}:\widehat{\Z/e\Z} \times \widehat{\Z/e\Z} \rightarrow \C$, $(\chi^a, \chi^b) \mapsto J_{p^r}^{\ast}(a,b)$. 
We recall the well-known relations \eqref{eq:Jr-Cycr} and \eqref{eq:Cycr-Jr}: 
\begin{align*}
J_{p^r}^{\ast}(a,b)=(-1)^{fa}\sum_{i=0}^{e-1} \sum_{j=0}^{e-1} {\rm Cyc}_{p^r}(i,j) \zeta_e^{ai+bj},\ \quad
{\rm Cyc}_{p^r}(i,j)=\frac{1}{e^2}\sum_{a=0}^{e-1} \sum_{b=0}^{e-1} (-1)^{fa} \zeta_e^{-(ai+bj)}J_{p^r}^{\ast}(a,b).
\end{align*}
If $f$ is even, 
then we find that $J_{p^r}^{\ast}$ and ${\rm Cyc}_{p^r}$
are each other's double finite Fourier transform with some twist. 
Indeed, we have
\begin{align*}
(\mathscr{F}({\rm Cyc}_{p^r}))(\chi^a, \chi^b)
&=\sum_{i \in \Z/e\Z} \sum_{j \in \Z/e\Z} {\rm Cyc}_{p^r}(i,j) \overline{\chi^a(i)\chi^b(j)}\\
&=\sum_{i=0}^{e-1} \sum_{j=0}^{e-1} {\rm Cyc}_{p^r}(i,j) \zeta_e^{-(ai+bj)}\\
&=J_{p^r}^{\ast}(\chi^{-a},\chi^{-b}).\\
\end{align*}
We also have
\begin{align*}
(\mathscr{F}^{-1}(J_{p^r}^{\ast}))(i,j)
&=\frac{1}{e^2}\sum_{\psi_1 \in \widehat{\Z/e\Z}} \sum_{\psi_2 \in \widehat{\Z/e\Z}}J_{p^r}^{\ast}(\psi_1,\psi_2) \psi_1(i)\psi_2(j)\\
&=\frac{1}{e^2}\sum_{a=0}^{e-1} \sum_{b=0}^{e-1}J_{p^r}^{\ast}(\chi^a,\chi^b) \chi^a(i)\chi^b(j)\\
&=\frac{1}{e^2}\sum_{a=0}^{e-1} \sum_{b=0}^{e-1} \zeta_e^{ai+bj}J_{p^r}^{\ast}(a,b)\\
&={\rm Cyc}_{p^r}(-i,-j).
\end{align*}
Together, these relations form the two-dimensional analogue of the relations between Gaussian periods $\eta_{p^r}$
and Gauss sums $G^*_{p^r}$ given by
\begin{align*}
(\mathscr{F}(\eta_r))(\chi^a)=G_{r}^{\ast}(\chi^{-a}), \quad
(\mathscr{F}^{-1}(G_{r}^{\ast}))(i)=\eta_{r}(-i)
\end{align*}
where $\mathscr{F}$ is the finite Fourier transform on $L^2(\Z/e\Z)$.

Therefore, in the restricted case $d=-1$ and $f$ even, equation \eqref{eq:comb} gives another proof of Theorem~\ref{mainth1} (2):
\begin{align*}
-\sum_{i=0}^{e-1}\sum_{j=0}^{e-1}({\rm Cyc}_{p^r}\overset{-1}{\ast}{\rm Cyc}_{q^s})[i,j]\,\zeta_e^{ai+bj}
&=-\widehat{({\rm Cyc}_{p^r}\ast {\rm Cyc}_{q^s})}(\chi^{-a},\chi^{-b})\\
&=-\widehat{{\rm Cyc}_{p^r}}(\chi^{-a},\chi^{-b})\widehat{{\rm Cyc}_{q^s}}(\chi^{-a},\chi^{-b})\\
&=-J_{p^r}^{\ast}(a,b)J_{q^s}^{\ast}(a,b)\\
&=J^{*}_{p^rq^s,-1}(a,b). 
\end{align*}

%------------------------------------S4
\section{Proof of Theorem 1.4: Multiplicative quadratic forms on algebraic varieties}\label{S4}

This section is devoted to the quadratic case. 
We begin by recalling the classical theory of multiplicative quadratic forms and its extension to algebraic $K$-varieties. 
We then prove Theorem~\ref{mainth2}, which gives multiplicative quadratic forms of dimension $l-1$ on complete intersections of quadrics, and conclude with applications to explicit lifts of solutions of Diophantine systems.

\subsection{Multiplicative quadratic forms and the level of fields}\label{S4-1}

Let $K$ be a field with char $K\neq 2$. 
Hurwitz \cite{Hur1898} proved that if there exists an identity of the type 
\[
(x_1^2+\cdots+x_n^2)(y_1^2+\cdots+y_n^2)=z_1^2+\cdots+z_n^2
\]
where the $z_k$'s are bilinear forms of the independent variables 
$x_i$ and $y_j$ over $K$, then $n=1,2,4,8$.

In general, for a regular quadratic form 
$q(\mathbf{x}):=q(x_1,\ldots,x_n)=\sum_{i,j=1}^n a_{ij}x_ix_j$ 
of dimension $n$ over $K$, 
i.e. $a_{i,j}\in K$ and det\! $[a_{ij}]_{1\leq i,j\leq n}\neq 0$, 
(i) $q(\mathbf{x})$ is called {\it multiplicative} if there exists a formula
\begin{align}\label{eq:multi}
q(\mathbf{x})q(\mathbf{y})=q(\textbf{\textit{z}})
\end{align}
where the $x_i$ and $y_j$ are independent variables and 
$z_k\in K(\mathbf{x},\mathbf{y})$ 
and 
(ii) 
$q(\mathbf{x})$ is called {\it strictly multiplicative} 
if there exists a formula 
\eqref{eq:multi} with $z_k$ linear in $y_j$ over $K(\mathbf{x})$.
It is known that if $q(\mathbf{x})$ is isotropic, then 
$q(\mathbf{x})$ is  
multiplicative and in this case 
$q(\mathbf{x})$ is strictly multiplicative if and only if 
$q(\mathbf{x})$ is hyperbolic 
(see Pfister \cite[Chapter 2]{Pfi95}, Rajwade \cite[Chapter 12]{Raj93}). 
A quadratic form $q$ is called an $m$-fold Pfister form if
$q=\langle 1,a_1\rangle \otimes \cdots \otimes \langle 1,a_m\rangle$,
i.e. a tensor product of binary quadratic forms of the type $\langle 1,a_i\rangle$.
Pfister \cite{Pfi65b} proved the following theorem: 
%
\begin{thm}[Pfister\ \cite{Pfi65b}]\label{thm:Pfister} 
If $q$ is a Pfister form, then $q$ is strictly multiplicative. 
Conversely, if $q$ is an anisotropic multiplicative form over $K$, 
then $q$ must be a Pfister form. 
\end{thm}
%
Let $D_K(n)$ be the set of values in $K^{\times}$ represented by 
a sum of $n$ squares in $K$:  
\[
D_K(n)=\{\alpha\in K^{\times}\ |\ \alpha=\alpha_1^2+\cdots+\alpha_n^2,\ \alpha_j\in K\}.
\]
The \textit{level} (or {\it Stufe}) of a field $K$ is defined as 
$s(K):=\text{Inf}\{n\in\mathbf{N} \mid -1\in D_K(n)\}$. 
By Theorem~\ref{thm:Pfister}, we see that if $n$ is a power of $2$, 
then $D_K(n)$ becomes a multiplicative group. 
Using this, Pfister \cite{Pfi65a} proved the following remarkable theorem 
(see also Pfister \cite[page 28, Historical Note]{Pfi95}): 

\begin{thm}[{Pfister \cite[Satz 4, Satz 5]{Pfi65a}, see also Rajwade \cite[Chapter 2, Theorem~2.2]{Raj93}, Pfister \cite[Chapter 3, Theorem~1.3, Theorem~1.4]{Pfi95}}]
The level $s(K)$ of a field $K$ is, if finite, always a power of $2$. 
Conversely, every power of $2$ is the level $s(K)$ of some field $K$. 
\end{thm}

For multiplicative quadratic forms, 
see also Pfister \cite{Pfi66}, \cite{Pfi67}, 
Knebusch and Scharlau \cite{KS80}, Scharlau \cite{Sch85}, 
Shapiro \cite{Sha00}, Lam \cite{Lam05}. 

Hoshi \cite{Hos03} defined multiplicative quadratic forms 
on algebraic $K$-varieties as follows:

\begin{defn}[{Hoshi \cite[Section 2]{Hos03}}]
Let $K$ be a field with {\rm char} $K\neq 2$ and 
$V \subset K^n$ be an algebraic $K$-variety. 
We say a regular
quadratic form $q({\bf x})$ is {\it multiplicative on $V$} 
if there exists a bilinear map 
$\varphi :K^n \times K^n \rightarrow K^n$ such that
\[
\varphi(V \times V) \subset V \quad \mbox{and}\quad q({\bf x})q({\bf y})=q(\varphi({\bf x,y}))\mbox{ for any }{\bf x,y} \in V.
\]
\end{defn}

\begin{remark}
Throughout the paper, all varieties and morphisms are defined over a
base field $K$. 
In the definition above, we formulated the notion of a multiplicative form in terms of $K$-rational points $V(K)$, since our main interest lies in arithmetic applications. 
In our concrete situation, however, the operations $\overset{d}{\ast}$ on $V$ arising from Thaine's $d$-composition are given by polynomial expressions with coefficients in $K$. 
In particular, they define morphisms of varieties over $K$ and extend functorially to $F$-rational points for any field extension $F/K$. Thus one can equivalently view a multiplicative quadratic (or $f$-ic) form as a morphism
\[
  q,h \colon V \longrightarrow \mathbb{A}^1
\]
satisfying the corresponding multiplicative identities
\[
  q \circ \overset{d}{\ast} = m \circ (q \times q),
  \qquad
  h \circ \overset{d}{\ast} = m \circ (h \times h)
\]
where \(m:\mathbb A^1\times\mathbb A^1\to\mathbb A^1\) is the usual multiplication morphism. 
On the open subset where $q$ or $h$ is nonzero, this restricts to the multiplication map on ${\mathbb G}_m$.
\end{remark}

Hoshi \cite[Proposition~1]{Hos03} showed that an $m$-fold Pfister form is multiplicative on a variety $V$ defined as the intersection of $2^{m-1}-1$ quadrics (see also \cite[Example 2]{Hos03}). 
Moreover, Hoshi \cite[Theorem~4, Theorem~6]{Hos03} 
constructed {\it non}-Pfister forms of dimension $4$ and $6$ 
which are multiplicative on certain algebraic $K$-varieties. 
%

The next subsection extends the construction of multiplicative
quadratic forms on algebraic $K$-varieties in Hoshi~\cite{Hos03} to dimension $l-1$, where $l$ is an odd prime, and proves Theorem~\ref{mainth2} via the multiplicative identity of Theorem~\ref{mainth1}.

%------------------------------------S4-2
\subsection{Proof of Theorem~\ref{mainth2}}\label{S4-2}

We now prove Theorem~\ref{mainth2}. 
The key point is that Theorem~\ref{mainth1} yields explicit bilinear formulas compatible with the quadratic relations defining $V$, so that the Jacobi-sum identities translate into closure of $V$ under the corresponding composition law. 
We first isolate the identities needed for the quadratic construction, then prove that the resulting bilinear map preserves the defining equations of $V$, and finally deduce Theorem~\ref{mainth2}.

We begin with an elementary lemma:

\begin{lem} [{see \cite[Theorem~3.9.1]{BEW98}, \cite[Theorem~(i), (ii), page 186]{KR85a}}] \label{lem:BEW}
Let $l$ be an odd prime. 
Let $K$ be a field with char $K\neq 2$, $l$ and 
$\Gal(K(\zeta_l)/K)\simeq (\Z/l\Z)^{\times}$.  
Let ${\bf p}\in K^{\times}$ and
\[
\alpha=\sum_{k=0}^{l-1}a_k\zeta_l^k \quad (a_k \in K)\\
\]
with
\[
\alpha\overline{\alpha}= {\bf p}
\]
where $\overline{\alpha}=\sigma_{-1}(\alpha)$ and $\sigma_{-1} \in \Gal(K(\zeta_l)/K)$ 
satisfies $\sigma_{-1}(\zeta_l)=\zeta_l^{-1}$. 
Then we have
\begin{enumerate}
\item ${\bf p} =\sum_{k=0}^{l-1}a_k^2 - \sum_{k=0}^{l-1}a_ka_{k+1}$,
\item $\sum_{k=0}^{l-1}a_ka_{k+1}=\sum_{k=0}^{l-1}a_ka_{k+m}$ $(2\leq m\leq (l-1)/2)$ 
\end{enumerate}
where each subscript of the $a_k$'s should be taken modulo $l$. 
Moreover, we may assume that $a_0=0$ without loss of generality.
\end{lem}
\begin{proof}
Since char $K\neq l$, we have $1+\zeta_l+\cdots+\zeta_l^{l-1}=0$.
Hence, for any $c\in K$,
\[
\sum_{k=0}^{l-1}(a_k-c)\zeta_l^k
=\sum_{k=0}^{l-1}a_k\zeta_l^k - c(1+\zeta_l+\cdots+\zeta_l^{l-1})
=\alpha.
\]
Taking $c=a_0$, we may replace $(a_k)$ by $(a_k-a_0)$ and thus assume $a_0=0$.

By $\alpha\overline{\alpha}= {\bf p}$, we have 
\[
\left( a_1\zeta_l +  \cdots +a_{l-1}\zeta_l^{l-1}\right) \left(a_1\zeta_l^{l-1} + \cdots + a_{l-1}\zeta_l\right) = {\bf p},
\]
that is,
\[
S_0 + S_1\left(\zeta_l + \zeta_l^{l-1}\right) +  \cdots + S_{(l-1)/2}\left(\zeta_l^{(l-1)/2} + \zeta_l^{l-(l-1)/2}\right)= {\bf p}
\]
where $S_j = \sum_{i-k=j}a_ia_{k}$ ($j=0,\ldots,l-1$).
Note that  $S_j =S_{l-j}$.
By $\zeta_l+\cdots+\zeta_l^{l-1}=-1$, we have
\[
(S_1 - S_0 + {\bf p})\left(\zeta_l + \zeta_l^{l-1}\right) + 
 \cdots + (S_{(l-1)/2} - S_0 + {\bf p})\left(\zeta_l^{(l-1)/2} + \zeta_l^{l-(l-1)/2}\right) = 0.
\]
Because it follows from the assumption 
$\Gal(K(\zeta_l)/K)\simeq (\Z/l\Z)^{\times}$ that 
$\zeta_l,\zeta_l^2, \dots, \zeta_l^{l-1}$ are linearly independent over $K$, 
we have
\[
S_1 - S_0 + {\bf p} = S_m - S_0 + {\bf p}=0\ \ (2\leq m\leq (l-1)/2).
\]
Hence ${\bf p}=S_0-S_1$ and $S_1=S_m$ for $2\le m\le (l-1)/2$, which proves the lemma.
\end{proof}

Theorem~\ref{mainth2} follows from the following theorem: 

\begin{thm}\label{thm:BF}
Let $l$ be an odd prime and $d \in (\Z/l\Z)^{\times}$. 
Let $K$ be a field with char $K\neq 2$, $l$ 
and $\Gal(K(\zeta_l)/K)\simeq (\Z/l\Z)^{\times}$. 
Let
$E_{{\bf p}_1}(a,b)=\sum_{i=0}^{l-1}a_i\zeta_l^i$ and $E_{{\bf p}_2}(a,b)=\sum_{j=0}^{l-1}b_j\zeta_l^j$
be
generalized Jacobi sums for 
${\bf p}_1$, ${\bf p}_2\in K^\times$ respectively. 
For $a,b\in \Z$ with $l\nmid a$, $l\nmid b$, and $l\nmid a+b$,
let $E_{{\bf p}_1{\bf p}_2,d}(a,b)=-\sigma_{-d}(E_{{\bf p}_1}(a,b))E_{{\bf p}_2}(a,b)$
be a generalized Jacobi sum for ${\bf p}_1{\bf p}_2\in K^\times$ 
$($see Theorem~\ref{mainth1}$)$.
Then $E_{{\bf p}_1{\bf p}_2,d}(a,b)=\sum_{k=0}^{l-1}c_k\zeta_l^k$ 
is  given by
$$c_k=-\sum_{i=0}^{l-1}a_{-d^{-1}i}b_{k-i}$$
and the coefficients $c_{0},\dots,c_{l-1}$ satisfy
\begin{enumerate}
\item ${\bf p}_1{\bf p}_2=\sum_{k=0}^{l-1}c_k^2 - \sum_{k=0}^{l-1}c_kc_{k+1}$,
\item $\sum_{k=0}^{l-1}c_kc_{k+1}
=\sum_{k=0}^{l-1}c_kc_{k+m}$\ $(2\leq m\leq (l-1)/2)$
\end{enumerate}
with all indices taken modulo $l$.
Moreover, we may assume $a_0=b_0=c_0=0$ 
without loss of generality and in this case 
$E_{{\bf p}_1{\bf p}_2,d}(a,b)=\sum_{k=1}^{l-1}c_k\zeta_l^k$ 
is  given by 
$$c_k=-\left( \sum_{i=1}^{l-1}a_{-d^{-1}i}b_{k-i} - \sum_{i=1}^{l-1}a_{-d^{-1}i}b_{-i}\right).$$
\end{thm}
\begin{proof}
By Theorem~\ref{mainth1} (1), the element
\[
E_{{\bf p}_1{\bf p}_2,d}(a,b):=-\sigma_{-d}(E_{{\bf p}_1}(a,b))E_{{\bf p}_2}(a,b)
\]
is a generalized Jacobi sum for ${\bf p}_1{\bf p}_2$. 
In particular, since $l\nmid a$, $l\nmid b$, and $l\nmid a+b$, condition (5) in
Definition~\ref{def:T9and10} gives
\[
E_{{\bf p}_1{\bf p}_2,d}(a,b)E_{{\bf p}_1{\bf p}_2,d}(-a,-b) =E_{{\bf p}_1{\bf p}_2,d}(a,b)\overline{E_{{\bf p}_1{\bf p}_2,d}(a,b)}= {\bf p}_1{\bf p}_2=p_1p_2.
\]
Moreover, writing
\[
E_{{\bf p}_1}(a,b)=\sum_{i=0}^{l-1} a_i\zeta_l^i,
\qquad
E_{{\bf p}_2}(a,b)=\sum_{j=0}^{l-1} b_j\zeta_l^j,
\]
we obtain
\[
\sigma_{-d}(E_{{\bf p}_1}(a,b))
=\sigma_{-d}\left(\sum_{i=0}^{l-1}a_i\zeta_l^i\right)
=\sum_{i=0}^{l-1}a_i\zeta_l^{-di}
=\sum_{i=0}^{l-1}a_{-d^{-1}i}\zeta_l^i.
\]
Thus we have
\begin{align*}
E_{{\bf p}_1{\bf p}_2,d}(a,b)
&= -\sigma_{-d}(E_{{\bf p}_1}(a,b))E_{{\bf p}_2}(a,b)\\
&=-\Bigl(\sum_{i=0}^{l-1}a_{-d^{-1}i}\zeta_l^i\Bigr) \Bigl(\sum_{j=0}^{l-1}b_j\zeta_l^j\Bigr)\\
&= -\sum_{k=0}^{l-1}\left(\sum_{i+j\equiv k \,({\rm mod}\,l)}a_{-d^{-1}i}b_j\right)\zeta_l^k \\
&=-\sum_{k=0}^{l-1}\left(\sum_{i=0}^{l-1}a_{-d^{-1}i}b_{k-i}\right)\zeta_l^k\\
&=\sum_{k=0}^{l-1}c_k\zeta_l^k.
\end{align*}

Hence the first assertion follows from Lemma~\ref{lem:BEW}. 
For the reduced form, assume $a_0=b_0=0$. 
Using $1=-(\zeta_l+\cdots+\zeta_l^{l-1})$,
we can eliminate the constant term and rewrite
\begin{align*}
E_{{\bf p}_1{\bf p}_2,d}(a,b)
&= -\sigma_{-d}(E_{{\bf p}_1}(a,b))E_{{\bf p}_2}(a,b)\\
&=-\Bigl(\sum_{i=1}^{l-1}a_{-d^{-1}i}\zeta_l^i\Bigr) \Bigl(\sum_{j=1}^{l-1}b_j\zeta_l^j\Bigr)\\
&= -\sum_{k=0}^{l-1}\left(\sum_{i+j\equiv k \,({\rm mod}\,l)}a_{-d^{-1}i}b_j\right)\zeta_l^k \\
&= -\sum_{k=1}^{l-1}\left(\sum_{i+j\equiv k \,({\rm mod}\,l)}a_{-d^{-1}i}b_j\right)\zeta_l^k 
- \left( \sum_{i+j=l}a_{-d^{-1}i}b_j\right)\\
&=-\sum_{k=1}^{l-1}\left(\sum_{i=1}^{l-1}a_{-d^{-1}i}b_{k-i} - \sum_{i=1}^{l-1}a_{-d^{-1}i}b_{-i}\right)\zeta_l^k\\
&=\sum_{k=1}^{l-1}c_k\zeta_l^k.
\end{align*}
\end{proof}

\begin{remark}\label{rem:base_change}
In the reduced case $a_0=b_0=c_0=0$, the formulas in Theorem~\ref{thm:BF}
are obtained by formal expansion and coefficient comparison, and hence are
polynomial identities in the coefficients $a_i$ and $b_j$.
Therefore they remain valid after any base change and specialization.
\end{remark}

{\it Proof of Theorem~\ref{mainth2}}. 

We first see that the quadratic form
\begin{align*}
q({\bf x})=f_0({\bf x})=\sum_{i=1}^{l-1}x_i^2-\sum_{i=1}^{l-1}x_ix_{i+1}
=(x_1,\ldots,x_{l-1})\,A_{l-1}
\left( 
\begin{array}{c}
x_1\\
\vdots\\
x_{l-1}
\end{array}
\right)
\end{align*}
is regular because it follows from a linear recurrence relation 
det\! $A_{n+2}=$ det\! $A_{n+1}-\frac{1}{4}$ det\! $A_n$ 
with det\! $A_2=\frac{3}{4}$ and det\! $A_1=1$ 
that det\! $A_{l-1}=l/2^{l-1}\neq 0$ 
where $A_{l-1}$ is the corresponding $(l-1)\times (l-1)$ symmetric matrix: 

\renewcommand{\arraystretch}{1.5}
\begin{align*}
A_{l-1}=\left( 
\begin{array}{ccccc}
1 & -\tfrac{1}{2} & 0 & \cdots & 0\\
-\tfrac{1}{2} & 1 & -\frac{1}{2} & \ddots & \vdots\\
0 & -\tfrac{1}{2} & 1 & \ddots & 0\\
\vdots & \ddots & \ddots & \ddots & -\tfrac{1}{2}\\
0 & \cdots & 0 & -\tfrac{1}{2} & 1
\end{array}
\right).
\end{align*}
\renewcommand{\arraystretch}{1.0}

By Remark~\ref{rem:base_change}, in the reduced case $a_0=b_0=c_0=0$
the formulas in Theorem~\ref{thm:BF} remain valid after base change to
$K(x_1,\ldots,x_{l-1},y_1,\ldots,y_{l-1})$ and the specialization
$a_i=x_i$, $b_j=y_j$.
Thus the map $\varphi_d$ defined by
\[
z_k=-\left(\sum_{i=1}^{l-1} x_{-d^{-1}i}y_{k-i}-\sum_{i=1}^{l-1} x_{-d^{-1}i}y_{-i}\right)
\]
satisfies $\varphi_d(V\times V)\subset V$ and $q(\varphi_d(x,y))=q(x)q(y)$.
Hence $q$ is multiplicative on $V$.
\qed

\begin{remark}\label{rem:reg}
We can consider an analog of Theorem~\ref{mainth2} 
when $a_0=x_0\neq0$ and $b_0=y_0\neq 0$ 
as in the first half of Theorem~\ref{thm:BF}. 
However, we see that the corresponding quadratic form
\begin{align*}
q({\bf x})=f_0(x_0,\ldots,x_{l-1})
=\sum_{i=0}^{l-1}x_i^2-\sum_{i=0}^{l-1}x_ix_{i+1}
=(x_0,\ldots,x_{l-1})\,B_l
\left( 
\begin{array}{c}
x_0\\
\vdots\\
x_{l-1}
\end{array}
\right)
\end{align*}
is {\it not} regular, since each row and column sums to zero, and hence det\! $B_l=0$ 
where $B_l$ is the corresponding $l\times l$ symmetric matrix:
%
\renewcommand{\arraystretch}{1.5}
\begin{align*}
B_l=\left(
\begin{array}{c|ccccc}
1 & -\tfrac{1}{2} & 0 & \cdots & 0 & -\tfrac{1}{2}\\\hline
-\tfrac{1}{2}& 1 & -\tfrac{1}{2} & 0 & \cdots & 0\\
0 & -\tfrac{1}{2} & 1 & -\frac{1}{2} & \ddots & \vdots\\
\vdots & 0 & -\tfrac{1}{2} & 1 & \ddots & 0\\
0 & \vdots & \ddots & \ddots & \ddots & -\tfrac{1}{2}\\
-\tfrac{1}{2} & 0 & \cdots & 0 & -\tfrac{1}{2} & 1
\end{array}
\right).
\end{align*}
\renewcommand{\arraystretch}{1.0}
%

Thus, to obtain a regular quadratic form, one has to eliminate one variable, for instance by imposing $x_0=0$. 
This is precisely the reduced situation in the second half of Theorem~\ref{thm:BF}, where one also normalizes the coefficients by $a_0=b_0=0$.
\end{remark}

\begin{example}[Examples of Theorem~\ref{mainth2} for $l=3,5,7$]\label{ex:mainth2}
Let $K$ be a field with char $K\neq 2$, $l$ and 
$\Gal(K(\zeta_l)/K)\simeq (\Z/l\Z)^{\times}$. 
Each subscript of the $x_i$'s and the $y_j$'s 
should be taken modulo $l$
and we take $x_0=y_0=0$ 
(for the regularity of $q({\bf x})$, see Remark \ref{rem:reg}).\\ 

(1) $l=3$. It follows from Theorem~\ref{mainth2} that 
the quadratic form $q({\bf x})=q(x_1,x_{2})$ defined by
\[
q({\bf x})=f_0({\bf x})
=\sum_{i=1}^{2}x_i^2-\sum_{i=1}^{2}x_ix_{i+1}
=(x_1^2+x_2^2) - (x_1x_2+x_2x_0)
=x_1^2+x_2^2 - x_1x_2
\]
(with the convention $x_0=0$) is multiplicative on $V=K^2$. 

Indeed, for $d \in (\Z/3\Z)^\times$, 
we obtain an explicit bilinear map $\varphi_d :K^2 \times K^2 \rightarrow K^2$ such that 
\begin{align*}
q({\bf x})q({\bf y})=q(\varphi_d ({\bf x}, {\bf y}))
\end{align*}
for ${\bf x},{\bf y} \in K^2$.
It is given explicitly by
\[
\varphi_d ({\bf x}, {\bf y})
={\bf z}_d
= (z_{1,d}, z_{2,d})
\]
where $z_{k,d}=-\left(\sum_{i=1}^{2}x_{-d^{-1}i}y_{k-i} - \sum_{i=1}^{2}x_{-d^{-1}i}y_{-i}\right)$.

For example, we have for $d=1 \in (\Z/3\Z)^\times$,
\begin{align*}
z_{1,1}
=- x_1y_2 + (x_2y_2 + x_1y_1),\\
z_{2,1}
=-x_2y_1+(x_2y_2 + x_1y_1).
\end{align*}

(2) $l=5$. 
We take an algebraic $K$-variety $V \subsetneq K^4$ defined by
$f_1({\bf x})=0$
where
\begin{align*}
f_1({\bf x})
&=\sum_{i=1}^{4}x_ix_{i+1}-\sum_{i=1}^{4}x_ix_{i+2}\\
&=(x_1x_2 + x_2x_3 + x_3x_4) - (x_1x_3 + x_2x_4 + x_4x_1).
\end{align*}
By Theorem~\ref{mainth2}, the quadratic form 
\begin{align*}
q({\bf x})=f_0({\bf x})
=\sum_{i=1}^{4}x_i^2-\sum_{i=1}^{4}x_ix_{i+1}
=(x_1^2+x_2^2+x_3^2+x_4^2) - (x_1x_2+x_2x_3 + x_3x_4)
\end{align*}
is multiplicative on $V=\{{\bf x}\in K^4\mid f_1({\bf x})=0\}$. 
Indeed, for $d \in (\Z/5\Z)^\times$, we obtain an explicit bilinear map  $\varphi_d :K^4 \times K^4 \rightarrow K^4$ such that 
\begin{align*}
q({\bf x})q({\bf y})=q(\varphi_d ({\bf x}, {\bf y}))
\end{align*}
for ${\bf x},{\bf y} \in V$.
It is given explicitly by
\[
\varphi_d ({\bf x}, {\bf y})
={\bf z}
= (z_{1,d}, z_{2,d}, z_{3,d}, z_{4,d})
\]
where $z_{k,d}=-\left(\sum_{i=1}^{4}x_{-d^{-1}i}y_{k-i} - \sum_{i=1}^{4}x_{-d^{-1}i}y_{-i}\right)$.

For example, we have for $d=1 \in (\Z/5\Z)^\times$,
\begin{align*}
z_{1,1}
&=-(x_3y_4 + x_2y_3 + x_1y_2 )
+(x_4y_4 + x_3y_3 + x_2y_2 + x_1y_1 ),\\
z_{2,1}
&=-(x_{4}y_1 + x_2y_4 + x_1y_3 )
+(x_{4}y_4 + x_3y_3 + x_2y_2 + x_1y_1 ),\\
z_{3,1}
&=-(x_4y_2 + x_3y_1 + x_1y_4 )
+(x_4y_4 + x_3y_3 + x_2y_2 + x_1y_1 ),\\
z_{4,1}
&=-(x_{4}y_3 + x_3y_2 + x_2y_1)
+(x_{4}y_4 + x_3y_3 + x_2y_2 + x_1y_1 ).
\end{align*}

Moreover, for $d=\pm 1\in (\Z/5\Z)^\times$, 
one can check the following identities, which in particular imply that $q({\bf x})$ is multiplicative on $V$, i.e. 
$q({\bf x})q({\bf y})=q(\varphi_d({\bf x}, {\bf y}))$ for 
${\bf x},{\bf y} \in V=\{{\bf x}\in K^4\mid f_1({\bf x})=0\}$:
\begin{align*}
q(\varphi_d({\bf x},{\bf y}))&=
q({\bf x})q({\bf y})+f_1({\bf x})f_1({\bf y}),\\
f_1(\varphi_d({\bf x},{\bf y}))&=
q({\bf x})f_1({\bf y})+f_1({\bf x})q({\bf y})+f_1({\bf x})f_1({\bf y}).
\end{align*}
Note that $q({\bf x})$ and $f_1({\bf x})$ are invariant 
under the change $d=1\mapsto -1$. 
For $d=\pm 2\in (\Z/5\Z)^\times$, we obtain that 
\begin{align*}
q(\varphi_d({\bf x},{\bf y}))&=
q({\bf x})q({\bf y})+f_1({\bf x})q({\bf y})-f_1({\bf x})f_1({\bf y}),\\
f_1(\varphi_d({\bf x},{\bf y}))&=
q({\bf x})f_1({\bf y})-f_1({\bf x})q({\bf y})
\end{align*}
because $q({\bf x})\mapsto q({\bf x})+f_1({\bf x})$, 
$f_1({\bf x})\mapsto -f_1({\bf x})$ after changing $d=1\mapsto \pm 2\in (\Z/5\Z)^\times$.\\

(3) $l=7$. 
We consider an algebraic $K$-variety $V \subsetneq K^6$ defined by
$f_1({\bf x})=0$, $f_2({\bf x})=0$
where
\begin{align*}
f_1({\bf x})
&=\sum_{i=1}^{6}x_ix_{i+1}-\sum_{i=1}^{6}x_ix_{i+2}\\
&=(x_1x_2+x_2x_3 + x_3x_4 + x_4x_5 + x_5x_6)\\
& - (x_1x_3+x_2x_4 + x_3x_5 + x_4x_6 + x_6x_1), \\
f_2({\bf x})
&=\sum_{i=1}^{6}x_ix_{i+2}-\sum_{i=1}^{6}x_ix_{i+3}\\
&= (x_1x_3+x_2x_4 + x_3x_5 + x_4x_6 + x_6x_1)\\
& - (x_1x_4+x_2x_5 + x_3x_6 + x_5x_1 + x_6x_2).
\end{align*}
By Theorem~\ref{mainth2}, the quadratic form 
\begin{align*}
q({\bf x})=f_0({\bf x})
=\sum_{i=1}^{6}x_i^2-\sum_{i=1}^{6}x_ix_{i+1}
&=(x_1^2+x_2^2+x_3^2+x_4^2+x_5^2+x_6^2)\\ 
&- (x_1x_2+x_2x_3 + x_3x_4 + x_4x_5 + x_5x_6)
\end{align*}
is multiplicative on $V=\{{\bf x}\in K^6\mid f_1({\bf x})=f_2({\bf x})=0\}$. 
Indeed, for $d \in (\Z/7\Z)^\times$, we obtain an explicit bilinear map  $\varphi_d :K^6 \times K^6 \rightarrow K^6$ such that 
\begin{align*}
q({\bf x})q({\bf y})=q(\varphi_d ({\bf x}, {\bf y}))
\end{align*}
for ${\bf x},{\bf y} \in V$.
It is given explicitly by
\[
\varphi_d ({\bf x}, {\bf y})
= (z_{1,d}, z_{2,d}, z_{3,d}, z_{4,d},z_{5,d},z_{6,d})
\]
where $z_{k,d}=-\left(\sum_{i=1}^{6}x_{-d^{-1}i}y_{k-i} - \sum_{i=1}^{6}x_{-d^{-1}i}y_{-i}\right)$.
For example, we have for $d=1 \in (\Z/7\Z)^\times$,
\begin{align*}
z_{1,1}
&=-( x_5y_6 + x_4y_5 + x_3y_4 + x_2y_3 + x_1y_2 )
+(x_6y_6 + x_5y_5 + x_4y_4 + x_3y_3 + x_2y_2 + x_1y_1 ),\\
z_{2,1}
&=-(x_6y_1 + x_4y_6 + x_3y_5 + x_2y_4 + x_1y_3 )
+(x_6y_6 + x_5y_5 + x_4y_4 + x_3y_3 + x_2y_2 + x_1y_1 ),\\
z_{3,1}
&=-(x_6y_2 + x_5y_1 + x_3y_6 + x_2y_5 + x_1y_4 )
+(x_6y_6 + x_5y_5 + x_4y_4 + x_3y_3 + x_2y_2 + x_1y_1 ),\\
z_{4,1}
&=-(x_6y_3 + x_5y_2 + x_4y_1 + x_2y_6 + x_1y_5 )
+(x_6y_6 + x_5y_5 + x_4y_4 + x_3y_3 + x_2y_2 + x_1y_1 ),\\
z_{5,1}
&=-( x_6y_4 + x_5y_3 + x_4y_2 + x_3y_1 + x_1y_6 )
+(x_6y_6 + x_5y_5 + x_4y_4 + x_3y_3 + x_2y_2 + x_1y_1 ),\\
z_{6,1}
&=-(x_6y_5 + x_5y_4 + x_4y_3 + x_3y_2 + x_2y_1)
+(x_6y_6 + x_5y_5 + x_4y_4 + x_3y_3 + x_2y_2 + x_1y_1 ). 
\end{align*}

Moreover, for $d=\pm 1\in (\Z/7\Z)^\times$, 
one can check the following identities, which in particular imply that $q({\bf x})$ is multiplicative on $V$, i.e. 
$q({\bf x})q({\bf y})=q(\varphi_d({\bf x}, {\bf y}))$ for 
${\bf x},{\bf y} \in V=\{{\bf x}\in K^6\mid f_1({\bf x})=f_2({\bf x})=0\}$:
\begin{align*}
q(\varphi_d({\bf x},{\bf y}))&=
q({\bf x})q({\bf y})+f_1({\bf x})f_1({\bf y})+f_2({\bf x})f_2({\bf y}),\\
f_1(\varphi_d({\bf x},{\bf y}))&=
q({\bf x})f_1({\bf y})+f_1({\bf x})q({\bf y})+f_1({\bf x})f_1({\bf y})+f_1({\bf x})f_2({\bf y})+f_2({\bf x})f_1({\bf y})+f_2({\bf x})f_2({\bf y}),\\
f_2(\varphi_d({\bf x},{\bf y}))&=
q({\bf x})f_2({\bf y})+f_2({\bf x})q({\bf y})+f_1({\bf x})f_1({\bf y})+f_1({\bf x})f_2({\bf y})+f_2({\bf x})f_1({\bf y}).
\end{align*}
Note that $q({\bf x})$, $f_1({\bf x})$, $f_2({\bf x})$ 
are invariant under the change $d=1\mapsto -1$. 
For $d=\pm 3\in (\Z/7\Z)^\times$, we obtain that
\begin{align*}
q(\varphi_d({\bf x},{\bf y}))&=
q({\bf x})q({\bf y})+f_1({\bf x})q({\bf y})-f_1({\bf x})f_2({\bf y})+f_2({\bf x})f_1({\bf y})-f_2({\bf x})f_2({\bf y}),\\
f_1(\varphi_d({\bf x},{\bf y}))&=
q({\bf x})f_1({\bf y})+f_2({\bf x})q({\bf y})-f_1({\bf x})f_2({\bf y}),\\
f_2(\varphi_d({\bf x},{\bf y}))&=
q({\bf x})f_2({\bf y})-f_1({\bf x})q({\bf y})-f_2({\bf x})q({\bf y})-f_1({\bf x})f_1({\bf y})+f_1({\bf x})f_2({\bf y})+f_2({\bf x})f_2({\bf y})
\end{align*}
because we get 
$\overline{\sigma}: q({\bf x})\mapsto q({\bf x})+f_1({\bf x})$, 
$f_1({\bf x})\mapsto f_2({\bf x})$, $f_2({\bf x})\mapsto -f_1({\bf x})-f_2({\bf x})$ by changing $d=1\mapsto \pm 3\in (\Z/7\Z)^\times$. 
Applying this twice, we get for $d=\pm 2$ $(=\pm 3^2)\in (\Z/7\Z)^\times$,
\begin{align*}
q(\varphi_d({\bf x},{\bf y}))&=
q({\bf x})q({\bf y})+f_1({\bf x})q({\bf y})+f_2({\bf x})q({\bf y})-f_1({\bf x})f_1({\bf y})+f_1({\bf x})f_2({\bf y})-f_2({\bf x})f_1({\bf y}),\\
f_1(\varphi_d({\bf x},{\bf y}))&=
q({\bf x})f_1({\bf y})-f_1({\bf x})q({\bf y})-f_2({\bf x})q({\bf y})+f_1({\bf x})f_1({\bf y})-f_2({\bf x})f_2({\bf y}),\\
f_2(\varphi_d({\bf x},{\bf y}))&=
q({\bf x})f_2({\bf y})+f_1({\bf x})q({\bf y})-f_2({\bf x})f_1({\bf y}) 
\end{align*}
where $\overline{\sigma}^{\,2}: 
q({\bf x})\mapsto q({\bf x})+f_1({\bf x})+f_2({\bf x})$, 
$f_1({\bf x})\mapsto -f_1({\bf x})-f_2({\bf x})$, 
$f_2({\bf x})\mapsto f_1({\bf x})$ by changing $d=1\mapsto \pm 2$ 
$(=\pm 3^2)\in (\Z/7\Z)^\times$. 
Note that $(\Z/7\Z)^\times/\{\pm 1\}$ is a cyclic group of order $3$ 
and $\overline{\sigma}^{\,3}={\rm id}$. 
\end{example}

%------------------------------------S4-3
\subsection{Explicit lifts of solutions for Diophantine equations}\label{S4-3}

We apply Theorem~\ref{thm:BF} to systems of Diophantine equations associated with Jacobi sums and cyclotomic numbers. 
This recovers explicit lifting formulas for the prime-power systems studied in Hoshi and Kanai~\cite{HK22}, and also explains the corresponding multiplicative quadratic forms on algebraic $K$-varieties considered in Hoshi~\cite{Hos03}.

With the notation $W_p$ introduced before Corollary~\ref{cor:lifting}, 
classical Jacobi sums in the prime-power case furnish integral points on the fibers $W_{p^r}$. 
The composition law then yields integral points on $W_{p^r q^s}$ from those on $W_{p^r}$ and $W_{q^s}$.

(1) $l=3$ (cf. Gauss \cite[Section 358]{Gau1801}, Katre and Rajwade \cite[Proposition~1]{KR85a}, \cite[Section 3.1, Section 10.10]{BEW98}). 
Let $p^r, q^s$ be prime powers with $p^r\equiv 1\ ({\rm mod}\ 3)$, $q^s\equiv 1\ ({\rm mod}\ 3)$. 
Note that we can consider both of the cases $p\neq q$ 
and $p=q$ (see the end of this section for the case where $p=q$).  
Let
$J^\ast_{p^r}(1,1)=\sum_{i=1}^{2}a_{i}\zeta_3^{i}$ $($resp. $J^\ast_{q^s}(1,1)=\sum_{j=1}^{2}b_{j}\zeta_3^{j}$$)$ 
be Jacobi sums for $\F_{p^r}$ $($resp. $\F_{q^s}$$)$. 
We see that $J^\ast_{p^r}(1,1)$ and $J^\ast_{q^s}(1,1)$ are given by 
\begin{align*}
&J^\ast_{p^r}(1,1)
=\frac{1}{2}(x_1+3x_2\sqrt{-3})
=\frac{-x_1+3x_2}{2}\zeta_3 + \frac{-x_1-3x_2}{2}\zeta_3^2,\\
&J^\ast_{q^s}(1,1)
=\frac{1}{2}(y_1+3y_2\sqrt{-3})
=\frac{-y_1+3y_2}{2}\zeta_3 + \frac{-y_1-3y_2}{2}\zeta_3^2
\end{align*}
where $x_1,x_2,y_1,y_2 \in\Z$ are given as integer solutions of the systems of Diophantine equations
\begin{align}\label{eq:DS3p}
\begin{cases}
4p^r=x_1^2+27x_2^2,\\
x_1\equiv 1\ ({\rm mod}\ 3),\ p\notmid x_1,
\end{cases}
\end{align}
\begin{align*} %\label{eq:DS3q}
\begin{cases}
4q^s=y_1^2+27y_2^2,\\
y_1\equiv 1\ ({\rm mod}\ 3),\ q\notmid y_1. 
\end{cases}
\end{align*}
Each system has two solutions $(x_1,\pm x_2)$ and $(y_1,\pm y_2)$ 
which depend on a choice of generator of 
$\F_{p^r}$ and $\F_{q^s}$ respectively.

By Theorem~\ref{mainth1} and Theorem~\ref{thm:BF},
for $d\in(\Z/3\Z)^{\times}$, 
generalized Jacobi sums
\[
J^\ast_{p^rq^s,d}(1,1)=-\sigma_{-d}(J^\ast_{p^r}(1,1))J^\ast_{q^s}(1,1)=\sum_{k=1}^{2}c_k\zeta_3^k
\]
for $p^rq^s$ are given by
$c_k=-(\sum_{i=1}^{2}a_{-d^{-1}i}b_{k-i} - \sum_{i=1}^{2}a_{-d^{-1}i}b_{-i})$ 
where $J^\ast_{p^r}(1,1)=\sum_{i=1}^2 a_i\zeta_3^i$, 
$J^\ast_{q^s}(1,1)=\sum_{j=1}^2 b_j\zeta_3^j$. 

By Theorem~\ref{thm:BF} again, we have
\begin{align}\label{eq:DS3c}
p^rq^s=c_1^2+c_2^2-c_1c_2.
\end{align}
We define
\begin{align*}
z_{1,d}&:=c_1+c_2=-\frac{1}{2}\sigma^i(x_1y_1+27x_2y_2),\\
z_{2,d}&:=\frac{1}{3}\left(c_1-c_2\right)=-\frac{1}{2}\sigma^i(x_1y_2-y_1 x_2)
\end{align*}
where $\sigma(x_1,x_2,y_1,y_2)=(x_1,-x_2,y_1,y_2)$ and $i=0,1$ if $d=1,2$. 
By the
equation (\ref{eq:DS3c}), we can verify $z_{1,d}, z_{2,d}$ satisfy
\begin{align*}
z_{1,d}^2+27z_{2,d}^2
=4(c_1^2+c_2^2-c_1c_2)
=4p^rq^s.
\end{align*} 
Thus $(z_{1,d},z_{2,d})$ yields solutions of the system of Diophantine equations for $p^rq^s$:
\begin{align*}
\begin{cases}
4p^rq^s=z_{1,d}^2+27z_{2,d}^2,\\
z_{1,d}\equiv 1\ ({\rm mod}\ 3).
\end{cases}
\end{align*}

From the argument above, if we take
${\bf z}^\prime=(z_{1,d}^{\prime}, z_{2,d}^{\prime}) = (2z_{1,d}, 2z_{2,d})$, 
then we get a quadratic form $q({\bf x})=x_1^2+27x_2^2$ which 
is multiplicative on $V=K^2$ with 
$\varphi_d ({\bf x}, {\bf y})={\bf z}^\prime=(z_{1,d}^{\prime}, z_{2,d}^{\prime})$. 
Indeed, we have
\begin{align*}
q(\varphi_d ({\bf x}, {\bf y}))=q({\bf z}^\prime)
=(z_{1,d}^{\prime})^2+27(z_{2,d}^{\prime})^2
=4(z_{1,d}^2+27z_{2,d}^2)
=4(4p^rq^s)
=(4p^r)(4q^s)
=q({\bf x})q({\bf y}).
\end{align*}
Thus we recover the well-known formula:
\begin{align*}
(x_1^2+27x_2^2)(y_1^2+27y_2^2)=(x_1y_1\pm 27 x_2y_2)^2+27(x_1y_2\mp y_1x_2)^2. 
\end{align*}

In the case
$p=q$ (i.e. $x_i=y_i$), 
a lift of Jacobi sums $J^\ast_{p^r}(1,1)$ to $J^\ast_{p^{nr}}(1,1)$ 
can be obtained by Davenport and Hasse's lifting theorem 
(see Theorem~\ref{thm:corMT1} and Section \ref{S2}). 
Hoshi and Kanai \cite[Section 7, (1) $e=3$]{HK22} 
gave such an explicit lift of solutions of the system of Diophantine equations (\ref{eq:DS3p}) corresponding to Davenport and Hasse's lifting theorem, cf. Theorem~\ref{thm:corMT1} with 
$d=2=-1 \in(\Z/3\Z)^{\times}$. 
Indeed, in such a case, 
\begin{align*}
z_{1,-1}&=-\frac{1}{2}(x_1^2-27x_2^2),\\
z_{2,-1}&=-x_1x_2
\end{align*}
are given as $c^{(2)}$, $d^{(2)}$ in Hoshi and Kanai \cite[page 19, line 8]{HK22}.\\

(2) $l=5$ (cf. Lehmer \cite[Equation (10)]{Leh51}, Berndt and Evans \cite[Section 5]{BE81},  Katre and Rajwade \cite{KR85b}, \cite[Section 3.7]{BEW98}, Hoshi \cite[Section 5]{Hos03}, \cite[Section 3]{Hos06}).
Let $p^r, q^s$ be prime powers with $p^r\equiv 1\ ({\rm mod}\ 5)$, $q^s\equiv 1\ ({\rm mod}\ 5)$. 
Note that we can consider both of the cases $p\neq q$ 
and $p=q$ (see the end of this section for the case where $p=q$). 
Let 
$J^\ast_{p^r}(1,1)=\sum_{i=1}^{4}a_{i}\zeta_5^{i}$ $($resp. $J^\ast_{q^s}(1,1)=\sum_{j=1}^{4}b_{j}\zeta_5^{j}$$)$ 
be 
Jacobi sums for $\F_{p^r}$ $($resp. $\F_{q^s}$$)$. 
We see that $J^\ast_{p^r}(1,1)$ and $J^\ast_{q^s}(1,1)$ are given by 
\begin{align*}
J^\ast_{p^r}(1,1)
=\frac{1}{4}\left(T\zeta_5+\sigma^3(T)\zeta_5^2+\sigma(T)\zeta_5^3+\sigma^2(T)\zeta_5^4\right),
\end{align*}
\begin{align*}
J^\ast_{q^s}(1,1)
=\frac{1}{4}\left(U\zeta_5+\tau^3(U)\zeta_5^2+\tau(U)\zeta_5^3+\tau^2(U)\zeta_5^4\right)
\end{align*}
where $T=-x_1+5x_2+4x_3+2x_4$, $U=-y_1+5y_2+4y_3+2y_4$, $\sigma(x_1,x_2,x_3,x_4)=(x_1,-x_2,-x_4,x_3)$, $\tau(y_1,y_2,y_3,y_4)=(y_1,-y_2,-y_4,y_3)$ and
$x_1,x_2,x_3,x_4, y_1,y_2,y_3,y_4\in\Z$ are obtained as 
integer solutions of the systems of Diophantine equations 
\begin{align}\label{eq:DS5p}
\begin{cases}
16p^r=x_1^2+125x_2^2+50x_3^2+50x_4^2,\\
x_1x_2=x_3^2-4x_3x_4-x_4^2, \\
x_1 \equiv 1\ ({\rm mod}\ 5),\ p\notmid x_1^2-125x_2^2,
\end{cases} %\label{eq:Eq5}
\end{align}
\begin{align*} %\label{eq:DS5q}
\begin{cases}
16q^s=y_1^2+125y_2^2+50y_3^2+50y_4^2,\\
y_1y_2=y_3^2-4y_3y_4-y_4^2, \\
y_1 \equiv 1\ ({\rm mod}\ 5),\ q\notmid y_1^2-125y_2^2.
\end{cases} %\label{eq:Eq5}
\end{align*}
Each system has four solutions $\sigma^i(x_1,x_2,x_3,x_4)$ and $\tau^i(y_1,y_2,y_3,y_4)$ $(i=0,1,2,3)$ 
which depend on a choice of generator of $\F_{p^r}$ and $\F_{q^s}$ 
respectively.

By Theorem~\ref{mainth1} and Theorem~\ref{thm:BF},
for $d\in(\Z/5\Z)^{\times}$, generalized Jacobi sums
\[
J^\ast_{p^rq^s,d}(1,1)=
-\sigma_{-d}(J^\ast_{p^r}(1,1))J^\ast_{q^s}(1,1)=\sum_{k=1}^{4}c_k\zeta_5^k
\]
for $p^rq^s$ are given by
$c_k=-(\sum_{i=1}^{l-1}a_{-d^{-1}i}b_{k-i} - \sum_{i=1}^{l-1}a_{-d^{-1}i}b_{-i})$, 
$a_1=T, a_2=\sigma^3(T), a_3=\sigma(T), a_4=\sigma^2(T), b_1=U, b_2=\tau^3(U), b_3=\tau(U), b_4=\tau^2(U)$.
The $c_k$'s satisfy
\begin{align}\label{eq:DS5c}
\begin{dcases}
p^rq^s
=c_1^2+c_2^2+c_3^2+c_4^2-c_1c_2-c_2c_3-c_3c_4, \\
c_1c_2+c_2c_3+c_3c_4 - (c_1c_3+c_2c_4+c_4c_1)=0.
\end{dcases}
\end{align}

We define 
\begin{align*} %\label{eq:BF5}
\begin{split}
z_{1,d}&:=c_1+c_2+c_3+c_4=-\frac{1}{4}\sigma^i\left(x_1y_1 + 125x_2y_2 + 50x_3y_3 + 50x_4y_4\right),\\
z_{2,d}&:=-\frac{1}{5}\left(c_1-c_2-c_3+c_4\right)=-\frac{1}{4}\sigma^i\left(x_1y_2 + x_2y_1 - 2x_3y_3 + 4x_3y_4 + 4x_4y_3  + 2x_4y_4\right),\\
z_{3,d}&:=-\frac{1}{5}\left(2c_1-c_2+c_3-2c_4\right)=-\frac{1}{4}\sigma^i\left( x_1y_3 - 5x_2y_3 + 10x_2y_4 - x_3y_1 + 5x_3y_2 - 10x_4y_2 \right),\\
z_{4,d}&:=-\frac{1}{5}\left(c_1+2c_2-2c_3-c_4\right)=-\frac{1}{4}\sigma^i\left(x_1y_4 + 10x_2y_3 + 5x_2y_4 - 10x_3y_2 - x_4y_1 - 5x_4y_2  \right)
\end{split}
\end{align*}
where $\sigma(x_1,x_2,x_3,x_4,y_1,y_2,y_3,y_4)=(x_1,-x_2,-x_4,x_3,y_1,y_2,y_3,y_4)$ and $i=0$ 
(resp. $1,3,2$) if $d=1$ (resp. $2,3,4$). 
By the equation (\ref{eq:DS5c}), 
we can verify that $z_{1,d},z_{2,d},z_{3,d}, z_{4,d}$ satisfy
\begin{align*}
z_{1,d}^2+125z_{2,d}^2+50z_{3,d}^2+50z_{4,d}^2
&=16(c_1^2 + c_2^2 + c_3^2 + c_4^2) - 8( (c_1c_2  +  c_2c_3 + c_3c_4) +( c_1c_3 + c_2c_4 + c_4c_1))\\
&=16(c_1^2+c_2^2+c_3^2+c_4^2-(c_1c_2+c_2c_3+c_3c_4))\\
&=16p^rq^s
\end{align*}
and
\begin{align*}
z_{1,d}z_{2,d}-z_{3,d}^2+4z_{3,d}z_{4,d}+z_{4,d}^2=\frac{5}{4}\left(c_1c_2 + c_2c_3 + c_3c_4 - (c_1c_3 + c_2c_4 + c_4c_1)\right)=0.
\end{align*}
Thus $(z_{1,d},z_{2,d},z_{3,d},z_{4,d})$
yields solutions of the system of Diophantine equations 
for $p^rq^s$:
\begin{align*}
\begin{cases}
16p^rq^s=z_{1,d}^2+125z_{2,d}^2+50z_{3,d}^2+50z_{4,d}^2,\\
z_{1,d}z_{2,d}=z_{3,d}^2-4z_{3,d}z_{4,d}-z_{4,d}^2, \\
z_{1,d} \equiv 1\ ({\rm mod}\ 5). 
\end{cases}
\end{align*}

For 
${\bf z}^{\prime} = (z_{1,d}^{\prime},z_{2,d}^{\prime},z_{3,d}^{\prime},z_{4,d}^{\prime})=(4z_{1,d},4z_{2,d},4z_{3,d},4z_{4,d})$, 
the quadratic form $q({\bf x})=x_1^2+125x_2^2+50x_3^2+50x_4^2$ is 
multiplicative on $V$ defined by $x_1x_2=x_3^2-4x_3x_4-x_4^2$
with 
$\varphi_d ({\bf x}, {\bf y})=(z_{1,d}^{\prime}, z_{2,d}^{\prime}, z_{3,d}^{\prime}, z_{4,d}^{\prime})$. 
Indeed, we have
\begin{align*}
q(\varphi_d ({\bf x}, {\bf y}))=q({\bf z}^{\prime})
&={(z_{1,d}^{\prime})}^2+125{(z_{2,d}^{\prime})}^2+50{(z_{3,d}^{\prime})}^2+50{(z_{4,d}^{\prime})}^2\\
&=16(z_{1,d}^2+125z_{2,d}^2+50z_{3,d}^2+50z_{4,d}^2)\\
&=16(16p^rq^s)
=(16p^r)(16q^s)\\
&=(x_1^2+125x_2^2+50x_3^2+50x_4^2)(y_1^2+125y_2^2+50y_3^2+50y_4^2)\\
&=q({\bf x})q({\bf y})
\end{align*}
and
\[
z_{1,d}^{\prime}z_{2,d}^{\prime}={(z_{3,d}^{\prime})}^2-4z_{3,d}^{\prime}z_{4,d}^{\prime}-{(z_{4,d}^{\prime})}^2.
\]
If we take $d=1$ and define $Z_{i}=-z^{\prime}_{i,1}$, 
then we recover the bilinear form $\varphi_1 ({\bf x}, {\bf y})=(Z_1,Z_2,Z_3,Z_4)$  
given in Hoshi \cite[Section 3, Equation (6)]{Hos03}:
\begin{align*}
Z_1&=x_1y_1 + 125x_2y_2 + 50x_3y_3 + 50x_4y_4,\\
Z_2&=x_1y_2 + x_2y_1 - 2x_3y_3 + 4x_3y_4 + 4x_4y_3  + 2x_4y_4,\\
Z_3&=x_1y_3 - 5x_2y_3 + 10x_2y_4 - x_3y_1 + 5x_3y_2 - 10x_4y_2,\\
Z_4&=x_1y_4 + 10x_2y_3 + 5x_2y_4 - 10x_3y_2 - x_4y_1 - 5x_4y_2.
\end{align*}

In the case
$p=q$ (i.e. $x_i=y_i$), 
a lift of Jacobi sums $J^\ast_{p^r}(1,1)$
to $J^\ast_{p^{nr}}(1,1)$
can be obtained by Davenport and Hasse's lifting theorem 
(see Theorem~\ref{thm:corMT1} and Section \ref{S2}).  
Hoshi and Kanai \cite[Section 7, (2) $e=5$]{HK22} 
gave such an explicit lift of solutions of the system of Diophantine equations (\ref{eq:DS5p}) corresponding to Davenport and Hasse's lifting theorem, cf. Theorem~\ref{thm:corMT1} with 
$d=4=-1 \in(\Z/5\Z)^{\times}$. 
Indeed,
\begin{align*}
z_{1,-1}&=-\frac{1}{4}(x_1^2+125x_2^2-50x_3^2-50x_4^2),\\
z_{2,-1}&=-\frac{1}{2}(x_1x_2+x_3^2-4x_3x_4-x_4^2),\\
z_{3,-1}&=-\frac{1}{2}(x_1x_3+10x_2x_4-5x_3x_2),\\
z_{4,-1}&=-\frac{1}{2}(x_1x_4+10x_2x_3+5x_4x_2)
\end{align*}
are given as
$x^{(2)}$, $w^{(2)}$, $v^{(2)}$, $u^{(2)}$ 
in Hoshi and Kanai \cite[page 21, lines 10--13]{HK22}.\\

(3) $l=7$ (cf. Leonard and Williams \cite{LW75a}, \cite[Section 3.9]{BEW98}).
Let $p^r, q^s$ be prime powers with $p^r\equiv 1\ ({\rm mod}\ 7)$, $q^s\equiv 1\ ({\rm mod}\ 7)$. 
Note that we can consider both of the cases $p\neq q$ 
and $p=q$ (see the end of this section for the case where $p=q$). 
Let
$J^\ast_{p^r}(1,1)=\sum_{k=1}^{6}a_k\zeta_7^k$ $($resp. $J^\ast_{q^s}(1,1)=\sum_{k=1}^{6}b_k\zeta_7^k$$)$ 
be
Jacobi sums for $\F_{p^r}$ $($resp. $\F_{q^s}$$)$. 
We see that 
$J^\ast_{p^r}(1,1)$ and
$J^\ast_{q^s}(1,1)$
are given by 
\begin{align*}
J^\ast_{p^r}(1,1)
&=\frac{1}{12}\left(T\zeta_7+\sigma^4(T)\zeta_7^2+\sigma^5(T)\zeta_7^3
+\sigma^2(T)\zeta_7^4+\sigma(T)\zeta_7^5+\sigma^3(T)\zeta_7^6\right),\\
J^\ast_{q^s}(1,1) &=\frac{1}{12}\left(U\zeta_7+\tau^4(U)\zeta_7^2+\tau^5(U)\zeta_7^3
+\tau^2(U)\zeta_7^4+\tau(U)\zeta_7^5+\tau^3(U)\zeta_7^6\right)
\end{align*}
where $T=-2x_1+6x_2+7x_5+21x_6$, $U=-2y_1+6y_2+7y_5+21y_6$, $\sigma(x_1,x_2,x_3,x_4,x_5,x_6)=(x_1,-x_3,x_4,x_2,(-x_5-3x_6)/2,(x_5-x_6)/2)$, $\tau(y_1,y_2,y_3,y_4,y_5,y_6)=(y_1,-y_3,y_4,y_2,(-y_5-3y_6)/2,(y_5-y_6)/2)$  and
$x_1,\dots,x_6, y_1,\dots,y_6\in\Z$ are obtained as 
the integer solutions of the systems of Diophantine equations 
\begin{align}\label{eq:DS7p}
\begin{cases}
72p^r=2x_1^2+42x_2^2+42x_3^2+42x_4^2+343x_5^2+1029x_6^2, \\
12x_2^2-12x_4^2+147x_5^2-441x_6^2+56x_1x_6+24x_2x_3
-24x_2x_4+48x_3x_4+98x_5x_6=0, \\
12x_3^2-12x_4^2+49x_5^2-147x_6^2+28x_1x_5+28x_1x_6+48x_2x_3+24x_2x_4+24x_3x_4+490x_5x_6=0,\\
x_1\equiv 1\ ({\rm mod}\ 7),\ (x_5,x_6)\neq (0,0),\\
\end{cases}
\end{align}
\begin{align*} %\label{eq:DS7q}
\begin{cases}
72q^s=2y_1^2+42y_2^2+42y_3^2+42y_4^2+343y_5^2+1029y_6^2, \\
12y_2^2-12y_4^2+147y_5^2-441y_6^2+56y_1y_6+24y_2y_3
-24y_2y_4+48y_3y_4+98y_5y_6=0, \\
12y_3^2-12y_4^2+49y_5^2-147y_6^2+28y_1y_5+28y_1y_6+48y_2y_3+24y_2y_4+24y_3y_4+490y_5y_6=0,\\
y_1\equiv 1\ ({\rm mod}\ 7),\ (y_5,y_6)\neq (0,0),\\
\end{cases}
\end{align*}
i.e. 
$a_1=T, a_2=\sigma^4(T), a_3=\sigma^5(T), a_4=\sigma^2(T),a_5=\sigma(T), a_6=\sigma^3(T),
b_1=U, b_2=\tau^4(U), b_3=\tau^5(U), b_4=\tau^2(U),b_5=\tau(U), b_6=\tau^3(U)$. 
Each system has six solutions \\$\sigma^i(x_1,x_2,x_3,x_4,x_5,x_6)$ and $\tau^i(y_1,y_2,y_3,y_4,y_5,y_6)$ $(i=0,\dots,5)$ 
which depend on the choice of generator of $\F_{p^r}$ and $\F_{q^s}$ respectively. 

By Theorem~\ref{mainth1} and Theorem~\ref{thm:BF},
for $d\in(\Z/7\Z)^{\times}$, 
generalized Jacobi sums
\[
J^\ast_{p^rq^s,d}(1,1)
=-\sigma_{-d}(J^\ast_{p^r}(1,1))J^\ast_{q^s}(1,1)=\sum_{k=1}^{7}c_k\zeta_7^k
\]
for $p^rq^s$ are given by 
$c_k=-(\sum_{i=1}^{6}a_{-d^{-1}i}b_{k-i} - \sum_{i=1}^{6}a_{-d^{-1}i}b_{-i})$
and the $c_k$'s satisfy
\begin{align}\label{eq:DS7c}
\begin{dcases}
p^rq^s=
c_1^2+c_2^2+c_3^2+c_4^2+c_5^2+c_6^2- c_1c_2-c_2c_3-c_3c_4-c_4c_5-c_5c_6,\\
c_1c_2+c_2c_3+c_3c_4+c_4c_5+c_5c_6-(c_1c_3+c_2c_4+c_3c_5+c_4c_6+c_6c_1)=0,\\
c_1c_2+c_2c_3+c_3c_4+c_4c_5+c_5c_6-(c_1c_4+c_2c_5+c_3c_6+c_5c_1+c_6c_2)=0.
\end{dcases}
\end{align}
We define
\begin{align}\label{eq:BF7}
\begin{split}
z_{1,d}&:=c_1+c_2+c_3+c_4+c_5+c_6\\
&=-\frac{1}{12}\sigma^i(2 x_1 y_1 + 42 x_2 y_2 + 42 x_3 y_3 + 42 x_4 y_4 + 343 x_5 y_5 + 1029 x_6 y_6),\\
z_{2,d}&:=-(c_1-c_6)\\
&=-\frac{1}{12}\sigma^i(
2 x_1 y_2 
- 2 x_2 y_1 
+ 7 x_2 y_5 
- 21 x_2 y_6
- 21 x_3 y_5 
- 21 x_3 y_6
- 21 x_4 y_5 
+ 21 x_4 y_6\\
&\quad 
- 7 x_5 y_2
+ 21 x_5 y_3 
+ 21 x_5 y_4 
+ 21 x_6 y_2 
+ 21 x_6 y_3
- 21 x_6 y_4
),\\
z_{3,d}&:=-(c_2-c_5)\\
&=-\frac{1}{12}\sigma^i(
2 x_1 y_3 
- 21 x_2 y_5
- 21 x_2 y_6
- 2 x_3 y_1
- 14 x_3 y_5
- 42 x_4 y_6
+ 21 x_5 y_2
+ 14 x_5 y_3\\
&\quad 
+ 21 x_6 y_2
+ 42 x_6 y_4 
),\\
z_{4,d}&:=-(c_3-c_4)\\
&=-\frac{1}{12}\sigma^i(
2 x_1 y_4
- 21 x_2 y_5
+ 21 x_2 y_6
- 42 x_3 y_6
-2 x_4 y_1
+ 7 x_4 y_5
+ 21 x_4 y_6
+ 21 x_5 y_2
- 7 x_5 y_4\\
&\quad 
- 21 x_6 y_2
+ 42 x_6 y_3
- 21 x_6 y_4
 ),\\
z_{5,d}&:=-\frac{1}{7}(c_1+c_2-2c_3-2c_4+c_5+c_6)\\
&=-\frac{1}{168}\sigma^i(
28 x_1 y_5
- 12 x_2 y_2
+ 36 x_2 y_3
+ 36 x_2 y_4
+ 36 x_3 y_2
+ 24 x_3 y_3
+ 36 x_4 y_2
- 12 x_4 y_4 \\
&\quad 
+ 28 x_5 y_1
- 49 x_5 y_5
+ 441 x_5 y_6
+ 441 x_6 y_5
+ 147 x_6 y_6
),\\
z_{6,d}&:=-\frac{1}{7}(c_1-c_2-c_5+c_6)\\
&=-\frac{1}{168}\sigma^i(
28 x_1 y_6
+ 12 x_2 y_2
+ 12 x_2 y_3
- 12 x_2 y_4
+ 12 x_3 y_2
+ 24 x_3 y_4 \\
&\quad 
- 12 x_4 y_2
+ 24 x_4 y_3
 - 12 x_4 y_4 
+ 147 x_5 y_5
+ 49 x_5 y_6 
+ 28 x_6 y_1
+ 49 x_6 y_5
- 441 x_6 y_6
) 
\end{split}
\end{align}
where $\sigma(x_1,x_2,x_3,x_4,x_5,x_6,y_1,y_2,y_3,y_4,y_5,y_6)=(x_1,-x_3,x_4,x_2,\frac{-x_5-3x_6}{2},\frac{x_5-x_6}{2},y_1,y_2,y_3,y_4,y_5,y_6)$ and $i=0$ 
(resp. $2,1,4,5,3$)
if $d=1$ (resp. $2,3,4,5,6$). 
By the equation (\ref{eq:DS7c}), 
we can verify that 
$z_{1,d},z_{2,d},z_{3,d}, z_{4,d},z_{5,d},z_{6,d}$ satisfy
\begin{align*}
&2z_{1,d}^2+42z_{2,d}^2+42z_{3,d}^2+42z_{4,d}^2+343z_{5,d}^2+1029z_{6,d}^2\\
&=72(c_1^2+ c_2^2 + c_3^2  + c_4^2 + c_5^2 + c_6^2)\\
&\quad  - 24(
(c_1c_2 + c_2c_3 + c_3c_4 + c_4c_5 + c_5c_6)+
(c_1c_3 + c_2c_4 + c_3c_5 + c_4c_6 + c_6c_1)\\
&\quad
+(c_1c_4 + c_2c_5 + c_3c_6 + c_5c_1 + c_6c_2))\\
&=72(c_1^2+ c_2^2 + c_3^2  + c_4^2 + c_5^2 + c_6^2- (c_1c_2 + c_2c_3 + c_3c_4 + c_4c_5 + c_5c_6))\\
&=72p^rq^s,
\end{align*}
and
\begin{align*}
&12z_{2,d}^2-12z_{4,d}^2+147z_{5,d}^2-441z_{6,d}^2+56z_{1,d}z_{6,d}+24z_{2,d}z_{3,d}-24z_{2,d}z_{4,d}+48z_{3,d}z_{4,d}+98z_{5,d}z_{6,d}\\
&=48(c_1c_2 + c_2c_3 + c_3c_4 + c_4c_5  + c_5c_6 - (c_1c_3 + c_2c_4 + c_3c_5 + c_4c_6  + c_6c_1))\\
&=0,\\
&12z_{3,d}^2-12z_{4,d}^2+49z_{5,d}^2-147z_{6,d}^2+28z_{1,d}z_{5,d}+28z_{1,d}z_{6,d}+48z_{2,d}z_{3,d}+24z_{2,d}z_{4,d}+24z_{3,d}z_{4,d}\\
&+490z_{5,d}z_{6,d}\\
&=48(c_1c_2 + c_2c_3 + c_3c_4 + c_4c_5 + c_5c_6 - (c_1c_4 + c_2c_5 + c_3c_6 + c_5c_1 + c_6c_2))\\
&=0.
\end{align*}
Thus $(z_{1,d},z_{2,d},z_{3,d},z_{4,d},z_{5,d},z_{6,d})$ yields solutions of the system of Diophantine equations 
for $p^rq^s$:
\begin{align*}
\begin{cases}
72p^rq^s=2z_{1,d}^2+42z_{2,d}^2+42z_{3,d}^2+42z_{4,d}^2+343z_{5,d}^2+1029z_{6,d}^2, \\
12z_{2,d}^2-12z_{4,d}^2+147z_{5,d}^2-441z_{6,d}^2+56z_{1,d}z_{6,d}+24z_{2,d}z_{3,d}
-24z_{2,d}z_{4,d}+48z_{3,d}z_{4,d}\\
+98z_{5,d}z_{6,d}=0, \\
12z_{3,d}^2-12z_{4,d}^2+49z_{5,d}^2-147z_{6,d}^2+28z_{1,d}z_{5,d}+28z_{1,d}z_{6,d}+48z_{2,d}z_{3,d}+24z_{2,d}z_{4,d}+24z_{3,d}z_{4,d}\\
+490z_{5,d}z_{6,d}=0,\\
z_{1,d}\equiv 1\ ({\rm mod}\ 7).
\end{cases}
\end{align*}

For
\[
{\bf z}^{\prime}=(z_{1,d}^{\prime},z_{2,d}^{\prime},z_{3,d}^{\prime},z_{4,d}^{\prime},z_{5,d}^{\prime},z_{6,d}^{\prime})
=(6z_{1,d},6z_{2,d},6z_{3,d},6z_{4,d},6z_{5,d},6z_{6,d}),
\]
consider the quadratic form
\[
q({\bf x})=2x_1^2+42x_2^2+42x_3^2+42x_4^2+343x_5^2+1029x_6^2
\]
on the variety $V$ defined by
\begin{align*}
\begin{dcases}
12x_2^2-12x_4^2+147x_5^2-441x_6^2+56x_1x_6+24x_2x_3
-24x_2x_4+48x_3x_4+98x_5x_6=0, \\
12x_3^2-12x_4^2+49x_5^2-147x_6^2+28x_1x_5+28x_1x_6+48x_2x_3+24x_2x_4+24x_3x_4+490x_5x_6=0.
\end{dcases}
\end{align*}
Then $q({\bf x})$ is not multiplicative on $V$ with respect to
\[
\varphi_d({\bf x},{\bf y})=(z_{1,d}^{\prime},z_{2,d}^{\prime},z_{3,d}^{\prime},z_{4,d}^{\prime},z_{5,d}^{\prime},z_{6,d}^{\prime}).
\]
Indeed, the point ${\bf z}^{\prime}$ still satisfies the defining equations of $V$:
\begin{align*}
\begin{cases}
12{(z^{\prime}_2)}^2-12{(z^{\prime}_4)}^2+147{(z^{\prime}_5)}^2-441{(z^{\prime}_6)}^2+56z^{\prime}_1z^{\prime}_6
+24z^{\prime}_2z^{\prime}_3-24z^{\prime}_2z^{\prime}_4+48z^{\prime}_3z^{\prime}_4+98z^{\prime}_5z^{\prime}_6=0, \\
12{(z^{\prime}_3)}^2-12{(z^{\prime}_4)}^2+49{(z^{\prime}_5)}^2-147{(z^{\prime}_6)}^2+28z^{\prime}_1z^{\prime}_5+28z^{\prime}_1z^{\prime}_6+48z^{\prime}_2z^{\prime}_3+24z^{\prime}_2z^{\prime}_4+24z^{\prime}_3z^{\prime}_4+490z^{\prime}_5z^{\prime}_6=0,
\end{cases}
\end{align*}
but
\begin{align*}
q(\varphi_d({\bf x},{\bf y}))=q({\bf z}^{\prime})
&=2{(z_{1,d}^{\prime})}^2+42{(z_{2,d}^{\prime})}^2+42{(z_{3,d}^{\prime})}^2+42{(z_{4,d}^{\prime})}^2+343{(z_{5,d}^{\prime})}^2+1029{(z_{6,d}^{\prime})}^2\\
&=36(2z_{1,d}^2+42z_{2,d}^2+42z_{3,d}^2+42z_{4,d}^2+343z_{5,d}^2+1029z_{6,d}^2)\\
&=36(72p^rq^s)\\
&=\frac{1}{2}(72p^r)(72q^s)\\
&=\frac{1}{2}q({\bf x})q({\bf y})
\neq q({\bf x})q({\bf y}).
\end{align*}

This suggests renormalizing the quadratic form. Define
\[
q^\prime({\bf x}):=\frac{1}{2}q({\bf x})
=x_1^2+21x_2^2+21x_3^2+21x_4^2+\frac{343}{2}x_5^2+\frac{1029}{2}x_6^2.
\]
Then
\begin{align*}
q^\prime(\varphi_d({\bf x},{\bf y}))
&=q^\prime({\bf z}^{\prime})
=\frac{1}{2}q({\bf z}^{\prime})\\
&={(z_{1,d}^{\prime})}^2+21{(z_{2,d}^{\prime})}^2+21{(z_{3,d}^{\prime})}^2+21{(z_{4,d}^{\prime})}^2+\frac{343}{2}{(z_{5,d}^{\prime})}^2+\frac{1029}{2}{(z_{6,d}^{\prime})}^2\\
&=\frac{1}{2}\left(\frac{1}{2}q({\bf x})q({\bf y})\right)\\
&=\left(\frac{1}{2}q({\bf x})\right)\left(\frac{1}{2}q({\bf y})\right)\\
&=q^\prime({\bf x})q^\prime({\bf y}).
\end{align*}
Hence $q^\prime({\bf x})$ is multiplicative on $V$ with respect to $\varphi_d({\bf x},{\bf y})=(z_{1,d}^{\prime},z_{2,d}^{\prime},z_{3,d}^{\prime},z_{4,d}^{\prime},z_{5,d}^{\prime},z_{6,d}^{\prime})$.

To compare this with Hoshi~\cite[Theorem 6]{Hos03}, we now specialize to $d=1$ and write
\begin{align*}
{\bf X}&=(X_1,X_2,X_3,X_4,X_5,X_6):=\left(x_1,x_2,x_3,x_4,\frac{7}{2}x_5,\frac{7}{2}x_6\right),\\
{\bf Y}&=(Y_1,Y_2,Y_3,Y_4,Y_5,Y_6):=\left(y_1,y_2,y_3,y_4,\frac{7}{2}y_5,\frac{7}{2}y_6\right),\\
{\bf Z}&=(Z_1,Z_2,Z_3,Z_4,Z_5,Z_6)
:=(z_{1,1}^{\prime},z_{2,1}^{\prime},z_{3,1}^{\prime},z_{4,1}^{\prime},\tfrac{7}{2}z_{5,1}^{\prime},\tfrac{7}{2}z_{6,1}^{\prime})\\
&\phantom{~=(Z_1,Z_2,Z_3,Z_4,Z_5,Z_6)}=(6z_{1,1},6z_{2,1},6z_{3,1},6z_{4,1},21z_{5,1},21z_{6,1}).
\end{align*}
Then
\[
q^\prime({\bf X})=X_1^2+21X_2^2+21X_3^2+21X_4^2+14X_5^2+42X_6^2,
\]
and
\begin{align*}
\begin{split}
Z_1
&=-(X_1Y_1+21X_2Y_2+21X_3Y_3+21X_4Y_4+14X_5Y_5+42X_6Y_6),\\
Z_2
&=-(X_1Y_2-X_2Y_1+X_2Y_5-3X_2Y_6-3X_3Y_5-3X_3Y_6-3X_4Y_5+X_4Y_6-X_5Y_2\\
&\qquad\quad +3X_5Y_3+3X_5Y_4+3X_6Y_2+3X_6Y_3-3X_6Y_4),\\
Z_3
&=-(X_1Y_3-3X_2Y_5-3X_2Y_6-X_3Y_1-2X_3Y_5-6X_4Y_6+3X_5Y_2+2X_5Y_3\\
&\qquad\quad +3X_6Y_2+6X_6Y_4),\\
Z_4
&=-(X_1Y_4-3X_2Y_5+3X_2Y_6-6X_3Y_6-2X_4Y_1+X_4Y_5+3X_4Y_6+3X_5Y_2-X_5Y_4\\
&\qquad\quad -3X_6Y_2+6X_6Y_3-3X_6Y_4),\\
Z_5
&=-\frac{1}{2}(2X_1Y_5-3X_2Y_2+9X_2Y_3+9X_2Y_4+9X_3Y_2+6X_3Y_3+9X_4Y_2-3X_4Y_4\\
&\qquad\qquad +2X_5Y_1-X_5Y_5+9X_5Y_6+9X_6Y_5+3X_6Y_6),\\
Z_6
&=-\frac{1}{2}(2X_1Y_6+3X_2Y_2+3X_2Y_3-3X_2Y_4+3X_3Y_2+6X_3Y_4-3X_4Y_2+6X_4Y_3\\
&\qquad\qquad -3X_4Y_4+3X_5Y_5+X_5Y_6+2X_6Y_1+X_6Y_5-9X_6Y_6).
\end{split}
\end{align*}
Therefore we recover the multiplicative quadratic form $q^\prime({\bf X})$ on the variety $V$
defined by
\begin{align*}
\begin{dcases}
3X_2^2-3X_4^2+3X_5^2-9X_6^2+4X_1X_6+6X_2X_3-6X_2X_4+12X_3X_4+2X_5X_6=0, \\
3X_3^2-3X_4^2+X_5^2-3X_6^2+2X_1X_5+2X_1X_6+12X_2X_3+6X_2X_4+6X_3X_4+10X_5X_6=0
\end{dcases}
\end{align*}
with $\varphi_1({\bf X},{\bf Y})=(Z_1,Z_2,Z_3,Z_4,Z_5,Z_6)$,
which coincides with the one given by Hoshi \cite[Section 2, Theorem~6]{Hos03}. (There is a typo in \cite[p.~74, line 18]{Hos03}: the term $-3X_6Y_6$ in the sixth component should be moved to the fifth component.)

In the case
$p=q$ (i.e. $x_i=y_i$), 
a lift of Jacobi sums $J^\ast_{p^r}(1,1)$ to $J^\ast_{p^{nr}}(1,1)$ 
can be obtained by Davenport and Hasse's lifting theorem 
(see Theorem~\ref{thm:corMT1} and Section \ref{S2}). 
Hoshi and Kanai \cite[Section 7, (3) $e=7$]{HK22} 
gave such an explicit lift of solutions of the system of Diophantine equations (\ref{eq:DS7p}) corresponding to Davenport and Hasse's lifting theorem, cf. Theorem~\ref{thm:corMT1} with 
$d=6=-1 \in(\Z/7\Z)^{\times}$. 
Indeed,   
\begin{align*}
z_{1,-1} &= -\frac{1}{12}(2x_1^2 - 42x_2^2 - 42x_3^2 - 42x_4^2 + 343x_5^2 +1029x_6^2),\\
z_{2,-1} &= -\frac{1}{12}(4x_1x_2 - 14x_2x_5 + 42x_3x_5 + 42x_4x_5 + 42x_2x_6 + 42x_3x_6 - 42x_4x_6),\\
z_{3,-1} &= -\frac{1}{12}(4x_1x_3 + 42x_2x_5 + 28x_3x_5 + 42x_2x_6 + 84x_4x_6),\\
z_{4,-1} &= -\frac{1}{12}(4x_1x_4 + 42x_2x_5 - 14x_4x_5 - 42x_2x_6 + 84x_3x_6 - 42x_4x_6),\\
z_{5,-1} &= -\frac{1}{168}(12x_2^2 - 72x_2x_3 - 24x_3^2 - 72x_2x_4 + 12x_4^2 + 56x_1x_5 - 49x_5^2 + 882x_5x_6 + 147x_6^2),\\
z_{6,-1} &= -\frac{1}{168}(-12x_2^2 - 24x_2x_3 + 24x_2x_4 - 48x_3x_4 + 12x_4^2 + 147x_5^2 + 56x_1x_6 + 98x_5x_6 - 441x_6^2)
\end{align*}
are given as $x_1^{(2)}$, $x_2^{(2)}$, $x_3^{(2)}$, 
$x_4^{(2)}$, $x_5^{(2)}$, $x_6^{(2)}$ 
in Hoshi and Kanai \cite[page 24, lines 5--10]{HK22}. 

\begin{remark}
A similar result can be obtained for $l=11$ 
by using Leonard and Williams \cite{LW75b}.
\end{remark}

\section{Proof of Theorem~\ref{mainth3}: Multiplicative $f$-ic forms on algebraic varieties}\label{S5}

In this section, we construct multiplicative $f$-ic forms on algebraic $K$-varieties for $f \ge 2$. 
The defining equations of $V$ are recast in terms of correlation sums arising from the cyclotomic setup, and this formulation is compatible with the multiplicative formula of Theorem~\ref{thm:BF2}, which yields the induced composition law and shows that $V$ is stable under this law. 
The regularity of the form $h$ is then obtained from its derived form and its relation to the norm form. 
This completes the proof of Theorem~\ref{mainth3}.

Let $h(x)=h(x_1,\dots,x_n)$ be an $f$-ic form, i.e. a homogeneous polynomial of total degree $f$ in $n$ variables. 
Assume that $h$ is regular, i.e. that the associated symmetric $f$-linear map $\theta:K^n\times\cdots\times K^n\to K$ has the property that $\theta({\bf v}_1,{\bf v}_2,\dots,{\bf v}_f)=0$ for any ${\bf v}_2,\dots, {\bf v}_f$ implies ${\bf v}_1=0$
(see Shapiro \cite[Chapter 16]{Sha00}, Schafer \cite{Sch70}, cf. Section \ref{S4} for $f=2$).

\begin{defn}
Let $f \geq 2$ be an integer and $K$ be a field with char $K\notmid f!$.  
Let $V \subset K^n$ be an algebraic $K$-variety. 
We say a regular $f$-ic form $h({\bf x})$ is {\it multiplicative on $V$} 
if there exists a bilinear map 
$\varphi :K^n \times K^n \rightarrow K^n$ such that
\[
\varphi(V \times V) \subset V \quad \mbox{and}\quad h({\bf x})h({\bf y})=h(\varphi({\bf x,y}))\mbox{ for any }{\bf x,y} \in V.
\]
\end{defn}

\begin{lem} \label{lem:BEW2}
Let $f\geq2$ be an integer and $l$ be an odd prime with $l \equiv 1\ ({\rm mod}\ f)$. 
Write $l=ef+1$.
Let $K$ be a field with char $K\notmid f!$, $l$ and 
$\Gal(K(\zeta_l)/K)\simeq (\Z/l\Z)^{\times}=\langle \gamma \rangle$. 
Define $\beta := \gamma^e$ and $M=K(\zeta_l)^{\langle \beta \rangle}$, 
i.e., $\Gal(K(\zeta_l)/M)\simeq \Z/f\Z = \langle \beta \rangle \leq  (\Z/l\Z)^{\times}$.
Let ${\bf p}\in K^{\times}$ and
\[
\alpha_{\bf p}=\sum_{k=0}^{l-1}a_k\zeta_l^k \quad (a_k \in K)\\
\quad \text{with} \quad
N_{K(\zeta_l)/M}(\alpha_{\bf p})= {\bf p}.
\]
Then we have
\begin{enumerate}
\item [$(1)$] ${\bf p} =S_0 - S_1$,
\item [$(2)$] $S_1=S_2=\cdots =S_{e}$
\end{enumerate}
where 
\begin{align*}
S_m = 
\left\{
\begin{array}{ll}
\displaystyle
\sum_{i_0 + \beta i_1 + \cdots +\beta^{f-1}i_{f-1}\equiv 0 \ (l)}a_{i_0}a_{i_1}\cdots a_{i_{f-1}} & (m = 0)\\
\displaystyle
\sum_{i_0 + \beta i_1 + \cdots +\beta^{f-1}i_{f-1}\equiv \gamma^{m} \ (l)}a_{i_0}a_{i_1}\cdots a_{i_{f-1}} & (m = 1,\ldots,e)
\end{array}
\right.
\end{align*}
and each subscript of the $a_k$'s should be taken modulo $l$. 
Moreover, we may assume that $a_0=0$ without loss of generality.
\end{lem}
\begin{proof}
By $N_{K(\zeta_l)/M}(\alpha_{\bf p})= {\bf p}$ and $[K(\zeta_l):M]=f$, we have 
\begin{align*}
{\bf p} = \Big( a_1\zeta_l +  \cdots +a_{l-1}\zeta_l^{l-1}\Big) \left( a_1\zeta_l^{\beta} +  \cdots +a_{l-1}\zeta_l^{\beta(l-1)}\right) \cdots \left( a_1\zeta_l^{\beta^{f-1}} +  \cdots +a_{l-1}\zeta_l^{\beta^{f-1}(l-1)}\right) ,
\end{align*}
that is,
\begin{align*}
{\bf p} 
&= S_0 + S_1\eta(0) + \cdots + S_{e-1}\eta(e-2) + S_e \eta(e-1)
\end{align*}
where, for $0\leq i \leq e-1$, $\eta(i)=\sum_{j=0}^{f-1}\zeta_l^{\gamma ^{ej+i}}$ are Gaussian $f$-periods of degree $e$.
Note that  $S_j =S_{j+ei} (1\leq i \leq e-1)$ 
because multiplying the equation $i_0 + \beta i_1 + \cdots +\beta^{f-1}i_{f-1}=\gamma^j$ by $\beta^i(=\gamma^{ei})$, 
the variables only change their order.
By $\zeta_l+\cdots+\zeta_l^{l-1}=-1$, we have
\[
(S_1 - S_0 + {\bf p})\eta(1)  + 
 \cdots + (S_{e-1} - S_0 + {\bf p})\eta(e-1) + (S_e - S_0 + {\bf p})\eta(0) = 0.
\]
Because $\{\eta(0), \ldots, \eta(e-1)\}$ forms a normal basis of $M/K$, 
we have
\[
S_1 - S_0 + {\bf p} = S_m - S_0 + {\bf p}=0\ \ (1\leq m\leq e).
\]
This implies that 
(1) ${\bf p}= S_0 - S_1$ and (2) $S_1=S_m$ 
$(2\leq m\leq e)$. 
\end{proof}

\begin{thm}\label{thm:BF2}
Let $f\geq2$ be an integer. 
Let $l$ be an odd prime with $l \equiv 1\ ({\rm mod}\ f)$ and $d \in (\Z/l\Z)^{\times}$. 
Write $l=ef+1$.
Let $K$ be a field with char $K\notmid f$, $l$ and 
$\Gal(K(\zeta_l)/K)\simeq (\Z/l\Z)^{\times}=\langle \gamma \rangle$. 
Define $\beta := \gamma^e$ and $M=K(\zeta_l)^{\langle \beta \rangle}$, 
i.e., $\Gal(K(\zeta_l)/M)\simeq \Z/f\Z = \langle \beta \rangle \leq  (\Z/l\Z)^{\times}$.
For ${\bf p}_1, {\bf p}_2\in K^{\times}$ we define
\begin{align*}
\alpha_{{\bf p}_1}=\sum_{i=0}^{l-1}a_i\zeta_l^i\quad \text{with} \quad
N_{K(\zeta_l)/M}(\alpha_{{\bf p}_1})= {\bf p}_1, \quad \alpha_{{\bf p}_2}=\sum_{j=0}^{l-1}b_j\zeta_l^j\quad \text{with} \quad
N_{K(\zeta_l)/M}(\alpha_{{\bf p}_2})= {\bf p}_2.
\end{align*}
Define also $\alpha_{{\bf p}_1{\bf p}_2,d}:=(-1)^{f-1}\sigma_{-d}(\alpha_{{\bf p}_1})\alpha_{{\bf p}_2}$ where $\sigma_{-d}\in \Gal(K(\zeta_l)/K)$ satisfies $\sigma_{-d}(\zeta_l)=\zeta_l^{-d}$.
Then $\alpha_{{\bf p}_1{\bf p}_2,d} = \sum_{k=0}^{l-1} c_k\zeta_l^k$ is given by
$$c_k = (-1)^{f-1}\sum_{i=0}^{l-1}a_{-d^{-1}i}b_{k-i},$$
and these coefficients satisfy
\begin{enumerate}
\item ${\bf p}_1{\bf p}_2 = S_0 - S_1$,
\item $S_1 = S_2 = \cdots = S_e$.
\end{enumerate}
Here, the sums $S_m$ are defined as
\begin{align*}
S_m = 
\left\{
\begin{array}{ll}
\displaystyle
\sum_{i_0 + \beta i_1 + \cdots +\beta^{f-1}i_{f-1}\equiv 0 \ (l)}c_{i_0}c_{i_1}\cdots c_{i_{f-1}} & (m = 0)\\
\displaystyle
\sum_{i_0 + \beta i_1 + \cdots +\beta^{f-1}i_{f-1}\equiv \gamma^{m} \ (l)}c_{i_0}c_{i_1}\cdots c_{i_{f-1}} & (m = 1,\ldots,e),
\end{array}
\right.
\end{align*}
with all indices taken modulo $l$.
Moreover, we may assume $a_0=b_0=c_0=0$ 
without loss of generality and in this case 
$\alpha_{{\bf p}_1{\bf p}_2,d}=\sum_{k=1}^{l-1}c_k\zeta_l^k$ 
is  given by 
$$c_k=(-1)^{f-1}\left( \sum_{i=1}^{l-1}a_{-d^{-1}i}b_{k-i} - \sum_{i=1}^{l-1}a_{-d^{-1}i}b_{-i}\right).$$
\end{thm}
\begin{proof}
By the definition of $\alpha_{{\bf p}_1}$ and $\alpha_{{\bf p}_2}$, we have
\begin{align*}
N_{K(\zeta_l)/M}(\alpha_{{\bf p}_1{\bf p}_2,d})
=N_{K(\zeta_l)/M}((-1)^{f-1}\sigma_{-d}(\alpha_{{\bf p}_1})
\alpha_{{\bf p}_2})
=(-1)^{f(f-1)}{\bf p}_1{\bf p}_2={\bf p}_1{\bf p}_2.
\end{align*}
Using a calculation similar to that in the proof of Theorem~\ref{thm:BF}, we obtain
\begin{align*}
\alpha_{{\bf p}_1{\bf p}_2,d}=(-1)^{f-1}\sigma_{-d}(\alpha_{{\bf p}_1})\alpha_{{\bf p}_2}=(-1)^{f-1}\Bigl(\sum_{i=1}^{l-1}a_{-d^{-1}i}\zeta_l^i\Bigr) \Bigl(\sum_{j=1}^{l-1}b_j\zeta_l^j\Bigr)=\sum_{k=0}^{l-1}c_k\zeta_l^k.
\end{align*}
where $c_k=(-1)^{f-1}\sum_{i=0}^{l-1}a_{-d^{-1}i}b_{k-i}$. 
Thus the assertion follows from Lemma \ref{lem:BEW2}.

The case $a_0=b_0=0$ also follows from the same argument as in the proof of Theorem~\ref{thm:BF}, yielding $c_k=(-1)^{f-1}(\sum_{i=1}^{l-1}a_{-d^{-1}i}b_{k-i} - \sum_{i=1}^{l-1}a_{-d^{-1}i}b_{-i})$ 
when $\alpha_{{\bf p}_1{\bf p}_2,d}=\sum_{k=1}^{l-1}c_k\zeta_l^k$.
\end{proof}

To consider the regularity of the $f$-ic form $h({\bf x})$, we define the associated symmetric $f$-linear map $\theta_h$ explicitly.
Following Dr\'apal and Vojt\v{e}chovsk\'y \cite[\S2.1]{DV09}, for a map $P\colon V\to F$
with $P(0)=0$ we define the $f$-th defect (derived form)
\[
\Delta^{f}P(u_1,\dots,u_f)
=\sum_{\emptyset\neq I\subset\{1,\dots,f\}}(-1)^{f-|I|}
\,P\Big(\sum_{i\in I}u_i\Big),
\qquad
\theta_P:=\frac{1}{f!}\Delta^{f}P.
\]
If char $F\nmid f!$ and $P$ is a homogeneous polynomial map of degree $f$, then
$\theta_P$ is a symmetric $f$-linear form and $P(x)=\theta_P(x,\dots,x)$.

\begin{prop}\label{prop:norm_polarization}
Let $E/F$ be a separable field extension of degree $f$, and assume that char $F \nmid f!$. 
Let $N_{E/F}\colon E\to F$ be the norm map.
Then
\[
\theta_N(u_1,\dots,u_f):=\frac{1}{f!}\,\Delta^f N_{E/F}(u_1,\dots,u_f)
\]
is a symmetric $f$-linear form on $E$.
\end{prop}

\begin{proof}
Regard $E$ as an $F$-vector space of dimension $f$. 
For $x\in E$, let $m_x\colon E\to E$ be multiplication by $x$. 
Then $x\mapsto m_x$ is $F$-linear and
\[
N_{E/F}(x)=\det(m_x),
\]
so $N_{E/F}$ is a homogeneous polynomial map of degree $f$ on $E$.

Now apply the basic properties of derived forms recalled above.
Since $N_{E/F}$ is a homogeneous polynomial map of degree $f$ and
char $F\nmid f!$, it follows from \cite[\S2.1]{DV09} that the $f$-th defect
$\Delta^f N_{E/F}$ is symmetric and $f$-additive, hence $(1/f!)\Delta^f N_{E/F}$ is $F$-linear
in each argument. 
Therefore $\theta_N$ is a symmetric $f$-linear form on $E$.
\end{proof}

\begin{lem}\label{lem:h_regularity}
With the notation of Theorem~\ref{mainth3}, the $f$-ic form $h({\bf x})$ on $K^{l-1}$ is regular.
\end{lem}

\begin{proof}
Set $L=K(\zeta_l)$ and $M=L^{\langle\beta\rangle}$, so $[L:M]=f$.
Via the $K$-linear isomorphism
\[
\iota:K^{l-1}\xrightarrow{\ \sim\ }L,\qquad
{\bf x}=(x_1,\dots,x_{l-1})\longmapsto
\alpha_{\bf x}:=\sum_{i=1}^{l-1}x_i\zeta_l^i,
\]
we identify $K^{l-1}$ with $L$. Let $\sigma\in\Gal(L/M)$ be the generator determined by
$\sigma(\zeta_l)=\zeta_l^\beta$. Then
\[
N_{L/M}(z)=\prod_{t=0}^{f-1}\sigma^t(z)\qquad(z\in L).
\]

\medskip\noindent
\textbf{Step 1: $h$ as a $K$-linear projection of the norm.}
As in the proof of Lemma~\ref{lem:BEW2}, expanding $N_{L/M}(\alpha_{\bf x})$ and grouping terms by
$\langle\beta\rangle$-orbits yields
\begin{equation}\label{eq:norm_period_expansion}
N_{L/M}(\alpha_{\bf x})
= S_0({\bf x})+\sum_{m=1}^{e}S_m({\bf x})\,\eta(m-1)
\end{equation}
where $S_m({\bf x})$ are the polynomials defined in Lemma~\ref{lem:BEW2} (with $x_0=0$ and indices modulo $l$).
Since $\eta(0)+\cdots+\eta(e-1)=\zeta_l+\cdots+\zeta_l^{l-1}=-1$, we can rewrite \eqref{eq:norm_period_expansion} as
\[
N_{L/M}(\alpha_{\bf x})
=\sum_{r=0}^{e-1}\bigl(S_{r+1}({\bf x})-S_0({\bf x})\bigr)\,\eta(r).
\]
Define a $K$-linear projection $\lambda:M\to K$ by
\[
\lambda(\eta(0))=-1,\qquad \lambda(\eta(r))=0\ \ (1\le r\le e-1).
\]
Then applying $\lambda$ gives
\[
\lambda\bigl(N_{L/M}(\alpha_{\bf x})\bigr)=-(S_1({\bf x})-S_0({\bf x}))=h({\bf x}),
\]
hence
\begin{equation}\label{eq:h_as_proj_norm}
h=\lambda\circ N_{L/M}\circ\iota.
\end{equation}

\medskip\noindent
\textbf{Step 2: Polarization and a trace specialization.}
Let $\theta_N$ be the symmetric $f$-linear form associated with $N_{L/M}$ over the base field $M$
(Proposition~\ref{prop:norm_polarization} applied to $E=L$, $F=M$), and let $\theta_h$ be the symmetric
$f$-linear form associated with $h$.
By \eqref{eq:h_as_proj_norm} and linearity of polarization,
\begin{equation}\label{eq:theta_h_comp}
\theta_h=\lambda\circ\theta_N\circ(\iota\times\cdots\times\iota).
\end{equation}
A standard computation from $N_{L/M}(z)=\prod_{t=0}^{f-1}\sigma^t(z)$ gives, for $\alpha,v\in L$,
\begin{equation}\label{eq:thetaN_trace}
\theta_N(v,\alpha,\dots,\alpha)
=\frac{1}{f}\sum_{k=0}^{f-1}\sigma^k(v)\prod_{j\ne k}\sigma^j(\alpha)
=\frac{1}{f}\Tr_{L/M}\!\bigl(vP(\alpha)\bigr),
\qquad
P(\alpha):=\prod_{j=1}^{f-1}\sigma^j(\alpha).
\end{equation}

\medskip\noindent
\textbf{Step 3: Regularity.}
It suffices to show that if $\alpha \in L$ satisfies $\theta_h(v,\alpha,\dots,\alpha)=0$ for all $v\in L$,
then $\alpha=0$.
Indeed, assuming this, let $u\in L$ be such that
$\theta_h(u,w_1,\dots,w_f)=0$ for all $w_1,\dots,w_f\in L$.
Then, in particular, by taking $w_1=v$ and $w_2=\cdots=w_f=u$, we obtain
$\theta_h(u,v,u,\dots,u)=0$ for all $v\in L$.
By symmetry of $\theta_h$, this implies
$\theta_h(v,u,\dots,u)=0$ for all $v\in L$.
Hence $u=0$.

Assume $\alpha\in L$ satisfies $\theta_h(v,\alpha,\dots,\alpha)=0$ for all $v\in L$.
Using \eqref{eq:theta_h_comp} and \eqref{eq:thetaN_trace} we obtain
\[
\lambda\!\left(\Tr_{L/M}(vP(\alpha))\right)=0\qquad(\forall v\in L).
\]
Set $\Psi:=\lambda\circ\Tr_{L/M}:L\to K$. 
Then $\Psi\neq 0$: indeed, $\eta(0)+\cdots+\eta(e-1)=-1$ implies $\lambda(1)=1$, hence
\[
\Psi(1)=\lambda\!\left(\Tr_{L/M}(1)\right)=\lambda(f\cdot 1)=f\neq 0
\]
since char $K\nmid f$.
Thus $\Psi(vP(\alpha))=0$ for all $v\in L$.

Because $L/K$ is finite separable, the trace pairing $(x,y)\mapsto\Tr_{L/K}(xy)$ is non-degenerate.
Therefore there exists $c\in L^\times$ such that $\Psi(z)=\Tr_{L/K}(cz)$ for all $z\in L$. 
Consequently,
\[
0=\Psi(vP(\alpha))=\Tr_{L/K}(c\,v\,P(\alpha))\qquad(\forall v\in L),
\]
and non-degeneracy forces $cP(\alpha)=0$, hence $P(\alpha)=0$. Since $L$ is a field and
$P(\alpha)=\prod_{j=1}^{f-1}\sigma^j(\alpha)$, we conclude $\alpha=0$.
Thus $\theta_h(v,\alpha,\dots,\alpha)\equiv 0$ implies $\alpha=0$, i.e.\ $h$ is regular.
\end{proof}

\begin{proof}[Proof of Theorem~\ref{mainth3}]
The regularity of the $f$-ic form $h({\bf x})$ follows from Lemma \ref{lem:h_regularity}.

By the definition of $V$ in Theorem~\ref{mainth3} and Lemma~\ref{lem:BEW2}, for ${\bf x}\in K^{l-1}$ we have
\[
{\bf x}\in V \quad\Longleftrightarrow\quad h_m({\bf x})=0\ (m=1,\dots,e-1)
\quad\Longleftrightarrow\quad S_1({\bf x})=S_2({\bf x})=\cdots=S_e({\bf x})
\]
where $S_j(x)$ denotes the $j$-th correlation sum of the coefficients of $\sum_{i=1}^{l-1}x_i\zeta_l^{i}$ (cf. proof of Lemma~\ref{lem:BEW2}).

Now let ${\bf x,y}\in V$ and define ${\bf z}=\varphi_d({\bf x,y})$.
By Theorem~\ref{thm:BF2}, the coefficients of $\alpha_{\bf z}=\sum_{i=1}^{l-1}c_i\zeta_l^i$ are given by $c_k=(-1)^{f-1}\left( \sum_{i=1}^{l-1}a_{-d^{-1}i}b_{k-i} - \sum_{i=1}^{l-1}a_{-d^{-1}i}b_{-i}\right)$, and in
particular the corresponding correlation sums satisfy
\[
S_1({\bf z})=S_2({\bf z})=\cdots=S_e({\bf z}).
\]
Hence ${\bf z}\in V$, i.e.\ $\varphi_d(V\times V)\subset V$.

Finally, Theorem~\ref{thm:BF2} (2) gives $h({\bf z})=h({\bf x})h({\bf y})$, i.e.
\[
h(\varphi_d({\bf x,y}))=h({\bf x})h({\bf y}),
\]
and Theorem~\ref{mainth3} follows.
\end{proof}

\begin{example}
[$l=5, f=4, e=1, \gamma =\beta =2$]
\begin{align*}
h({\bf x})=h(x_1,\ldots,x_4)&=S_0 - S_1
\end{align*}
where
\begin{align*}
S_0({\bf x})&= \sum_{s + 2t + 4u + 3v\equiv 0 \ (5)} x_{s}x_{t}x_{u}x_{v}\\
&= x_1^4 + x_2^4 + x_3^4 + x_4^4\\
&+ 4(
x_1x_2x_3x_4 
+x_1^2x_2x_3 
+ x_1x_2x_3^2 
+ x_1^2x_2x_4 
+ x_1x_2^2x_4 
+ x_1x_3^2x_4
+ x_1x_3x_4^2 
+ x_2^2x_3x_4 
+ x_2x_3x_4^2 
) \\
&+ 2(x_1^2x_2^2  
+ x_1^2x_3^2 
+ x_1^2x_4^2 
+ x_2^2x_3^2  
+ x_2^2x_4^2 
+ x_3^2x_4^2 ),\\
%
S_1({\bf x})&= \sum_{s + 2t + 4u + 3v\equiv 2 \ (5)} x_{s}x_{t}x_{u}x_{v}\\
&= 5x_1x_2x_3x_4 + x_1^2x_2^2 + x_1^2x_3^2 + x_1^2x_4^2 + x_2^2x_3^2+ x_2^2x_4^2 + x_3^2x_4^2 \\
&+ x_1x_2^3 + x_1x_3^3 + x_1x_4^3 
+ x_1^3x_2 + x_1^3x_3 + x_1^3x_4 
+ x_2x_3^3 + x_2x_4^3 
+ x_2^3x_3 + x_2^3x_4 
+ x_3^3x_4   + x_3x_4^3 \\
& + 2(x_1x_2x_3^2  +  x_1^2x_2x_3 
+ x_1x_2^2x_4  +  x_1^2x_2x_4 
+ x_1x_3^2x_4 + x_1x_3x_4^2 
+ x_2^2x_3x_4 + x_2x_3x_4^2 ) \\
&+ 3(x_1x_2^2x_3 + x_1x_2x_4^2 + x_1^2x_3x_4 + x_2x_3^2x_4 ).
%
\end{align*}
Thus
\begin{align*}
h({\bf x})&= - x_1x_2x_3x_4 
+ x_1^4 + x_2^4 + x_3^4  + x_4^4 
+ x_1^2x_3^2 + x_1^2x_2^2  + x_1^2x_4^2 + x_2^2x_3^2 + x_2^2x_4^2 + x_3^2x_4^2 \\
&- (x_1x_2^3 + x_1x_3^3 + x_1x_4^3  
+ x_1^3x_2   + x_1^3x_3   + x_1^3x_4  
+ x_2x_3^3 + x_2x_4^3
+ x_2^3x_4 + x_2^3x_3
+ x_3x_4^3 + x_3^3x_4  )\\
&+ 2( x_1x_2x_3^2  + x_1x_2^2x_4 + x_1x_3^2x_4 + x_1x_3x_4^2 
+ x_1^2x_2x_3 + x_1^2x_2x_4 + x_2^2x_3x_4  + x_2x_3x_4^2)\\
&- 3(x_1x_2^2x_3 + x_1x_2x_4^2 + x_1^2x_3x_4 + x_2x_3^2x_4).
\end{align*}

Indeed, for $d \in (\Z/5\Z)^\times$, we obtain an explicit bilinear map
$\varphi_d : K^4 \times K^4 \to K^4$
such that
\[
h({\bf x})h({\bf y}) = h(\varphi_d({\bf x},{\bf y}))
\]
for ${\bf x},{\bf y}\in V$.
It is given explicitly by $\varphi_d({\bf x},{\bf y})
= (z_{1,d}, z_{2,d}, z_{3,d}, z_{4,d})$
where
\[
z_{k,d}
= -\left(\sum_{i=1}^{4}x_{-d^{-1}i}y_{k-i} - \sum_{i=1}^{4}x_{-d^{-1}i}y_{-i}\right).
\]
For example, we have for $d=1 \in (\Z/5\Z)^\times$,
\begin{align*}
z_{1,1}
&=-(x_3y_4 + x_2y_3 + x_1y_2 )
+(x_4y_4 + x_3y_3 + x_2y_2 + x_1y_1 ),\\
z_{2,1}
&=-(x_{4}y_1 + x_2y_4 + x_1y_3 )
+(x_{4}y_4 + x_3y_3 + x_2y_2 + x_1y_1 ),\\
z_{3,1}
&=-(x_4y_2 + x_3y_1 + x_1y_4 )
+(x_4y_4 + x_3y_3 + x_2y_2 + x_1y_1 ),\\
z_{4,1}
&=-(x_{4}y_3 + x_3y_2 + x_2y_1)
+(x_{4}y_4 + x_3y_3 + x_2y_2 + x_1y_1 ).
\end{align*}
\end{example}

\begin{example}[$l=7, f=3, e=2, \gamma =3, \beta =2$]
\begin{align*}
h({\bf x})=h(x_1,\ldots,x_6)&=S_0 - S_1
\end{align*}
and
\begin{align*}
f_1({\bf x})=f_1(x_1,\ldots,x_6)=S_1-S_2
\end{align*}
where
\begin{align*}
S_0({\bf x})&= \sum_{s + 2t + 4u\equiv 0 \ (7)} x_{s}x_{t}x_{u}\\
&= x_1^3 + x_2^3 + x_4^3 + x_3^3 + x_5^3 + x_6^3\\
&+ 3(x_1x_2x_4) + 3(x_3x_5x_6) + 3(x_1x_2x_6 + x_2x_4x_5+x_1x_3x_4)  + 3(x_1x_5x_6 + x_2x_3x_5 +  x_3x_4x_6),\\
%
S_1({\bf x})&= \sum_{s + 2t + 4u\equiv 3 \ (7)} x_{s}x_{t}x_{u}\\
&=x_1x_2^2 + 
x_1x_3^2 + 
x_1x_5^2 +
x_2x_4^2 +
x_2x_3^2 +
x_2x_6^2 +
x_3x_4^2 + 
x_3x_5^2 +
x_4x_5^2 + 
x_4x_6^2 + 
x_5x_6^2 \\
%
&+x_1^2x_4 +
x_1^2x_6 + 
x_2^2x_5 +
x_3^2x_6 \\
%
&+x_1x_2x_3 + 
x_1x_2x_5 +
x_1x_3x_5 + 
x_1x_3x_6 +
x_1x_4x_6\\
&+ x_2x_3x_4 +
x_2x_3x_6 +
x_2x_4x_6 +
x_2x_5x_6 +
x_3x_4x_5 + 
x_4x_5x_6
\\
%
&+x_1x_2x_4 +
x_1x_4x_5 +
x_1x_5x_6
%
+x_2x_3x_5 +
x_3x_4x_6,\\
%
%
S_2({\bf x})&= \sum_{s + 2t + 4u\equiv 2 \ (7)} x_{s}x_{t}x_{u}\\
&= 
x_1x_4^2 +
x_1x_6^2 +
x_2 x_5^2+
x_3x_6^2 \\
&+ x_1^2x_2 +
x_1^2x_3 +
x_1^2x_5 +
%
x_2^2x_3 +
x_2^2x_4 +
x_2^2x_6 +
%
x_3^2x_4 +
x_3^2x_5 +
%
x_4^2x_6 +
x_4^2x_5 +
x_5^2x_6 +\\
%
&+ x_1x_2x_3 +
x_1x_2x_5 +
x_1x_3x_5 +
x_1x_3x_6 +
x_1x_4x_6\\
&+ x_2x_3x_4 +
x_2x_3x_6 +
x_2x_4x_6 +
x_2x_5x_6 +
x_3x_4x_5 +
x_4x_5x_6
\\
%
&+x_1x_2x_6 +
x_1x_3x_4 +
x_1x_4x_5
%
+ 
x_2x_4x_5 +
x_3x_5x_6.
%
\end{align*}
Thus
\begin{align*}
h({\bf x})&=x_1^3 + x_2^3 + x_3^3 + x_4^3 + x_5^3 + x_6^3\\
&- (x_1^2x_2 + x_1^2x_5 + x_1^2x_3 + x_1x_6^2+ x_1x_4^2   
%
+ x_2x_5^2 + x_2^2x_6  + x_2^2x_3 + x_2^2x_4 \\
%
&+ x_3x_6^2  + x_3^2x_4 + x_3^2x_5  + x_4^2x_6  + x_4^2x_5  + x_5^2x_6\\
%
&+ x_1x_2x_3 + x_1x_2x_5 + x_1x_3x_5  + x_1x_3x_6 + x_1x_4x_5  + x_1x_4x_6 \\
&+ x_2x_3x_4 + x_2x_3x_6 + x_2x_4x_6  + x_2x_5x_6 + x_3x_4x_5  + x_4x_5x_6)\\
&+ 2(x_1x_2x_4 + x_1x_5x_6+ x_2x_3x_5  + x_3x_4x_6)\\
&+ 3(x_1x_3x_4 + x_2x_4x_5 + x_1x_2x_6 + x_3x_5x_6),\\
\end{align*}
\begin{align*}
f_1({\bf x})&=x_1^2x_2 - x_1x_2^2 + x_1^2x_3 + x_2^2x_3 - x_1x_3^2 - x_2x_3^2 - x_1^2x_4 - x_1x_2x_4 + x_2^2x_4 + x_1x_3x_4 + x_3^2x_4 \\
&+ x_1x_4^2 - x_2x_4^2 - x_3x_4^2 + x_1^2x_5 - x_2^2x_5 - x_2x_3x_5 + x_3^2x_5 + x_2x_4x_5 + x_4^2x_5 - x_1x_5^2 + x_2x_5^2\\
&- x_3x_5^2 - x_4x_5^2 - x_1^2x_6 + x_1x_2x_6 + x_2^2x_6 - x_3^2x_6 - x_3x_4x_6 + x_4^2x_6 - x_1x_5x_6 + x_3x_5x_6 \\
&+ x_5^2x_6 + x_1x_6^2 - x_2x_6^2 + x_3x_6^2 - x_4x_6^2 - x_5x_6^2.
\end{align*}

Indeed, for $d \in (\Z/7\Z)^\times$, we obtain an explicit bilinear map $\varphi_d :K^6 \times K^6 \rightarrow K^6$ such that 
\begin{align*}
h({\bf x})h({\bf y})=h(\varphi_d ({\bf x}, {\bf y}))
\end{align*}
for ${\bf x},{\bf y} \in V$.
It is given explicitly by
$\varphi_d ({\bf x}, {\bf y}) = (z_{1,d}, z_{2,d}, z_{3,d}, z_{4,d},z_{5,d},z_{6,d})$
where 
\[
z_{k,d}=(-1)^{f-1}\left(\sum_{i=1}^{6}x_{-d^{-1}i}y_{k-i} - \sum_{i=1}^{6}x_{-d^{-1}i}y_{-i}\right).
\]
For example, we have for $d=1 \in (\Z/7\Z)^\times$,
\begin{align*}
z_{1,1}
&=( x_5y_6 + x_4y_5 + x_3y_4 + x_2y_3 + x_1y_2 )
-(x_6y_6 + x_5y_5 + x_4y_4 + x_3y_3 + x_2y_2 + x_1y_1 ),\\
z_{2,1}
&=(x_6y_1 + x_4y_6 + x_3y_5 + x_2y_4 + x_1y_3 )
-(x_6y_6 + x_5y_5 + x_4y_4 + x_3y_3 + x_2y_2 + x_1y_1 ),\\
z_{3,1}
&=(x_6y_2 + x_5y_1 + x_3y_6 + x_2y_5 + x_1y_4 )
-(x_6y_6 + x_5y_5 + x_4y_4 + x_3y_3 + x_2y_2 + x_1y_1 ),\\
z_{4,1}
&=(x_6y_3 + x_5y_2 + x_4y_1 + x_2y_6 + x_1y_5 )
-(x_6y_6 + x_5y_5 + x_4y_4 + x_3y_3 + x_2y_2 + x_1y_1 ),\\
z_{5,1}
&=( x_6y_4 + x_5y_3 + x_4y_2 + x_3y_1 + x_1y_6 )
-(x_6y_6 + x_5y_5 + x_4y_4 + x_3y_3 + x_2y_2 + x_1y_1 ),\\
z_{6,1}
&=(x_6y_5 + x_5y_4 + x_4y_3 + x_3y_2 + x_2y_1)
-(x_6y_6 + x_5y_5 + x_4y_4 + x_3y_3 + x_2y_2 + x_1y_1 ). 
\end{align*}

Moreover, for $d=1\in (\Z/7\Z)^\times$, 
we obtain the following relations 
which also guarantee that $h({\bf x})$ is multiplicative on $V$, i.e. 
$h({\bf x})h({\bf y})=h(\varphi_d({\bf x}, {\bf y}))$ for 
${\bf x},{\bf y} \in V=\{{\bf x}\in K^6\mid f_1({\bf x})=0\}$:
\begin{align*}
h(\varphi_d({\bf x},{\bf y}))&=
h({\bf x})h({\bf y})+f_1({\bf x})h({\bf y})+2f_1({\bf x})f_1({\bf y}),\\
f_1(\varphi_d({\bf x},{\bf y}))&=
h({\bf x})f_1({\bf y})-f_1({\bf x})h({\bf y}).
\end{align*}
\end{example}

\section{Proof of Theorem~\ref{thm:structureV}: The algebraic group $W$}\label{S6}

In this section we explain how the composition law constructed in Section~\ref{S5} gives rise, after restricting to a suitable dense open subset, to an algebraic $K$-group structure. 
Identifying $W$ with a norm-defined subgroup of the Weil restriction $R_{L/K}(\mathbb{G}_m)$ makes the operation $\overset{-1}{\ast}$ transparent and leads to the torus structure asserted in Theorem~\ref{thm:structureV}. 
Thus the torus structure records the group-theoretic content of the coefficient composition constructed in Section~\ref{S5}.
Once this description is in place, Corollary~\ref{cor:CI} follows from the resulting dimension argument.

Throughout this section, we retain the notation and assumptions of Theorem~\ref{mainth3}.
In particular, let $K$, $L=K(\zeta_l)$, $M$, and
\[
V=\{{\bf x}\in K^{l-1}\mid h_m({\bf x})=0 \ (m=1,\dots,e-1)\}
\]
where
\begin{align*}
h_m({\bf x})&=h_m(x_1,\ldots,x_{l-1})\\
&=\left(\sum_{i_0 + \beta i_1 + \cdots +\beta^{f-1}i_{f-1} \equiv \gamma^m \ (l)}x_{i_0}x_{i_1}\cdots x_{i_{f-1}}\right) - \left(\sum_{i_0 + \beta i_1 + \cdots +\beta^{f-1}i_{f-1} \equiv \gamma^{m+1} \ (l)}x_{i_0}x_{i_1}\cdots x_{i_{f-1}}\right)
\end{align*}
with subscripts taken modulo $l$ and with the convention $x_0=0$. 
For each $d\in(\Z/l\Z)^{\times}$, let
$
\varphi_d: V(K)\times V(K)\to V(K)
$
be the corresponding binary operation on $V(K)$, and define
\[
{\bf x}\overset d\ast {\bf y}:=\varphi_d({\bf x},{\bf y}).
\]
For ${\bf x}=(x_1,\dots,x_{l-1})\in V(K)$, we define
\[
\alpha_{\bf x}:=(-1)^{f-1}\sum_{i=1}^{l-1}x_i\zeta_l^{\,i}\in L,
\]
and
\[
W:=\{{\bf x}\in V\mid N_{L/K}(\alpha_{\bf x})\neq 0\}.
\]
The next proposition shows that $V(K)$ is a left unital groupoid with respect to $\overset{d}{\ast}$. 
In the special case $d=-1$, the operation becomes commutative and associative, which leads to the algebraic $K$-group structure on the open subset $W\subset V$ appearing in Theorem~\ref{thm:structureV} (cf. Hoshi and Kanai \cite[Section 3]{HK22}).

\begin{prop}\label{prop:CA}
For ${\bf x}, {\bf y} \in V(K)$ and $d,d_1,d_2\in (\Z/l\Z)^\times$,  the following hold: 
\begin{enumerate}
\item [$(1)$]  $({\bf 1}\overset{d}{\ast}{\bf x})_k=x_k$ and $({\bf x}\overset{d}{\ast}{\bf 1})_k=x_{-d^{-1}k}$ where ${\bf 1} := (-1)^{f}(1,1,\dots,1)$. 
The operation $\overset{d}{\ast}$ is left unital, i.e. ${\bf 1}\overset{d}{\ast}{\bf x}= {\bf x}$, and the operation $\overset{-1}{\ast}$ is unital, i.e. ${\bf 1}\overset{-1}{\ast}{\bf x}= {\bf x}\overset{-1}{\ast}{\bf 1}={\bf x}$.
\item [$(2)$] ${\bf x}\overset{d_1}{\ast}({\bf y}\overset{d_2}{\ast}{\bf z})=({\bf x}\overset{-d_2^{-1}d_1}{\ast}{\bf y})\overset{d_2}{\ast}{\bf z}$ where $d_1, d_2 \in (\Z/l\Z)^{\times}$. 
In particular, ${\bf x}\overset{d_1}{\ast}({\bf y}\overset{-1}{\ast}{\bf z})=({\bf x}\overset{d_1}{\ast}{\bf y})\overset{-1}{\ast}{\bf z}$.
\item [$(3)$] $({\bf y}\overset{d}{\ast}{\bf x})_k=({\bf x}\overset{d^{-1}}{\ast}{\bf y})_{-d^{-1}k}$. 
In particular, ${\bf y}\overset{d}{\ast}{\bf x}={\bf x}\overset{d}{\ast}{\bf y}$ if and only if $d=-1$.
\end{enumerate}
\end{prop}
\begin{proof}
(1) By the convention $y_0=0$, we have
\begin{align*}
({\bf x}\overset{d}{\ast}{\bf y})_k &= (-1)^{f-1}\left( \sum_{i=1}^{l-1}x_{-d^{-1}i}y_{k-i} - \sum_{i=1}^{l-1}x_{-d^{-1}i}y_{-i} \right)\\
&=(-1)^{f-1}\left( x_{-d^{-1}k}y_0 + \sum_{\substack{i=1\\i\neq k}}^{l-1}x_{-d^{-1}i}y_{k-i} - \sum_{i=1}^{l-1}x_{-d^{-1}i}y_{-i} \right)\\
&=(-1)^{f-1}\left( \sum_{\substack{i=1\\i\neq k}}^{l-1}x_{-d^{-1}i}y_{k-i} - \sum_{i=1}^{l-1}x_{-d^{-1}i}y_{-i} \right).
\end{align*}
Hence we have
\begin{align*}
({\bf x}\overset{d}{\ast}{\bf 1})_k &=(-1)^{f-1}\left( \sum_{\substack{i=1\\i\neq k}}^{l-1}x_{-d^{-1}i}\cdot (-1)^{f} - \sum_{i=1}^{l-1}x_{-d^{-1}i}\cdot (-1)^{f} \right)\\
&=(-1)^{(f-1)(f+1)}x_{-d^{-1}k}\\
&=x_{-d^{-1}k}
\end{align*}
and
\begin{align*}
({\bf 1}\overset{d}{\ast}{\bf x})_k &=(-1)^{f-1}\left( \sum_{\substack{i=1\\i\neq k}}^{l-1}(-1)^{f}\cdot x_{k-i} - \sum_{i=1}^{l-1}(-1)^{f}\cdot x_{-i} \right)\\
&=x_{k}.
\end{align*}

(2) We calculate the left-hand side and the right-hand side separately. 
Set ${\bf v}:={\bf y}\overset{d_2}{\ast}{\bf z}$ and ${\bf w}:={\bf x}\overset{-d_2^{-1}d_1}{\ast}{\bf y}$. 
\begin{align*}
(({\bf x}\overset{-d_2^{-1}d_1}{\ast}{\bf y})\overset{d_2}{\ast}{\bf z})_k&=({\bf w}\overset{d_2}{\ast}{\bf z})_k\\
&=(-1)^{f-1}\left( \sum_{j=1}^{l-1}w_{-d_2^{-1}j}z_{k-j} - \sum_{j=1}^{l-1}w_{-d_2^{-1}j}z_{-j} \right)\\
&=(-1)^{f-1}\left( \sum_{j=1}^{l-1}(-1)^{f-1}\left( \sum_{i=1}^{l-1}x_{d_2d_1^{-1}i}y_{(-d_2^{-1}j)-i} - \sum_{i=1}^{l-1}x_{d_2d_1^{-1}i}y_{-i} \right)z_{k-j}\right.\\
& \left. - \sum_{j=1}^{l-1}(-1)^{f-1}\left( \sum_{i=1}^{l-1}x_{d_2d_1^{-1}i}y_{(-d_2^{-1}j)-i} - \sum_{i=1}^{l-1}x_{d_2d_1^{-1}i}y_{-i} \right)z_{-j} \right)\\
&=\sum_{i=1}^{l-1}\sum_{j=1}^{l-1}x_{d_2d_1^{-1}i}y_{-d_2^{-1}j-i}z_{k-j} - \sum_{i=1}^{l-1}\sum_{j=1}^{l-1}x_{d_2d_1^{-1}i}y_{-i}z_{k-j} \\
&- \sum_{i=1}^{l-1}\sum_{j=1}^{l-1}x_{d_2d_1^{-1}i}y_{-d_2^{-1}j-i}z_{-j} + \sum_{i=1}^{l-1}\sum_{j=1}^{l-1}x_{d_2d_1^{-1}i}y_{-i}z_{-j}\\
&=\sum_{i=1}^{l-1}\sum_{j=1}^{l-1}x_{d_2d_1^{-1}i}y_{-d_2^{-1}j-i}z_{k-j} - \sum_{i=1}^{l-1}\sum_{j=1}^{l-1}x_{d_2d_1^{-1}i}y_{-d_2^{-1}j-i}z_{-j}.
\end{align*}
On the other hand, by setting $j^\prime :=i+j$ and $i^\prime:=-d_2^{-1}i$, we have
\begin{align*}
({\bf x}\overset{d_1}{\ast}({\bf y}\overset{d_2}{\ast}{\bf z}))_k&=({\bf x}\overset{d_1}{\ast}{\bf v})_k\\
&=(-1)^{f-1}\left( \sum_{i=1}^{l-1}x_{-d_1^{-1}i}v_{k-i} - \sum_{i=1}^{l-1}x_{-d_1^{-1}i}v_{-i} \right)\\
&=(-1)^{f-1} \left( \sum_{i=1}^{l-1}x_{-d_1^{-1}i}(-1)^{f-1}\left( \sum_{j=1}^{l-1}y_{-d_2^{-1}j}z_{(k-i)-j} - \sum_{j=1}^{l-1}y_{-d_2^{-1}j}z_{-j} \right) \right.\\
&- \left. \sum_{j=1}^{l-1}x_{-d^{-1}j}(-1)^{f-1} \left( \sum_{j=1}^{l-1}y_{-d_2^{-1}j}z_{(-i)-j} - \sum_{j=1}^{l-1}y_{-d_2^{-1}j}z_{-j} \right) \right)\\
&=\sum_{i=1}^{l-1}\sum_{j=1}^{l-1}x_{-d_1^{-1}i}y_{-d_2^{-1}j}z_{k-(i+j)} - \sum_{i=1}^{l-1}\sum_{j=1}^{l-1}x_{-d_1^{-1}i}y_{-d_2^{-1}j}z_{-j} \\
&- \sum_{i=1}^{l-1}\sum_{j=1}^{l-1}x_{d_1^{-1}i} y_{-d_2^{-1}j}z_{-(i+j)} + \sum_{i=1}^{l-1}\sum_{j=1}^{l-1}x_{-d_1^{-1}i}y_{-d_2^{-1}j}z_{-j}\\
&=\sum_{i=1}^{l-1}\sum_{j=1}^{l-1}x_{-d_1^{-1}i}y_{-d_2^{-1}j}z_{k-(i+j)} - \sum_{i=1}^{l-1}\sum_{j=1}^{l-1}x_{-d_1^{-1}i} y_{-d_2^{-1}j}z_{-(i+j)}\\
&=\sum_{i=1}^{l-1}\sum_{j^\prime=1}^{l-1}x_{-d_1i}y_{-d_2^{-1}(j^\prime-i)}z_{k-j^\prime} - \sum_{i=1}^{l-1}\sum_{j=1}^{l-1}x_{-d_1^{-1}i} y_{-d_2^{-1}(j^\prime-i)}z_{-j^\prime}\\
&=\sum_{i^\prime=1}^{l-1}\sum_{j^\prime=1}^{l-1}x_{d_2d_1^{-1}i^\prime}y_{-d_2^{-1}j^\prime-i^\prime}z_{k-j^\prime} - \sum_{i^\prime=1}^{l-1}\sum_{j=1}^{l-1}x_{d_2d_1^{-1}i^\prime} y_{-d_2^{-1}j^\prime + i^\prime}z_{-j^\prime}\\
&=(({\bf x}\overset{-d_2^{-1}d_1}{\ast}{\bf y})\overset{d_2}{\ast}{\bf z})_k.
\end{align*}

(3) By substituting $i:=-d^{-1}(k-i^\prime)$ and $j:=d^{-1}j^\prime$, we have
\begin{align*}
({\bf x}\overset{d^{-1}}{\ast}{\bf y})_{-d^{-1}k}&=
 (-1)^{f-1}\left( \sum_{i=1}^{l-1}x_{-di}y_{-d^{-1}k-i} - \sum_{j=1}^{l-1}x_{-dj}y_{-j} \right)\\
 &= (-1)^{f-1}\left( \sum_{i^\prime=1}^{l-1}x_{k-i^\prime}y_{-d^{-1}i^\prime} - \sum_{j^\prime=1}^{l-1}x_{-j^\prime}y_{-d^{-1}j^\prime} \right)\\
 &= (-1)^{f-1}\left( \sum_{i^\prime=1}^{l-1}y_{-d^{-1}i^\prime}x_{k-i^\prime} - \sum_{j^\prime=1}^{l-1}y_{-d^{-1}j^\prime}x_{-j^\prime} \right)\\
 &= ({\bf y}\overset{d}{\ast}{\bf x})_{k}.
\end{align*}
\end{proof}

\begin{remark}
We notice that the operation $\overset{d}{\ast}$ on $V(K)$ is parallel to the $d$-composition $\overset{d}{\ast}$ of matrices in \cite[Section 3]{HK22}. 
\end{remark}


\begin{example}
The multiplication $\overset{-1}{\ast}$ is given by 
\[
\overset{-1}{\ast} : V \times V \to V\ ; \ ( x_1 , \dots , x_{l-1}, y_1 , \dots , y_{l-1} ) \mapsto (z_1 , \dots , z_{l-1})
\]
where $z_{k} := z_{k,-1} = (-1)^{f-1}\left(\sum_{i=1}^{l-1}x_{i}y_{k-i} - \sum_{i=1}^{l-1}x_{i}y_{-i}\right)$.
When $l=3$, we have
\begin{align*}
z_{1}
=(-1)^{f-1} (x_2y_2 - (x_1y_2 + x_2y_1)),\\
z_{2}
=(-1)^{f-1} (x_1y_1 - (x_1y_2 + x_2y_1)).
\end{align*}
When $l=5$, we have
\begin{align*}
z_{1}
&=(-1)^{f-1}( x_2y_4 + x_3y_3 + x_4y_2
-(x_1y_4 + x_2y_3 + x_3y_2 + x_4y_1 )),\\
z_{2}
&=(-1)^{f-1}(x_1y_1 + x_3y_4 + x_4y_3
-(x_1y_4 + x_2y_3 + x_3y_2 + x_4y_1 )),\\
z_{3}
&=(-1)^{f-1}(x_1y_2 + x_2y_1 + x_4y_4
-(x_1y_4 + x_2y_3 + x_3y_2 + x_4y_1 )),\\
z_{4}
&=(-1)^{f-1}(x_1y_3 + x_2y_2 + x_3y_1
-(x_1y_4 + x_2y_3 + x_3y_2 + x_4y_1 )).
\end{align*}
When $l=7$, we have
\begin{align*}
z_{1}
&=(-1)^{f-1}( x_2y_6 + x_3y_5 + x_4y_4 + x_5y_3 + x_6y_2
-(x_1y_6 + x_2y_5 + x_3y_4 + x_4y_3 + x_5y_2 + x_6y_1 )),\\
z_{2}
&=(-1)^{f-1}( x_1y_1 +  x_3y_6 + x_4y_5 + x_5y_4 + x_6y_3
-(x_1y_6 + x_2y_5 + x_3y_4 + x_4y_3 + x_5y_2 + x_6y_1 )),\\
z_{3}
&=(-1)^{f-1}( x_1y_2 + x_2y_1 +  x_4y_6 + x_5y_5 + x_6y_4
-(x_1y_6 + x_2y_5 + x_3y_4 + x_4y_3 + x_5y_2 + x_6y_1 )),\\
z_{4}
&=(-1)^{f-1}( x_1y_3 + x_2y_2 + x_3y_1 + x_5y_6 + x_6y_5
-(x_1y_6 + x_2y_5 + x_3y_4 + x_4y_3 + x_5y_2 + x_6y_1 )),\\
z_{5}
&=(-1)^{f-1}( x_1y_4 + x_2y_3 + x_3y_2 + x_4y_1 + x_6y_6
-(x_1y_6 + x_2y_5 + x_3y_4 + x_4y_3 + x_5y_2 + x_6y_1 )),\\
z_{6}
&=(-1)^{f-1}( x_1y_5 + x_2y_4 + x_3y_3 + x_4y_2 + x_5y_1 
-(x_1y_6 + x_2y_5 + x_3y_4 + x_4y_3 + x_5y_2 + x_6y_1 )).\\ 
\end{align*}
\end{example}

Next, we relate the representation matrix $M_d$ to the field norm and derive an explicit formula for the inverse element.
In what follows, when matrix notation is used, we regard vectors in $K^{l-1}$ as row vectors.

\begin{lem}\label{lem:det-Md}
Let $l$ be an odd prime, let $d\in (\Z/l\Z)^\times$, and let ${\bf x}=(x_1,\dots,x_{l-1})\in K^{l-1}$.
Let $M_d$ be the $(l-1) \times (l-1)$ matrix with entries
\[
M_d[k,j]:=(-1)^{f-1}(x_{d^{-1}(k-j)}-x_{d^{-1}k}),\qquad 1\le k,j\le l-1
\]
where the indices are taken modulo $l$ with the convention $x_0=0$.
Then
\[
\det M_d = N_{L/K}(\alpha_{\bf x}).
\]
\end{lem}

\begin{proof}
We introduce an auxiliary $l \times l$ matrix $\\widehat{M}$ indexed by $0 \le k, j \le l-1$, defined by the same formula as $M_d$.
By definition (with indices taken modulo l and the convention $x_0 = 0$),
the $0$-th column of $\widehat{M}$ is identically zero: for every $k$ we have
\[
\widehat{M}[k,0]
= (-1)^{f-1}(x_{d^{-1}(k-0)}-x_{d^{-1}k})
= (-1)^{f-1}(x_{d^{-1}k}-x_{d^{-1}k})=0.
\]
Let $U=\{v\in \overline{K}^{l}\mid v_0=0\}$. Then $U$ is $\widehat M$-stable, and the induced endomorphism on the quotient $\overline{K}^{l}/U$ is zero. 
Moreover, the restriction $\widehat M|_U$ is represented by the $(l-1)\times(l-1)$ minor $M_d$. 
Hence
\[
\det(tI_{l}-\widehat M)=t\,\det(tI_{l-1}-M_d),
\]
so the eigenvalues of $\widehat M$ consist of $0$ together with the eigenvalues of $M_d$ (with multiplicity).
Let $\zeta_l$ be a primitive $l$-th root of unity. 
For each $a \in \{0, \dots, l-1\}$, consider the vector ${\bf v}_a = (1, \zeta_l^a, \zeta_l^{2a}, \dots, \zeta_l^{(l-1)a})$ .
A direct computation yields
\[
{\bf v}_a\widehat{M} = \lambda_a ({\bf v}_a - {\bf v}_0) \quad \text{where } \lambda_a = (-1)^{f-1} \sum_{i=1}^{l-1} x_i \zeta_l^{adi}.
\]
Using the basis $\{{\bf v}_0, {\bf v}_1, \dots, {\bf v}_{l-1}\}$, the matrix representation of $\widehat{M}$ is upper triangular, and its non-zero eigenvalues are precisely $\lambda_1, \dots, \lambda_{l-1}$.
Therefore,
\[
\det M_d = \prod_{a=1}^{l-1} \lambda_a = (-1)^{(f-1)(l-1)} \prod_{a=1}^{l-1} \left( \sum_{i=1}^{l-1} x_i \zeta_l^{adi} \right).
\]
Since $l$ is an odd prime, $l-1$ is even, so the sign factor $(-1)^{(f-1)(l-1)}$ reduces to $1$.
Furthermore, since $d$ is coprime to $l$, the map $a \mapsto ad \pmod l$ is a permutation of $(\Z/l\Z)^\times$. 
Thus, the product runs over all Galois conjugates of $\alpha_{\bf x}$, yielding the norm $N_{L/K}(\alpha_{\bf x})$.
\end{proof}

Using the determinant formula, we obtain the explicit form of the inverse element.

\begin{prop}\label{prop:inverse}
Fix $d\in (\Z/l\Z)^\times$. Let ${\bf x}\in V(K)$ and assume that $N_{L/K}(\alpha_{\bf x})\neq 0$.
Then there exists a unique element ${\bf x}^{-1}\in V(K)$ such that
\[
{\bf x}\overset{d}{\ast}{\bf x}^{-1}={\bf 1},\qquad {\bf 1}:=(-1)^f(1,\dots,1),
\]
and it is given by
\[
{\bf x}^{-1}={\bf 1}\,\frac{\widetilde{M_d}}{N_{L/K}(\alpha_{\bf x})},
\]
where $\widetilde{M_d}$ denotes the adjugate matrix of $M_d$.
\end{prop}

\begin{proof}
The map ${\bf y}\mapsto{\bf x}\overset{d}{\ast}{\bf y}$ is $K$-linear and, in the row-vector convention above, is represented by the matrix $M_d$; hence
${\bf x}\overset{d}{\ast}{\bf x}^{-1}={\bf 1}$ is equivalent to ${\bf x}^{-1}M_d ={\bf 1}$.
By Lemma~\ref{lem:det-Md}, $\det M_d=N_{L/K}(\alpha_{\bf x})\neq 0$, so $M_d$ is invertible and
\[
{\bf x}^{-1}={\bf 1}M_d^{-1}={\bf 1}\,\frac{\widetilde{M_d}}{\det M_d}
={\bf 1}\,\frac{\widetilde{M_d}}{N_{L/K}(\alpha_{\bf x})}.
\]

To verify ${\bf x}^{-1}\in V(K)$, use the identity
\[
\alpha_{{\bf x}\overset{d}{\ast}{\bf y}}=\sigma_{-d}(\alpha_{\bf x})\,\alpha_{\bf y}
\]
where $\alpha_{\bf x}:=(-1)^{f-1}\sum_{i=1}^{l-1}x_i\zeta_l^{\,i},\ 
\alpha_{\bf y}:=(-1)^{f-1}\sum_{i=1}^{l-1}y_i\zeta_l^{\,i}$.
This follows by a direct calculation from the definition of ${\bf x}\overset{d}{\ast}{\bf y}:=\varphi_{d}({\bf x},{\bf y})$.
By substituting ${\bf y}={\bf x}^{-1}$ and ${\bf x}\overset{d}{\ast}{\bf x}^{-1}={\bf 1}$, we get
$\alpha_{\bf 1}=\sigma_{-d}(\alpha_{\bf x})\alpha_{{\bf x}^{-1}}$.
Since $\alpha_{\bf 1}=1$, we have $\alpha_{{\bf x}^{-1}}=\sigma_{-d}(\alpha_{\bf x})^{-1}$.
Therefore,
\[
N_{L/M}(\alpha_{{\bf x}^{-1}})
= N_{L/M}\bigl(\sigma_{-d}(\alpha_{\bf x})\bigr)^{-1}
= \sigma_{-d}\bigl(N_{L/M}(\alpha_{\bf x})\bigr)^{-1}\in K^\times,
\]
because ${\bf x}\in V(K)$ implies $N_{L/M}(\alpha_{\bf x})\in K^\times$ and $\sigma_{-d}$ acts trivially on $K$.
Hence ${\bf x}^{-1}\in V(K)$.
\end{proof}

\begin{example}[Case $l=5$ and $f=4$]
Consider the case where $l=5$ and $f=4$. 
We have $L=K(\zeta_5)$ and $(-1)^{f-1}=-1$. 
Let $d=1$.
For ${\bf x} = (x_1, x_2, x_3, x_4)$, the matrix $M_1$ is given by
\[
M_1=
\begin{pmatrix}
x_1 & x_1-x_4 & x_1-x_3 & x_1-x_2\\
x_2-x_1 & x_2 & x_2-x_4 & x_2-x_3\\
x_3-x_2 & x_3-x_1 & x_3 & x_3-x_4\\
x_4-x_3 & x_4-x_2 & x_4-x_1 & x_4
\end{pmatrix}
\]
By Lemma~\ref{lem:det-Md}, the determinant is simply the norm:
\[
    \det M_1 = N_{L/K}(x_1\zeta_5 + x_2\zeta_5^2 + x_3\zeta_5^3 + x_4\zeta_5^4).
\]
The inverse vector ${\bf x}^{-1}$ is the solution to ${\bf x}^{-1} M_1  = (1,1,1,1)$.
\end{example}

Note that the multiplication map $({\bf x},{\bf y})\mapsto {\bf x}\overset{-1}{\ast}{\bf y}$ is given by polynomials in the coordinates, hence is a regular morphism on $V\times V$. 
The only possible denominators occur in the inverse map, via the factor $N_{L/K}(\alpha_{\bf x})^{-1}$ in Proposition~\ref{prop:inverse}. 
Accordingly, we remove the boundary $S:=\{x\in V \mid N_{L/K}(\alpha_x)=0\}$
and set $W:=V\setminus S$.
Then inversion is regular on $W$, and hence $(W,\overset{-1}{\ast})$ is an algebraic $K$-group.
That is, $W=V^\times=V\setminus S\quad {\text where}\quad S:=\{{\bf x}\in V \mid N_{L/K}(\alpha_{\bf x})=0\}$.
By Proposition~\ref{prop:CA}, $\overset{-1}{\ast}$ is commutative and associative.
Thus $W$ is an algebraic $K$-group with respect to $\overset{-1}{\ast}$.

We may therefore view $(W,\overset{-1}{\ast})$ as a commutative algebraic $K$-group. 
Let $T\subset R_{L/K}(\Gm)$ be the $K$-subgroup defined by the norm condition $N_{L/M}(\alpha)\in {\mathbb G}_{m,K}$. 
In the next proof we construct an explicit isomorphism of algebraic $K$-groups $W\simeq T$, and then verify that $T$ is an algebraic $K$-torus.

\begin{proof}[Proof of Theorem~\ref{thm:structureV}]
\quad\par
{\bf Step 1: $W\simeq T$ is an algebraic $K$-group.}\ Keep the notation of Theorem~\ref{thm:structureV}. 
Put $f:=[L:M]$; then $[L:K]=ef=l-1$.
Define a $K$-subgroup of $R_{L/K}(\mathbb{G}_m)$ by
\[
T:=N_{L/M}^{-1}(\mathbb{G}_{m,K})=\Bigl\{\alpha\in R_{L/K}(\mathbb{G}_m)\ \Bigm|\ N_{L/M}(\alpha)\in \mathbb{G}_{m,K}\Bigr\}\subset R_{L/K}(\mathbb{G}_m)
\]
where $\mathbb{G}_{m,K}$ is regarded as the diagonal subtorus of $R_{M/K}(\mathbb{G}_m)$ induced by $K \subset M$.

Let $\Psi:\A^{l-1}_K\to R_{L/K}(\A^1)$ be the $K$-linear isomorphism given on $K$-points by $\Psi({\bf x})=\alpha_{\bf x}$.
By Lemma~\ref{lem:BEW2}, the defining equations of $V$ are equivalent to the condition $N_{L/M}(\Psi({\bf x}))\in K$; in particular, for ${\bf x}\in W$ one has $N_{L/M}(\Psi({\bf x}))\in K^\times$.
Hence $\Psi$ restricts to a morphism $\Psi:W\to T$.
Moreover, by Theorem~\ref{thm:BF2} with $d=-1$, the operation $\overset{-1}{\ast}$ on $W$ is transported by $\Psi$ to multiplication in $R_{L/K}(\mathbb{G}_m)$, and the condition $N_{L/M}(\cdot)\in \mathbb{G}_{m,K}$ is preserved. 
Therefore $\Psi:W\to T$ is an isomorphism of algebraic $K$-groups.
\quad\par
{\bf Step 2: $T$ is an algebraic $K$-torus.}\ We have shown that $\Psi:W\to T$ is an isomorphism of algebraic $K$-groups over $K$.
To complete the proof, it remains to verify that $T$ is an algebraic $K$-torus; equivalently, after base change to an algebraic closure $\overline{K}$ of $K$, we construct an explicit isomorphism
\[
\Phi:(\mathbb{G}_{m,\overline{K}})^{\,1+e(f-1)} \xrightarrow{\ \sim\ } T_{\overline{K}}.
\]

For this, we use the standard description of Weil restriction after base change to $\overline{K}$.
We have
\[
L\otimes_K \overline{K}\simeq \prod_{\sigma\in\Hom_K(L, \overline{K})} \overline{K},\qquad
M\otimes_K \overline{K}\simeq \prod_{\rho\in\Hom_K(M, \overline{K})} \overline{K},
\]
and hence
\[
(R_{L/K}(\mathbb{G}_m))_{ \overline{K}}\simeq (\mathbb{G}_{m, \overline{K}})^{[L:K]}=(\mathbb{G}_{m, \overline{K}})^{ef},\qquad
(R_{M/K}(\mathbb{G}_m))_{ \overline{K}}\simeq (\mathbb{G}_{m, \overline{K}})^{e}.
\]
Choose an ordering of $\Hom_K(L, \overline{K})$ compatible with restriction to $M$, so that for each $\rho\in\Hom_K(M, \overline{K})$ the set of extensions $\sigma$ of $\rho$ has cardinality $f$.
Thus we may write coordinates on $(R_{L/K}(\mathbb{G}_m))_{ \overline{K}}\simeq (\mathbb{G}_{m, \overline{K}})^{ef}$ as $(x_{j,k})$ with $1\le j\le e$ and $1\le k\le f$.
Under these identifications, the base change of the norm map becomes
\[
(N_{L/M})_{ \overline{K}}:(\mathbb{G}_{m, \overline{K}})^{ef}\longrightarrow (\mathbb{G}_{m, \overline{K}})^e,\qquad
(x_{j,k})\longmapsto \Bigl(\prod_{k=1}^f x_{j,k}\Bigr)_{1\le j\le e},
\]
and the inclusion $\mathbb{G}_{m,K}\hookrightarrow R_{M/K}(\mathbb{G}_m)$ becomes the diagonal embedding
\[
\mathbb{G}_{m, \overline{K}}\hookrightarrow (\mathbb{G}_{m, \overline{K}})^e.
\]
Consequently,
\[
T_{ \overline{K}}
=\Bigl\{(x_{j,k})\in(\mathbb{G}_{m, \overline{K}})^{ef}\ \Bigm|\ \prod_{k=1}^f x_{1,k}=\cdots=\prod_{k=1}^f x_{e,k}\Bigr\}.
\]

Define a morphism $\Phi:(\mathbb{G}_{m, \overline{K}})^{\,1+e(f-1)}\to T_{ \overline{K}}$ by coordinates $(t,u_{j,k})$ and
\[
x_{j,k}=u_{j,k}\ (k=1,\dots,f-1),\qquad
x_{j,f}=\frac{t}{\prod_{k=1}^{f-1}u_{j,k}}.
\]
Then $\prod_{k=1}^f x_{j,k}=t$ for every $j$, and $\Phi$ takes values in $T_{ \overline{K}}$.
Conversely, given $(x_{j,k})\in T_{ \overline{K}}$, set $t=\prod_{k=1}^f x_{1,k}$ and $u_{j,k}=x_{j,k}$ for $k\le f-1$; the defining equalities imply that each $x_{j,f}$ is recovered by the same formula.
These constructions are regular and mutually inverse, hence $\Phi$ is an isomorphism. Therefore
\[
T_{ \overline{K}}\simeq (\mathbb{G}_{m, \overline{K}})^{\,1+e(f-1)}=(\mathbb{G}_{m, \overline{K}})^{\,l-e},
\]
and $T$ is an algebraic $K$-torus. Since $W\simeq T$, the theorem follows.
\end{proof}

\begin{proof}[Proof of Corollary~\ref{cor:CI}]
Set
\[
W:=\{{\bf x}\in V \mid N_{L/K}(\alpha_{\bf x})\neq 0\}\subset V.
\]
Consider the $K$-linear isomorphism $\Psi:\A_K^{l-1}\xrightarrow{\ \sim\ }R_{L/K}(\A^1)$, given on $K$-points by $\Psi({\bf x})=\alpha_{\bf x}$, and set
\[
\overline T:=\{\alpha\in R_{L/K}(\A^1)\mid N_{L/M}(\alpha)\in \A^1_K\}.
\]
Then we have $T=\overline T\cap R_{L/K}({\mathbb G}_m)$.
By Lemma~\ref{lem:BEW2} (equivalently, Step~1 of Theorem~\ref{thm:structureV}), the defining equations of $V$ are equivalent to $N_{L/M}(\Psi({\bf x}))\in K$, hence $\Psi(V)=\overline T$, and by definition $\Psi(W)=T$.
Since $R_{L/K}({\mathbb G}_m)\subset R_{L/K}(\A^1)$ is a nonempty open subscheme, it follows that $T$ is a nonempty open dense subscheme of $\overline T$; transporting via $\Psi$, we conclude that $W$ is a nonempty Zariski open dense subset of $V$.
Since $W$ is irreducible, $V$ is also irreducible.
Therefore $\dim V=\dim W$.

By Theorem~\ref{thm:structureV}, $W\simeq T$ is an algebraic $K$-torus of dimension $l-e$; hence $\dim V=l-e$.
Thus
\[
{\rm codim}_{\A_K^{l-1}}(V)=(l-1)-(l-e)=e-1.
\]
Since $V\subset \A_K^{l-1}$ is defined by exactly $e-1$ equations $h_m({\bf x})=0$ $(m=1,\dots,e-1)$,
each of which is an $f$-ic form, it follows that $V$ is a complete intersection in $\mathbb{A}^{l-1}_K$, cut out by $e-1$ $f$-ics over $K$.
\end{proof}

\begin{proof}[Proof of Corollary~\ref{cor:lifting}]
Let $x\in W_{{\bf p}}(K)$ and $y\in W_{{\bf q}}(K)$. By Theorem~\ref{mainth3} we have $h(x\overset{-1}{\ast}y)=h(x)h(y)={\bf pq}$ , hence $x\overset{-1}{\ast}y\in W_{{\bf pq}}(K)$.
\end{proof}

\begin{remark}\label{rem:fibers}
Since $h$ is multiplicative with respect to the group law $\overset{-1}{\ast}$
on $W(K)$, the map $h:W(K)\to K^\times$ is a group homomorphism, and
$W_{\bf p}(K)=h^{-1}({\bf p})$ for ${\bf p} \in K^{\times}$.
Under the isomorphism of Theorem~\ref{thm:structureV}, the fiber $W_1$ is the
norm one torus
\[
W_1=
\ker\bigl(N_{L/M}:R_{L/K}(\mathbb{G}_m)\longrightarrow R_{M/K}(\mathbb{G}_m)\bigr)=R_{M/K}(R_{L/M}^{(1)}(\mathbb{G}_m)).
\]
In particular,
\[
W_1(K)=\{\alpha\in L^\times\mid N_{L/M}(\alpha)=1\}.
\]

The group $W_1(K)$ acts on $W_{\bf p}(K)\times W_{\bf q}(K)$ by
\[
u\cdot(x,y):=
\bigl(x\overset{-1}{\ast}u,\,
y\overset{-1}{\ast}u^{-1}\bigr),
\]
where inverses are taken with respect to the group law $\overset{-1}{\ast}$ on
$W(K)$. 
This action preserves the fibers of the map in
Corollary~\ref{cor:lifting}, and each nonempty fiber is a principal homogeneous
space under $W_1(K)$.

Indeed, if $(x,y)$ and $(x',y')$ have the same image in $W_{{\bf p}{\bf q}}(K)$, then the
unique element carrying $(x,y)$ to $(x',y')$ is
\[
u=x^{-1}\overset{-1}{\ast}x'\in W_1(K).
\]
Moreover, if $y_0\in W_{\bf q}(K)$ is fixed, then translation by $y_0$ gives a
bijection of $K$-rational fibers
\[
W_{\bf p}(K)\xrightarrow{\ \sim\ }W_{{\bf p}{\bf q}}(K),\qquad
x\longmapsto x\overset{-1}{\ast}y_0,
\]
with inverse
\[
z\longmapsto z\overset{-1}{\ast}y_0^{-1}.
\]
Thus, whenever $W_{\bf q}(K)$ is nonempty, the map in Corollary~\ref{cor:lifting} is
surjective, and its fibers are precisely the orbits of the $W_1(K)$-action
defined above.
\end{remark}

\begin{remark}[The boundary $S$ and torsion points]\label{rem:boundary-torsion}
\quad \\
\begin{enumerate}
\item
Under the standing assumption $[K(\zeta_l):K]=l-1$, the map
$\iota : K^{l-1}\to L\ ;\ {\bf x}\mapsto \alpha_{\bf x}$ is a $K$-linear isomorphism.
Hence
\[
S(K)=\{x\in V\mid x\in K^{\,l-1},\ N_{L/K}(\alpha_x)=0\}=\{0\}.
\]

\item
Fix an algebraic closure $ \overline{K}$ of $K$. Since
\[
N_{L/K}(\alpha)=\prod_{\sigma\in\Hom_K(L, \overline{K})}\sigma(\alpha)
\qquad(\alpha\in L),
\]
we have
\[
\quad \quad \quad S( \overline{K})=\{x\in V( \overline{K})\mid N_{L/K}(\alpha_x)=0\}
\subset
\bigcup_{\sigma\in\Hom_K(L, \overline{K})}\{x\in K^{\,l-1}\otimes_K \overline{K}\mid \sigma(\alpha_x)=0\}.
\]
Thus, viewed inside $K^{\,l-1}\otimes_K \overline{K}\simeq  \overline{K}^{\,l-1}$, the set of
solutions of $N_{L/K}(\alpha_x)=0$ is a union of $ \overline{K}$-hyperplanes given by the
linear equations $\sigma(\alpha_x)=0$.

On the other hand, using the coordinates $(x_{j,k})$ ($1\le j\le e,\ 1\le k\le f$)
from Step~2 in the proof of Theorem~\ref{thm:structureV}, one has
\[
V( \overline{K})=
\Bigl\{(x_{j,k})\in  \overline{K}^{\,ef}\ \Big|\ \prod_{k=1}^f x_{1,k}
=\prod_{k=1}^f x_{2,k}
=\cdots
=\prod_{k=1}^f x_{e,k}\Bigr\}.
\]
If we define $t:=\prod_{k=1}^f x_{1,k}$, then on $V( \overline{K})$ we obtain
\[
N_{L/K}(\alpha_x)=\prod_{j=1}^e\prod_{k=1}^f x_{j,k}
=\prod_{j=1}^e\Bigl(\prod_{k=1}^f x_{j,k}\Bigr)=t^e.
\]
In particular,
\[
S( \overline{K})=\{(x_{j,k})\in V( \overline{K})\mid t=0\}.
\]
Equivalently, $S( \overline{K})$ is the union of the subsets
\[
B_{\mathbf k}( \overline{K}):=
\Bigl\{(x_{j,k})\in V( \overline{K})\ \Big|\ x_{1,k_1}=x_{2,k_2}=\cdots=x_{e,k_e}=0\Bigr\}
\]
where $\mathbf k=(k_1,\dots,k_e)\in\{1,\dots,f\}^e$,
so there are $f^e$ maximal boundary pieces inside $V( \overline{K})$ (each obtained by
forcing one coordinate to be $0$ in each block).

\item
Under the identification $W_{ \overline{K}}\simeq T_{ \overline{K}}\simeq
\mathbb{G}_{m, \overline{K}}\times(\mathbb{G}_{m, \overline{K}})^{e(f-1)}$ from the proof of Theorem~\ref{thm:structureV},
an element of order $n$ in $W( \overline{K})$ has all coordinates in the group
$\mu_n( \overline{K})=\{\xi\in \overline{K}^\times\mid \xi^n=1\}$.
In particular, if $x\in W( \overline{K})$ has order $n$, then $N_{L/K}(\alpha_x)\in \mu_n( \overline{K})$.
Hence a point $x\in W( \overline{K})$ with $N_{L/K}(\alpha_x)$ not a root of unity
cannot be a torsion point of $W( \overline{K})$.

\item
In the examples of Section~~\ref{S4-3}, the preceding criterion shows
that the classical trivial solutions need not be torsion points.  
For example, in Dickson's system for $l=5$, written in the normalization
$16P=x_1^2+125x_2^2+50x_3^2+50x_4^2$, the trivial solution
$(x_1,x_2,x_3,x_4)=(4p^r,0,0,0)$ corresponds to the parameter $P=p^{2r}$.
Its norm $N_{L/K}(\alpha_x)$ is a nonzero $p$-power up to normalization, so
$x\in W$ and $x$ is not torsion in $W$ by (3).
\end{enumerate}

\end{remark}

Corollary~\ref{cor:lifting} and Remark~\ref{rem:fibers} describe the compatibility between the group law on $W$ and 
the fibers $W_{\bf p}(K)$ of the multiplicative form $h$: 
the product map sends $W_{\bf p}(K)\times W_{\bf q}(K)$ to $W_{{\bf p}{\bf q}}(K)$, 
and each nonempty fiber of this map is a principal homogeneous space under $W_1(K)$, 
where $W_1$ is the norm one torus.

For every $f\ge 2$, the dense open torus $W\subset V$ provides a group-theoretic way to produce new elements of the fibers $W_{\bf p}(K)$: they are transported by the power maps $[n]:W\to W$, and more generally by the group law on $W$. 
Thus the explicit multiplicative $f$-ic forms constructed here give rise to higher-degree Diophantine systems whose solutions over $K$, on the dense open part $W$, are organized by this torus structure.

In the quadratic case $f=2$, this recovers the solution-lifting formulas of
Hoshi--Kanai~\cite{HK22}. 
More precisely, in the prime-power setting considered there, for instance when ${\bf p}={\bf q}=p^r$, 
iteration of the power maps gives the geometric counterpart of the lifting procedure obtained
from Davenport--Hasse's lifting theorem for Jacobi sums and from the corresponding $n$-fold products of multiplication matrices.

\begin{acknowledgment}
Before getting a Ph.D., the first-named author stayed at the
University of Regensburg for one year: 2004--2005. 
He thanks his teacher Manfred Knebusch 
who led him to the extensive world of quadratic forms. 
\end{acknowledgment}

%%%%%%%%%%%%%%%%%%%%%%%%%%%%%%%%%%%%
\begin{thebibliography}{Gau1801}

\bibitem[BE81]{BE81}
B. C. Berndt, R. J. Evans, 
{\it The determination of Gauss sums}, 
Bull. Amer. Math. Soc. (N.S.) \textbf{5} (1981), 107--129.

\bibitem[BEW98]{BEW98}
B. C. Berndt, R. J. Evans, K. S. Williams,
\textit{Gauss and Jacobi sums},
Canadian Mathematical Society Series of Monographs and Advanced Texts, 
Wiley, New York, 1998.

\bibitem[DH35]{DH35}
H. Davenport, H. Hasse,
{\it Die Nullstellen der Kongruenzzetafunktionen in gewissen 
zyklischen F\"{a}llen}, (German) 
J. Reine Angew. Math. \textbf{172} (1935), 151--182. 

\bibitem[DV09]{DV09}
A. Dr\'apal, P. Vojt\v{e}chovsk\'y,
{\it Symmetric multilinear forms and polarization of polynomials}, (German) 
Linear Algebra Appl. \textbf{431} (2009), 998--1012. 

\bibitem[Gau1801]{Gau1801}
C. F. Gauss, 
\textit{Disquisitiones Arithmeticae}, 1801.

\bibitem[HH05a]{HH05a}
K. Hashimoto, A. Hoshi, 
{\it Families of cyclic polynomials obtained from geometric generalization 
of Gaussian period relations}, 
Math. Comp. \textbf{74} (2005), 1519--1530.

\bibitem[HH05b]{HH05b}
K. Hashimoto, A. Hoshi, 
{\it Geometric generalization of Gaussian period relations with application 
to Noether's problem for meta-cyclic groups}, 
Tokyo J. Math. \textbf{28} (2005), 13--32.

\bibitem[Hos03]{Hos03}
A. Hoshi, 
{\it Multiplicative quadratic forms on algebraic varieties}, 
Proc. Japan Acad. Ser. A Math. Sci. \textbf{79} (2003), 71--75.

\bibitem[Hos06]{Hos06}
A. Hoshi, 
{\it Explicit lifts of quintic Jacobi sums and period polynomials 
for $\F_q$}, 
Proc. Japan Acad. Ser. A Math. Sci. \textbf{82} (2006), 87--92.

\bibitem[HK22]{HK22}
A. Hoshi, K. Kanai, 
{\it Davenport and Hasse's theorems and lifts of multiplication 
matrices of Gaussian periods}, 
Finite Fields Appl. \textbf{84} (2022) Paper No. 102101, 26 pp.

\bibitem[Hur1898]{Hur1898} 
A. Hurwitz, 
{\it \"Uber die Composition der quadratischen Formen von beliebig vielen 
Variablen}, 
Nachrichten von der Gesellschaft der Wissenschaften zu G\"ottingen, 
Mathematisch--Physikalische Klasse (1898), 309--316.

\bibitem[KR85a]{KR85a}
S. A. Katre, A. R. Rajwade,
\textit{Complete solution of the cyclotomic problem in $\F_q^{\ast}$ 
for any prime modulus $l$, $q = p^{\alpha}$, 
$p \equiv 1\ ({\rm mod}\ l)$},
Acta Arith. \textbf{45} (1985), 183--199.

\bibitem[KR85b]{KR85b}
S. A. Katre, A. R. Rajwade,
\textit{Unique determination of cyclotomic numbers of order five},
Manuscripta Math. \textbf{53} (1985), 65--75.

\bibitem[KS80]{KS80}
M. Knebusch, W. Scharlau, 
{\it Algebraic theory of quadratic forms. 
Generic methods and Pfister forms}, 
Notes taken by Heisook Lee. DMV Seminar, 1, 
Birkh\"auser, Boston, Mass., 1980. i+44 pp.

\bibitem[Lam05]{Lam05}
T. Y. Lam, 
{\it Introduction to quadratic forms over fields}, 
Graduate Studies in Mathematics, 67, American Mathematical Society, 
Providence, RI, 2005, xxii+550 pp.

\bibitem[Leh51]{Leh51} 
E. Lehmer, 
{\it The quintic character of $2$ and $3$}, 
Duke Math. J. \textbf{18} (1951), 11--18. 

\bibitem[Leh88]{Leh88} 
E. Lehmer, {\it Connection between Gaussian periods and cyclic units}, 
Math. Comp. \textbf{50} (1988), 535--541.

\bibitem[LW75a]{LW75a}
P. A. Leonard, K. S. Williams,
\textit{The cyclotomic numbers of order seven},
Proc. Amer. Math. Soc. \textbf{51} (1975), 295--300.

\bibitem[LW75b]{LW75b}
P. A. Leonard, K. S. Williams,
\textit{The cyclotomic numbers of order eleven},
Acta Arith. \textbf{26} (1975), 365--383.

\bibitem[Man74]{Man74}
Yu. I. Manin, 
\textit{Cubic forms: algebra, geometry, arithmetic}, 
%Translated from the Russian by M. Hazewinkel. 
%North-Holland Mathematical Library, Vol. 4. 
%North-Holland Publishing Co., Amsterdam-London; 
%American Elsevier Publishing Co., New York, 1974. vii+292 pp.
North-Holland Mathematical Library 4, North-Holland, Amsterdam, 1974.

\bibitem[Pfi65a]{Pfi65a}
A. Pfister, 
{\it Zur Darstellung von $-1$ als Summe von Quadraten in einem K\"orper}, 
(German) J. London Math. Soc. \textbf{40} (1965), 159--165. 

\bibitem[Pfi65b]{Pfi65b}
A. Pfister, 
{\it Multiplikative quadratische Formen}, (German) 
Arch. Math. (Basel) \textbf{16} (1965), 363--370.

\bibitem[Pfi66]{Pfi66}
A. Pfister, 
{\it Quadratische Formen in beliebigen K\"orpern}, (German) 
Invent. Math. \textbf{1} (1966), 116--132.

\bibitem[Pfi67]{Pfi67}
A. Pfister, 
{\it Zur Darstellung definiter Funktionen als Summe von Quadraten}, (German) 
Invent. Math. \textbf{4} (1967), 229--237.

\bibitem[Pfi95]{Pfi95}
A. Pfister, 
{\it Quadratic forms with applications to algebraic geometry and topology}, 
London Mathematical Society Lecture Note Series, 217, 
Cambridge University Press, Cambridge, 1995, viii+179 pp.

\bibitem[Raj93]{Raj93}
A. R. Rajwade, 
{\it Squares}, 
London Mathematical Society Lecture Note Series, 171, 
Cambridge University Press, Cambridge, 1993, xii+286 pp. 

\bibitem[Sch85]{Sch85}
W. Scharlau, 
{\it Quadratic and Hermitian forms}, 
Grundlehren der mathematischen Wissenschaften, 270, 
Springer-Verlag, Berlin, 1985. 

\bibitem[Sch70]{Sch70}
R.D. Schafer, 
{\it Forms permitting composition},
Adv. Math. \textbf{4} (1970), 127--148.

\bibitem[Sha00]{Sha00}
D. B. Shapiro, 
{\it Compositions of quadratic forms}, 
De Gruyter Expositions in Mathematics, 33, 
Walter de Gruyter \& Co., Berlin, 2000, xiv+417 pp.

\bibitem[Sto67]{Sto67}
T. Storer, 
{\it Cyclotomy and difference sets}, 
Lectures in Advanced Mathematics, 
No. 2 Markham Publishing Co., Chicago, Ill. 1967, vii+134 pp.

\bibitem[Sto69]{Sto69}
T. Storer, 
{\it On the arithmetic structure of Galois domains}, 
Acta Arith. \textbf{15} (1969), 139--159.

\bibitem[SW88]{SW88}
R. Schoof, L. C. Washington, 
{\it Quintic polynomials and real cyclotomic fields with large class numbers}, 
Math. Comp. \textbf{50} (1988), 543--556. 

\bibitem[Ter99]{Ter99}
A. Terras, 
{\it Fourier analysis on finite groups and applications}, 
London Mathematical Society Student Texts, 43, 
Cambridge University Press, Cambridge, 1999, x+442 pp. 

\bibitem[Tha96]{Tha96}
F. Thaine, 
\textit{Properties that characterize Gaussian periods and cyclotomic numbers}, 
Proc. Amer. Math. Soc. \textbf{124} (1996), 35--45.  

\bibitem[Tha99]{Tha99}
F. Thaine, 
\textit{On the coefficients of Jacobi sums in prime cyclotomic fields},
Trans. Amer. Math. Soc. \textbf{351} (1999), 4769--4790.

\bibitem[Tha00]{Tha00}
F. Thaine, 
\textit{Families of irreducible polynomials of Gaussian periods and matrices of cyclotomic numbers},
Math. Comp. \textbf{69} (2000), 1653--1666.

\bibitem[Tha01]{Tha01}
F. Thaine, 
\textit{Jacobi sums and new families of irreducible polynomials of Gaussian periods},
Math. Comp. \textbf{70} (2001), 1617--1640.

\bibitem[Tha04]{Tha04}
F. Thaine, 
\textit{Cyclic polynomials and the multiplication matrices of their roots},
J. Pure Appl. Algebra \textbf{188} (2004), 247--286.

\bibitem[Tha08]{Tha08} 
F. Thaine, 
{\it On the construction of families of cyclic polynomials whose roots 
are units},
Experiment. Math. \textbf{17} (2008), 315--331.

\bibitem[Vos98]{Vos98} 
V. E. Voskresenskii, 
{\it Algebraic groups and their birational invariants}, 
Translated from the Russian manuscript by Boris Kunyavskii, 
Translations of Mathematical Monographs, 179. 
American Mathematical Society, Providence, RI, 1998.

\bibitem[Wam02]{Wam02}
P. van Wamelen, 
\textit{Jacobi sums over finite fields},
Acta Arith. \textbf{102} (2002), 1--20.

\bibitem[Wei49]{Wei49}
A. Weil, 
{\it Numbers of solutions of equations in finite fields}, 
Bull. Amer. Math. Soc. \textbf{55} (1949), 497--508. 
\end{thebibliography}
\end{document}
